\chapter{\texorpdfstring{Geometry}{7. Geometry}}

\section{\texorpdfstring{Geometry Primitives}{7.1 Geometry Primitives}}
\setcounter{section}{1}
\setcounter{subsection}{0}
\subsection{Point}
\label{sec:7.1.1}
A 2D point class templated on the coordinate type \inlinecode{T}. Algebraic operations preserve \inlinecode{T}, while metric operations use \inlinecode{fp\_t}, which is \inlinecode{T} for floating-point coordinates and \inlinecode{double} otherwise.

Non-floating-point coordinate types such as \inlinecode{Modular} or \inlinecode{Rational} compose for exact operations (arithmetic, \inlinecode{dot}, \inlinecode{cross}, \inlinecode{sqnorm}, comparisons, and \inlinecode{EQ}) using the coordinate type's own operators. Their metric operations convert coordinates explicitly to \inlinecode{double}; for \inlinecode{Modular}, this uses the stored representative and is not finite-field geometry.

\begin{compactitem}
  \item \inlinecode{TPoint\textless{}T\textgreater{}()} constructs the origin, while \inlinecode{TPoint\textless{}T\textgreater{}(x, y)} constructs point $(x, y)$. A pair of values can also be converted explicitly to a point.
  \item Operators \inlinecode{+}, \inlinecode{-}, \inlinecode{*}, \inlinecode{/}, and their compound forms act element-wise on points or apply a scalar. Comparisons are exact and lexicographic, as required by standard algorithms and containers. \inlinecode{EQ(p, q)} instead uses \inlinecode{EPS} for floating-point coordinates and remains exact for other coordinate types.
  \item \inlinecode{sqnorm()}, \inlinecode{dot(p)}, and \inlinecode{cross(p)} return the exact algebraic results in coordinate type \inlinecode{T}.
  \item \inlinecode{norm()}, \inlinecode{arg()}, \inlinecode{proj(p)}, and \inlinecode{normalize()} return metric results in \inlinecode{fp\_t}; division also promotes to \inlinecode{fp\_t} to avoid integer truncation.
  \item \inlinecode{rotate90()}, \inlinecode{rotate180()}, and \inlinecode{rotate270()} rotate by the corresponding counter-clockwise cardinal angle; overloads taking \inlinecode{p} rotate about point \inlinecode{p}.
  \item \inlinecode{rotate\_cw(t)} and \inlinecode{rotate\_ccw(t)} rotate by an arbitrary angle \inlinecode{t} in radians; overloads taking (\inlinecode{p}, \inlinecode{t}) rotate about point \inlinecode{p}.
  \item \inlinecode{reflect(p)} reflects across point \inlinecode{p}, while \inlinecode{reflect(p, q)} reflects across the line through points \inlinecode{p} and \inlinecode{q}.
  \item \inlinecode{to\_double()} and \inlinecode{to\_ldouble()} explicitly convert the coordinate type. Integral points also convert implicitly to floating-point point types, so \inlinecode{PointI} can be passed where \inlinecode{PointD} is expected.
  \item \inlinecode{PointI}, \inlinecode{PointL}, \inlinecode{PointD}, and \inlinecode{PointLD} use \inlinecode{int}, \inlinecode{long long}, \inlinecode{double}, and \inlinecode{long double} coordinates, respectively. \inlinecode{Point} aliases \inlinecode{PointD}.
\end{compactitem}

Overflow warning: the exact products \inlinecode{dot()}, \inlinecode{cross()}, and \inlinecode{sqnorm()} grow like the squared coordinate magnitude. With \inlinecode{PointI} these overflow a 32-bit \inlinecode{int} once coordinates exceed roughly a few tens of thousands, so use \inlinecode{PointL} for larger integer coordinates.

\textbf{Time Complexity}: $O(1)$ per call to the constructor and all other operations.

\Needspace*{3\baselineskip}
\textbf{Space Complexity}:
\begin{compactitem}
  \item $O(1)$ for storage of the point.
  \item $O(1)$ auxiliary for all operations.
\end{compactitem}

\begin{lstlisting}[style=cpp]
#include <cmath>
#include <ostream>
#include <tuple>
#include <type_traits>
#include <utility>

const double EPS = 1e-9;

// Epsilon-aware for floating-point coordinates; exact for int and for coordinate types like Modular
// or Rational, which therefore compose for all of the predicates below.
template<typename T, typename U, typename C = std::common_type_t<T, U>>
bool EQ(T a, U b) {
  if constexpr (std::is_floating_point_v<C>) return C(a) == C(b) || std::fabs(C(a) - C(b)) <= EPS;
  return C(a) == C(b);
}

template<typename T, typename U, typename C = std::common_type_t<T, U>>
bool LT(T a, U b) {
  if constexpr (std::is_floating_point_v<C>) return C(a) < C(b) - EPS;
  return C(a) < C(b);
}

template<typename T>
struct TPoint {
  // Metric operations preserve floating-point coordinates and otherwise promote to double.
  using fp_t = std::conditional_t<std::is_floating_point<T>::value, T, double>;

  T x, y;

  TPoint() : x(0), y(0) {}
  TPoint(T x, T y) : x(x), y(y) {}
  explicit TPoint(const std::pair<T, T> &p) : x(p.first), y(p.second) {}

  // Implicit conversion from integral to floating-point point (optional, can just use to_double()).
  template<
      typename U,
      typename = std::enable_if_t<std::is_integral<U>::value && std::is_floating_point<T>::value>>
  TPoint(const TPoint<U> &p) : x(static_cast<T>(p.x)), y(static_cast<T>(p.y)) {}

  friend bool EQ(const TPoint &a, const TPoint &b) { return EQ(a.x, b.x) && EQ(a.y, b.y); }

  bool operator==(const TPoint &p) const { return x == p.x && y == p.y; }
  bool operator!=(const TPoint &p) const { return !(*this == p); }
  bool operator<(const TPoint &p) const { return std::tie(x, y) < std::tie(p.x, p.y); }
  bool operator>(const TPoint &p) const { return p < *this; }
  bool operator<=(const TPoint &p) const { return !(*this > p); }
  bool operator>=(const TPoint &p) const { return !(*this < p); }
  TPoint operator+(const TPoint &p) const { return {x + p.x, y + p.y}; }
  TPoint operator-(const TPoint &p) const { return {x - p.x, y - p.y}; }
  TPoint operator+(T v) const { return {x + v, y + v}; }
  TPoint operator-(T v) const { return {x - v, y - v}; }
  TPoint operator*(T v) const { return {x * v, y * v}; }

  // Division always promotes to fp_t to avoid integer truncation.
  TPoint<fp_t> operator/(fp_t v) const { return {fp_t(x) / v, fp_t(y) / v}; }

  TPoint &operator+=(const TPoint &p) { x += p.x; y += p.y; return *this; }
  TPoint &operator-=(const TPoint &p) { x -= p.x; y -= p.y; return *this; }
  TPoint &operator+=(T v) { x += v; y += v; return *this; }
  TPoint &operator-=(T v) { x -= v; y -= v; return *this; }
  TPoint &operator*=(T v) { x *= v; y *= v; return *this; }
  friend TPoint operator+(T v, const TPoint &p) { return p + v; }
  friend TPoint operator*(T v, const TPoint &p) { return p * v; }

  // --- Exact operations: return T or TPoint<T>, work for any coordinate type ---

  T sqnorm() const { return x * x + y * y; }                    // Overflow warning.
  T dot(const TPoint &p) const { return x * p.x + y * p.y; }    // Overflow warning.
  T cross(const TPoint &p) const { return x * p.y - y * p.x; }  // Overflow warning.

  // --- Floating-point operations: return fp_t or TPoint<fp_t> ---

  fp_t norm() const { return std::hypot(fp_t(x), fp_t(y)); }
  fp_t arg() const { return std::atan2(fp_t(y), fp_t(x)); }
  fp_t proj(const TPoint &p) const { return (fp_t(x) * p.x + fp_t(y) * p.y) / p.norm(); }

  TPoint<fp_t> normalize() const {
    fp_t n = norm();
    if (n < 1e-30) {              // guard against dividing by a near-zero norm, not a
      return {fp_t{0}, fp_t{0}};  // geometric tolerance (EPS is too coarse here)
    }
    return {fp_t(x) / n, fp_t(y) / n};
  }

  // --- Cardinal rotations: exact for all coordinate types including int ---

  // Returns (x, y) rotated 90/180/270 degrees counter-clockwise about the origin.
  TPoint rotate90() const { return {-y, x}; }
  TPoint rotate180() const { return {-x, -y}; }
  TPoint rotate270() const { return {y, -x}; }

  // Returns (x, y) rotated 90/180/270 degrees counter-clockwise about point p.
  TPoint rotate90(const TPoint &p) const { return (*this - p).rotate90() + p; }
  TPoint rotate180(const TPoint &p) const { return (*this - p).rotate180() + p; }
  TPoint rotate270(const TPoint &p) const { return (*this - p).rotate270() + p; }

  // --- Arbitrary-angle rotations: always return floating-point ---

  // Returns (x, y) rotated t radians clockwise about the origin.
  TPoint<fp_t> rotate_cw(fp_t t) const {
    fp_t fx(x), fy(y);
    return {fx * std::cos(t) + fy * std::sin(t), fy * std::cos(t) - fx * std::sin(t)};
  }

  // Returns (x, y) rotated t radians counter-clockwise about the origin.
  TPoint<fp_t> rotate_ccw(fp_t t) const {
    fp_t fx(x), fy(y);
    return {fx * std::cos(t) - fy * std::sin(t), fx * std::sin(t) + fy * std::cos(t)};
  }

  // Returns (x, y) rotated t radians clockwise about point p.
  TPoint<fp_t> rotate_cw(const TPoint &p, fp_t t) const {
    return TPoint<fp_t>{fp_t(x) - p.x, fp_t(y) - p.y}.rotate_cw(t) +
           TPoint<fp_t>{fp_t(p.x), fp_t(p.y)};
  }

  // Returns (x, y) rotated t radians counter-clockwise about point p.
  TPoint<fp_t> rotate_ccw(const TPoint &p, fp_t t) const {
    return TPoint<fp_t>{fp_t(x) - p.x, fp_t(y) - p.y}.rotate_ccw(t) +
           TPoint<fp_t>{fp_t(p.x), fp_t(p.y)};
  }

  // --- Reflections ---

  // Returns (x, y) reflected across point p. Exact for any coordinate type.
  TPoint reflect(const TPoint &p) const { return {2 * p.x - x, 2 * p.y - y}; }

  // Returns (x, y) reflected across the line containing points p and q.
  // Always returns floating-point coordinates.
  TPoint<fp_t> reflect(const TPoint &p, const TPoint &q) const {
    TPoint<fp_t> fp{fp_t(p.x), fp_t(p.y)};
    if (EQ(p, q)) {
      return TPoint<fp_t>{fp_t(x), fp_t(y)}.reflect(fp);
    }
    TPoint<fp_t> r{fp_t(x) - p.x, fp_t(y) - p.y};
    TPoint<fp_t> s{fp_t(q.x) - p.x, fp_t(q.y) - p.y};
    fp_t ssq = s.sqnorm();
    r = TPoint<fp_t>{(r.x * s.x + r.y * s.y) / ssq, (r.x * s.y - r.y * s.x) / ssq};
    return TPoint<fp_t>{r.x * s.x - r.y * s.y + fp.x, r.x * s.y + r.y * s.x + fp.y};
  }

  // --- Explicit type conversions ---

  TPoint<double> to_double() const { return {static_cast<double>(x), static_cast<double>(y)}; }

  TPoint<long double> to_ldouble() const {
    return {static_cast<long double>(x), static_cast<long double>(y)};
  }

  // --- Friend free-function versions ---

  friend T sqnorm(const TPoint &p) { return p.sqnorm(); }
  friend fp_t norm(const TPoint &p) { return p.norm(); }
  friend fp_t arg(const TPoint &p) { return p.arg(); }
  friend T dot(const TPoint &p, const TPoint &q) { return p.dot(q); }
  friend T cross(const TPoint &p, const TPoint &q) { return p.cross(q); }
  friend fp_t proj(const TPoint &p, const TPoint &q) { return p.proj(q); }
  friend TPoint<fp_t> normalize(const TPoint &p) { return p.normalize(); }
  friend TPoint rotate90(const TPoint &p) { return p.rotate90(); }
  friend TPoint rotate180(const TPoint &p) { return p.rotate180(); }
  friend TPoint rotate270(const TPoint &p) { return p.rotate270(); }
  friend TPoint rotate90(const TPoint &p, const TPoint &q) { return p.rotate90(q); }
  friend TPoint rotate180(const TPoint &p, const TPoint &q) { return p.rotate180(q); }
  friend TPoint rotate270(const TPoint &p, const TPoint &q) { return p.rotate270(q); }
  friend TPoint<fp_t> rotate_cw(const TPoint &p, fp_t t) { return p.rotate_cw(t); }
  friend TPoint<fp_t> rotate_ccw(const TPoint &p, fp_t t) { return p.rotate_ccw(t); }
  friend TPoint<fp_t> rotate_cw(const TPoint &p, const TPoint &q, fp_t t) { return p.rotate_cw(q, t); }
  friend TPoint<fp_t> rotate_ccw(const TPoint &p, const TPoint &q, fp_t t) { return p.rotate_ccw(q, t); }
  friend TPoint reflect(const TPoint &p, const TPoint &q) { return p.reflect(q); }
  friend TPoint<fp_t> reflect(const TPoint &p, const TPoint &a, const TPoint &b) { return p.reflect(a, b); }

  friend std::ostream &operator<<(std::ostream &out, const TPoint &p) {
    if constexpr (std::is_floating_point<T>::value) {
      return out << "(" << (std::fabs(p.x) < EPS ? 0 : p.x) << ","
                 << (std::fabs(p.y) < EPS ? 0 : p.y) << ")";
    }
    return out << "(" << p.x << "," << p.y << ")";
  }
};

using PointI = TPoint<int>;
using PointL = TPoint<long long>;
using PointD = TPoint<double>;
using PointLD = TPoint<long double>;
using Point = PointD;  // Default point type is double.

// Can compose with numerical types from chapter 6:
// using PointB = TPoint<BigInt>;
// using PointR = TPoint<Rational<int64_t>>;
// using PointM = TPoint<Modular<1000000007>>;
\end{lstlisting}
\Needspace*{4\baselineskip}
\textbf{Example Usage}\par
\begin{lstlisting}[style=cpp]
#include <cassert>
using namespace std;

const double PI = acos(-1.0);

int main() {
  Point p(-10, 3);
  assert(EQ(Point(-18, 29), p + Point(-3, 9) * 6.0 / 2.0 - Point(-1, 1)));
  assert(EQ(109, p.sqnorm()));
  assert(EQ(10.44030650891, p.norm()));
  assert(EQ(2.850135859112, p.arg()));
  assert(EQ(0, p.dot(Point(3, 10))));
  assert(EQ(0, p.cross(Point(10, -3))));
  assert(EQ(10, p.proj(Point(-10, 0))));
  assert(EQ(1, p.normalize().norm()));
  assert(EQ(Point(-3, -10), p.rotate90()));
  assert(EQ(Point(10, -3), p.rotate180()));
  assert(EQ(Point(3, 10), p.rotate270()));
  assert(EQ(Point(3, 12), p.rotate_cw(Point(1, 1), PI / 2)));
  assert(EQ(Point(1, -10), p.rotate_ccw(Point(2, 2), PI / 2)));
  assert(EQ(Point(10, -3), p.reflect(Point(0, 0))));
  assert(EQ(Point(-10, -3), p.reflect(Point(-2, 0), Point(5, 0))));

  // Integer point - exact arithmetic, float-only ops return PointD.
  PointI a(3, 4), b(1, 0);
  assert(a + b == PointI(4, 4));
  assert(a - b == PointI(2, 4));
  assert(a * 2 == PointI(6, 8));
  assert(a.sqnorm() == 25);
  assert(a.dot(b) == 3);
  assert(a.cross(b) == -4);
  assert(EQ(a.norm(), 5.0));
  assert(a.rotate90() == PointI(-4, 3));
  assert(a.rotate180() == PointI(-3, -4));
  assert(a.rotate270() == PointI(4, -3));
  PointD anorm = a.normalize();  // returns PointD
  assert(EQ(anorm.norm(), 1.0));
  PointD adiv = a / 2.0;  // returns PointD (division always promotes)
  assert(EQ(adiv.x, 1.5) && EQ(adiv.y, 2.0));
  auto rotated = PointI(1, 0).rotate_ccw(PI / 4);  // returns PointD (arbitrary rotation promotes)
  double root2over2 = sqrt(2.0) / 2.0;
  assert(EQ(rotated.x, root2over2) && EQ(rotated.y, root2over2));

  // Implicit conversion PointI -> PointD.
  PointD promoted = PointI(1, 2);
  assert(promoted == PointD(1, 2));
  return 0;
}
\end{lstlisting}
\subsection{Line}
\label{sec:7.1.2}
A straight line in two dimensions, represented by coefficients of $ax + by + c = 0$ and templated on their type \inlinecode{T}. Coefficients are canonicalized by default so proportional valid coefficient vectors have the same representation, up to floating-point rounding.

The template flag \inlinecode{CANONICALIZE} controls whether construction normalizes the coefficients and defaults to \inlinecode{true}. When enabled, \inlinecode{EXACT = true} divides coefficients by their Euclidean GCD and fixes their sign, while \inlinecode{EXACT = false} scales a nonvertical line to $b = 1$ and a vertical line to $a = 1$. \inlinecode{EXACT} defaults to \inlinecode{true} for built-in integral types; set it explicitly for custom exact types such as \inlinecode{BigInt} and \inlinecode{Rational}. Exact canonicalization requires \inlinecode{\%}, comparisons, and basic arithmetic, and predicates remain exact for integral \inlinecode{T}. Floating-point and \inlinecode{Modular} coefficients use the scaling form.

\begin{compactitem}
  \item \inlinecode{TLine\textless{}T\textgreater{}(a, b, c)} constructs a line from its coefficients, while \inlinecode{TLine\textless{}T\textgreater{}(p, q)} constructs the line through distinct points \inlinecode{p} and \inlinecode{q}. For floating-point \inlinecode{T}, \inlinecode{TLine\textless{}T\textgreater{}(slope, p)} constructs the line of the given slope through \inlinecode{p}. The default constructor produces an invalid line.
  \item \inlinecode{canonicalized()} returns a canonical copy. \inlinecode{operator==} compares stored coefficients exactly, while \inlinecode{EQ(l1, l2)} canonicalizes raw lines if necessary and uses \inlinecode{EPS} for floating-point coefficients.
  \item \inlinecode{valid()}, \inlinecode{horizontal()}, and \inlinecode{vertical()} classify the line.
  \item \inlinecode{slope()}, \inlinecode{x(y)}, and \inlinecode{y(x)} return \inlinecode{fp\_t}, which is \inlinecode{T} for floating-point coefficients and \inlinecode{double} otherwise. An undefined result is returned as \inlinecode{NaN}.
  \item \inlinecode{contains(p)}, \inlinecode{is\_parallel(l)}, and \inlinecode{is\_perpendicular(l)} test geometric relationships.
  \item \inlinecode{parallel(p)} and \inlinecode{perpendicular(p)} return the corresponding line through point \inlinecode{p}.
  \item \inlinecode{operator\textless{}\textless{}} outputs the line in the form \inlinecode{ax+by+c=0}.
  \item \inlinecode{LineI}, \inlinecode{LineL}, \inlinecode{LineD}, and \inlinecode{LineLD} use \inlinecode{int}, \inlinecode{int64\_t}, \inlinecode{double}, and \inlinecode{long double} coefficients, respectively. \inlinecode{Line} aliases \inlinecode{LineD}.
\end{compactitem}

Overflow warning: the exact predicates form products of coefficients, which grow like the squared coordinate magnitude. With \inlinecode{LineI} these may overflow once coordinates reach a few tens of thousands, so use \inlinecode{LineL} for larger integer coordinates.

\Needspace*{3\baselineskip}
\textbf{Time Complexity}:
\begin{compactitem}
  \item $O(\log C)$ per integral canonicalizing constructor or call to \inlinecode{canonicalize()}, where $C$ is the largest absolute coefficient.
  \item $O(1)$ per call to all other operations.
\end{compactitem}

\Needspace*{3\baselineskip}
\textbf{Space Complexity}:
\begin{compactitem}
  \item $O(1)$ for storage of the line.
  \item $O(1)$ auxiliary for all operations.
\end{compactitem}

\begin{lstlisting}[style=cpp]
#include <cmath>
#include <cstdint>
#include <ios>
#include <limits>
#include <ostream>
#include <type_traits>
#include <utility>

const double EPS = 1e-9;
const double M_NAN = std::numeric_limits<double>::quiet_NaN();

// Epsilon-aware for floating-point coordinates; exact for int and for coordinate types like Modular
// or Rational, which therefore compose for all of the predicates below.
template<typename T, typename U, typename C = std::common_type_t<T, U>>
bool EQ(T a, U b) {
  if constexpr (std::is_floating_point_v<C>) return C(a) == C(b) || std::fabs(C(a) - C(b)) <= EPS;
  return C(a) == C(b);
}

// SFINAE guard: valid only when Pt exposes numeric .x/.y members, so templated point constructors
// don't hijack scalar-argument calls.
template<typename Pt>
using if_point = decltype(std::declval<const Pt &>().x, std::declval<const Pt &>().y, void());

template<typename T, bool CANONICALIZE = true, bool EXACT = std::is_integral<T>::value>
struct TLine {
  // Slope and solve operations preserve floating-point coordinates and otherwise promote to double.
  using fp_t = std::conditional_t<std::is_floating_point<T>::value, T, double>;

  T a, b, c;

  void canonicalize() {
    if constexpr (EXACT) {
      auto gcd = [](T x, T y) {
        if (x < T{0}) x = -x;
        if (y < T{0}) y = -y;
        while (y != T{0}) {
          x %= y;
          std::swap(x, y);
        }
        return x;
      };
      T g = gcd(a, gcd(b, c));
      if (g != 0) {
        a /= g;
        b /= g;
        c /= g;
      }
      if (a < 0 || (a == 0 && b < 0)) {
        a = -a;
        b = -b;
        c = -c;
      }
    } else if (!EQ(b, T{0})) {
      a /= b;
      c /= b;
      b = T{1};
    } else if (!EQ(a, T{0})) {
      c /= a;
      a = T{1};
      b = T{0};
    }
  }

  TLine() : a(0), b(0), c(0) {}  // Invalid or uninitialized line.

  TLine(T a, T b, T c) : a(a), b(b), c(c) {
    if constexpr (CANONICALIZE) canonicalize();
  }

  // Line through two points. Coefficients are exact for integral points (overflow warning for c).
  template<typename Pt, typename = if_point<Pt>>
  TLine(const Pt &p, const Pt &q) : a(q.y - p.y), b(p.x - q.x), c(-(a * p.x + b * p.y)) {
    if constexpr (CANONICALIZE) canonicalize();
  }

  // Line of a given slope through a point. Intended for floating-point T (a generic slope
  // produces non-integer coefficients), so this constructor is disabled for integer T.
  template<
      class Pt, typename = if_point<Pt>, typename U = T,
      class = std::enable_if_t<std::is_floating_point<U>::value>>
  TLine(fp_t slope, const Pt &p) : a(-slope), b(1), c(slope * p.x - p.y) {
    if constexpr (CANONICALIZE) canonicalize();
  }

  TLine<T, true, EXACT> canonicalized() const { return TLine<T, true, EXACT>(a, b, c); }

  bool operator==(const TLine &l) const { return a == l.a && b == l.b && c == l.c; }
  bool operator!=(const TLine &l) const { return !(*this == l); }

  friend bool EQ(const TLine &l1, const TLine &l2) {
    if constexpr (CANONICALIZE) {
      return EQ(l1.a, l2.a) && EQ(l1.b, l2.b) && EQ(l1.c, l2.c);
    }
    return EQ(l1.canonicalized(), l2.canonicalized());
  }

  // Returns whether the line is a valid line (not both a and b zero).
  bool valid() const { return !(EQ(a, 0) && EQ(b, 0)); }
  bool horizontal() const { return valid() && EQ(a, 0); }
  bool vertical() const { return valid() && EQ(b, 0); }

  // Slope -a/b, or NaN if the line is vertical or invalid.
  fp_t slope() const { return (!valid() || EQ(b, 0)) ? M_NAN : -fp_t(a) / b; }

  // Solve for x at a given y (NaN if horizontal or invalid).
  fp_t x(fp_t y) const { return (!valid() || EQ(a, 0)) ? M_NAN : -(fp_t(b) * y + c) / a; }

  // Solve for y at a given x (NaN if vertical or invalid).
  fp_t y(fp_t x) const { return (!valid() || EQ(b, 0)) ? M_NAN : -(fp_t(a) * x + c) / b; }

  // Whether point p lies on the line. Exact for integer T and integer p.
  template<typename Pt>
  bool contains(const Pt &p) const {
    return valid() && EQ(a * p.x + b * p.y + c, 0);  // Overflow warning.
  }

  // Parallel iff the normals are parallel (cross == 0); perpendicular iff normals are perpendicular
  // (dot == 0). Both exact for integer T. Overflow warning.
  bool is_parallel(const TLine &l) const { return valid() && l.valid() && EQ(a * l.b, l.a * b); }

  bool is_perpendicular(const TLine &l) const {
    return valid() && l.valid() && EQ(a * l.a, -(b * l.b));
  }

  // Parallel/perpendicular line through point p. Exact for integer T and integer p.
  template<typename Pt, typename = if_point<Pt>>
  TLine parallel(const Pt &p) const {
    return TLine(a, b, -(a * p.x + b * p.y));  // Overflow warning.
  }

  template<typename Pt, typename = if_point<Pt>>
  TLine perpendicular(const Pt &p) const {
    return TLine(-b, a, b * p.x - a * p.y);  // Overflow warning.
  }

  friend std::ostream &operator<<(std::ostream &out, const TLine &l) {
    auto clean = [](T v) {
      if constexpr (std::is_floating_point_v<T>) {
        return std::fabs(v) < EPS ? T{0} : v;
      }
      return v;
    };
    auto flags = out.flags();
    out << clean(l.a) << "x" << std::showpos << clean(l.b) << "y" << clean(l.c) << "=0";
    out.flags(flags);
    return out;
  }
};

using LineI = TLine<int>;
using LineL = TLine<int64_t>;
using LineD = TLine<double>;
using LineLD = TLine<long double>;
using Line = LineD;  // Default line type is double.

// Can compose with numerical types from chapter 6:
// using LineB = TLine<BigInt, true, true>;
// using LineR = TLine<Rational<int64_t>, true, true>;
\end{lstlisting}
\Needspace*{4\baselineskip}
\textbf{Example Usage}\par
\begin{lstlisting}[style=cpp]
#include <cassert>

struct Point {
  double x, y;
  Point(double x = 0, double y = 0) : x(x), y(y) {}
};

struct PointI {
  int x, y;
  PointI(int x = 0, int y = 0) : x(x), y(y) {}
};

int main() {
  Line invalid;
  assert(!invalid.contains(Point(0, 0)) && !invalid.is_parallel(invalid));

  // Floating-point line.
  Line l(2, -5, -8);
  Line para = Line(2, -5, -8).parallel(Point(-6, -2));
  Line perp = Line(2, -5, -8).perpendicular(Point(-6, -2));
  assert(l.is_parallel(para) && l.is_perpendicular(perp));
  assert(l.slope() == 0.4);
  assert(EQ(para, Line(-0.4, 1, -0.4)));
  assert(EQ(perp, Line(2.5, 1, 17)));
  assert(l.contains(Point(1.5, -1)));  // 2(1.5) - 5(-1) - 8 = 0.

  // Integer line through integer points - exact canonical coefficients and predicates.
  LineI il(PointI(0, 0), PointI(2, 1));  // a=1, b=-2, c=0  ->  x - 2y = 0
  assert(il.a == 1 && il.b == -2 && il.c == 0);
  assert(il.contains(PointI(4, 2)));  // exact integer check
  assert(!il.contains(PointI(4, 3)));
  assert(il.is_parallel(LineI(PointI(1, 1), PointI(3, 2))));        // both slope 1/2
  assert(il.is_perpendicular(LineI(PointI(0, 0), PointI(1, -2))));  // dot of normals = 0
  assert(il.slope() == 0.5);                                        // fp_t result

  LineI reduced(2, -4, -6);
  assert(reduced == LineI(1, -2, -3));

  using RawLine = TLine<double, false>;
  RawLine raw(2, -5, -8), scaled(-0.4, 1, 1.6);
  assert(raw != scaled);    // Stored coefficients differ.
  assert(EQ(raw, scaled));  // Canonicalized copies represent the same line.
  RawLine canonical = raw;
  canonical.canonicalize();
  assert(canonical == scaled);
  assert(EQ(raw.canonicalized(), Line(-0.4, 1, 1.6)));
  return 0;
}
\end{lstlisting}
\subsection{Circle}
\label{sec:7.1.3}
A circle in two dimensions with epsilon-aware operations. The circle centered at $(h, k)$ is represented by the relation $(x - h)^2 + (y - k)^2 = r^2$, where the radius $r$ is normalized to a nonnegative number.

This implementation stores and computes circle parameters in \inlinecode{double}. Point-accepting constructors and predicates are templated on the point type \inlinecode{Pt} and only read \inlinecode{.x}/\inlinecode{.y}, so they accept \inlinecode{Point}, \inlinecode{PointD}, or \inlinecode{PointI} from \secref{7.1.1} or any struct with numeric \inlinecode{.x} and \inlinecode{.y} fields.

\begin{compactitem}
  \item \inlinecode{Circle()} constructs the zero-radius circle centered at the origin.
  \item \inlinecode{Circle(r)} constructs a circle of radius \inlinecode{abs(r)} centered at the origin.
  \item \inlinecode{Circle(h, k, r)} constructs a circle of radius \inlinecode{abs(r)} centered at (\inlinecode{h}, \inlinecode{k}).
  \item \inlinecode{Circle(o, r)} constructs a circle of radius \inlinecode{abs(r)} centered at point \inlinecode{o}.
  \item \inlinecode{Circle(a, b)} constructs the circle whose diameter is segment \inlinecode{a}-\inlinecode{b}.
  \item \inlinecode{Circle(a, b, c)} constructs the circumcircle through non-collinear points \inlinecode{a}, \inlinecode{b}, and \inlinecode{c}, throwing if the points are collinear.
  \item \inlinecode{Circle(a, b, r)} constructs one circle of radius \inlinecode{abs(r)} passing through points \inlinecode{a} and \inlinecode{b}, throwing if no finite circle is determined by the inputs.
  \item \inlinecode{operator==} and \inlinecode{operator!=} compare stored centers and radii exactly; \inlinecode{EQ(a, b)} compares them using \inlinecode{EPS}.
  \item \inlinecode{contains(p)} returns whether point \inlinecode{p} lies inside or on the circle.
  \item \inlinecode{on\_boundary(p)} returns whether point \inlinecode{p} lies on the circle boundary.
  \item \inlinecode{incircle(a, b, c)} returns the circle inscribed in triangle \inlinecode{abc}.
  \item \inlinecode{in\_circumcircle(a, b, c, d)} returns $1$ if point \inlinecode{d} lies strictly inside the circle through points \inlinecode{a}, \inlinecode{b}, and \inlinecode{c}, $0$ if it lies on the circle, or $-1$ if it lies outside. This has exact reasoning for integral points by evaluating the determinant directly without constructing a \inlinecode{Circle} or taking square roots. Points \inlinecode{a}, \inlinecode{b}, and \inlinecode{c} must not be collinear. For integer inputs, the sign test is exact provided the chosen intermediate type does not overflow.
\end{compactitem}

Overflow warning: the determinant in \inlinecode{in\_circumcircle()} grows like the fourth power of the coordinate magnitude. For large integer coordinates, use a point or intermediate type wide enough for that determinant.

\textbf{Time Complexity}: $O(1)$ per call to the constructors and all other operations.

\Needspace*{3\baselineskip}
\textbf{Space Complexity}:
\begin{compactitem}
  \item $O(1)$ for storage of the circle.
  \item $O(1)$ auxiliary for all operations.
\end{compactitem}

\begin{lstlisting}[style=cpp]
#include <cmath>
#include <cstdint>
#include <ios>
#include <ostream>
#include <stdexcept>
#include <type_traits>

const double EPS = 1e-9;

template<typename T, typename U, typename C = std::common_type_t<T, U>>
bool EQ(T a, U b) {
  if constexpr (std::is_floating_point_v<C>) return C(a) == C(b) || std::fabs(C(a) - C(b)) <= EPS;
  return C(a) == C(b);
}

template<typename T, typename U, typename C = std::common_type_t<T, U>>
bool LT(T a, U b) {
  if constexpr (std::is_floating_point_v<C>) return C(a) < C(b) - EPS;
  return C(a) < C(b);
}

template<typename T, typename U>
bool LE(T a, U b) {
  return !LT(b, a);
}

// SFINAE guard: valid only when Pt exposes numeric .x/.y members. This keeps the templated point
// constructors from hijacking calls like Circle(double, double, double).
template<typename Pt>
using if_point = decltype(std::declval<const Pt &>().x, std::declval<const Pt &>().y, void());

struct Circle {
  double h, k, r;

  Circle() : h(0), k(0), r(0) {}
  explicit Circle(double r) : h(0), k(0), r(fabs(r)) {}
  Circle(double h, double k, double r) : h(h), k(k), r(fabs(r)) {}

  template<typename Pt, typename = if_point<Pt>>
  Circle(const Pt &o, double r) : h(o.x), k(o.y), r(fabs(r)) {}

  template<typename Pt, typename = if_point<Pt>>
  Circle(const Pt &a, const Pt &b) {
    h = (a.x + b.x) / 2.0;
    k = (a.y + b.y) / 2.0;
    r = std::hypot(a.x - h, a.y - k);
  }

  template<typename Pt, typename = if_point<Pt>>
  Circle(const Pt &a, const Pt &b, const Pt &c) {
    double ax = a.x, ay = a.y, bx = b.x, by = b.y, cx = c.x, cy = c.y;
    double asq = ax * ax + ay * ay, bsq = bx * bx + by * by, csq = cx * cx + cy * cy;
    double d = 2.0 * (ax * (by - cy) + bx * (cy - ay) + cx * (ay - by));
    if (EQ(d, 0)) {
      throw std::runtime_error("No circumcircle from collinear points.");
    }
    h = (asq * (by - cy) + bsq * (cy - ay) + csq * (ay - by)) / d;
    k = (asq * (cx - bx) + bsq * (ax - cx) + csq * (bx - ax)) / d;
    r = std::hypot(ax - h, ay - k);
  }

  template<typename Pt, typename = if_point<Pt>>
  Circle(const Pt &a, const Pt &b, double r) : r(fabs(r)) {
    if (EQ(this->r, 0) && EQ(a.x, b.x) && EQ(a.y, b.y)) {
      h = a.x;
      k = a.y;
      return;
    }
    double d = std::hypot(b.x - a.x, b.y - a.y);
    if (EQ(d, 0)) {
      throw std::runtime_error("Identical points, infinite circles.");
    }
    if (LT(this->r * 2.0, d)) {
      throw std::runtime_error("Points too far away to make Circle.");
    }
    double z = this->r * this->r - d * d / 4.0;
    double v = std::sqrt(z < 0 ? 0 : z) / d;
    double mx = (a.x + b.x) / 2.0, my = (a.y + b.y) / 2.0;
    h = mx + v * (a.y - b.y);
    k = my + v * (b.x - a.x);
    // The other answer is (h, k) = (mx - v*(a.y - b.y), my - v*(b.x - a.x)).
  }

  bool operator==(const Circle &c) const { return h == c.h && k == c.k && r == c.r; }
  bool operator!=(const Circle &c) const { return !(*this == c); }

  friend bool EQ(const Circle &a, const Circle &b) {
    return EQ(a.h, b.h) && EQ(a.k, b.k) && EQ(a.r, b.r);
  }

  template<typename Pt>
  bool contains(const Pt &p) const {
    return LE(std::hypot(p.x - h, p.y - k), r);
  }

  template<typename Pt>
  bool on_boundary(const Pt &p) const {
    return EQ(std::hypot(p.x - h, p.y - k), r);
  }

  friend std::ostream &operator<<(std::ostream &out, const Circle &c) {
    auto flags = out.flags();
    out << std::showpos << "(x" << -(fabs(c.h) < EPS ? 0 : c.h) << ")^2+"
        << "(y" << -(fabs(c.k) < EPS ? 0 : c.k) << ")^2" << std::noshowpos << "="
        << (fabs(c.r) < EPS ? 0 : c.r * c.r);
    out.flags(flags);
    return out;
  }
};

template<typename Pt>
Circle incircle(const Pt &a, const Pt &b, const Pt &c) {
  double al = std::hypot(b.x - c.x, b.y - c.y);
  double bl = std::hypot(a.x - c.x, a.y - c.y);
  double cl = std::hypot(a.x - b.x, a.y - b.y);
  double l = al + bl + cl;
  double px = a.x - c.x, py = a.y - c.y, qx = b.x - c.x, qy = b.y - c.y;
  return EQ(l, 0) ? Circle(a.x, a.y, 0)
                  : Circle(
                        (al * a.x + bl * b.x + cl * c.x) / l, (al * a.y + bl * b.y + cl * c.y) / l,
                        fabs(px * qy - py * qx) / l
                    );
}

template<typename Pt>
int in_circumcircle(const Pt &a, const Pt &b, const Pt &c, const Pt &d) {
  using T = decltype(a.x + a.x);
  using W = std::conditional_t<
      std::is_floating_point<T>::value, long double,
      std::conditional_t<std::is_integral<T>::value, int64_t, T>>;
  W adx = (W)a.x - d.x, ady = (W)a.y - d.y;
  W bdx = (W)b.x - d.x, bdy = (W)b.y - d.y;
  W cdx = (W)c.x - d.x, cdy = (W)c.y - d.y;
  // Overflow warning.
  W det = (adx * adx + ady * ady) * (bdx * cdy - cdx * bdy) +
          (bdx * bdx + bdy * bdy) * (cdx * ady - adx * cdy) +
          (cdx * cdx + cdy * cdy) * (adx * bdy - bdx * ady);
  W orient = ((W)b.x - a.x) * ((W)c.y - a.y) - ((W)b.y - a.y) * ((W)c.x - a.x);
  W val = W{0} < orient ? det : -det;
  return EQ(val, W{0}) ? 0 : (W{0} < val ? 1 : -1);
}
\end{lstlisting}
\Needspace*{4\baselineskip}
\textbf{Example Usage}\par
\begin{lstlisting}[style=cpp]
#include <cassert>

struct Point {
  double x, y;
  Point(double x = 0, double y = 0) : x(x), y(y) {}
};

struct PointI {
  int x, y;
  PointI(int x = 0, int y = 0) : x(x), y(y) {}
};

int main() {
  Circle c(-2, 5, sqrt(10));
  assert(EQ(c, Circle(Point(-2, 5), sqrt(10))));
  assert(EQ(c, Circle(Point(1, 6), Point(-5, 4))));
  assert(EQ(c, Circle(Point(-3, 2), Point(-3, 8), Point(-1, 8))));
  assert(EQ(Circle(Point(0, 0), Point(0, 2), -1), Circle(0, 1, 1)));
  assert(EQ(c, incircle(Point(-12, 5), Point(3, 0), Point(0, 9))));
  assert(c.contains(Point(-2, 8)) && !c.contains(Point(-2, 9)));
  assert(c.on_boundary(Point(-1, 2)) && !c.on_boundary(Point(-1.01, 2)));

  // Integer-coordinate points are accepted; the Circle is computed in double.
  assert(EQ(c, Circle(PointI(1, 6), PointI(-5, 4))));
  assert(EQ(c, Circle(PointI(-3, 2), PointI(-3, 8), PointI(-1, 8))));
  assert(EQ(c, incircle(PointI(-12, 5), PointI(3, 0), PointI(0, 9))));
  assert(c.contains(PointI(-2, 8)) && !c.contains(PointI(-2, 9)));
  assert(c.on_boundary(PointI(-1, 2)));

  // Exact integer in-circle predicate: unit circle through (1,0), (0,1), (-1,0).
  assert(in_circumcircle(PointI(1, 0), PointI(0, 1), PointI(-1, 0), PointI(0, 0)) == 1);   // inside
  assert(in_circumcircle(PointI(1, 0), PointI(0, 1), PointI(-1, 0), PointI(2, 0)) == -1);  // out
  assert(in_circumcircle(PointI(1, 0), PointI(0, 1), PointI(-1, 0), PointI(0, -1)) == 0);  // edge

  // Orientation-independent: reversing a, b, c gives the same answer.
  assert(in_circumcircle(PointI(-1, 0), PointI(0, 1), PointI(1, 0), PointI(0, 0)) == 1);
  return 0;
}
\end{lstlisting}
\subsection{Triangle}
\label{sec:7.1.4}
Common triangle calculations in two dimensions. The functions are templated on the point type \inlinecode{Pt}, which should accept \inlinecode{Point}, \inlinecode{PointD}, or \inlinecode{PointI} from \secref{7.1.1} or any struct with numeric \inlinecode{.x} and \inlinecode{.y} fields. \inlinecode{triangle\_area()} returns \inlinecode{double} regardless of input type (the area is fractional via the \inlinecode{/2}). \inlinecode{same\_side()} and \inlinecode{point\_in\_triangle()} do all arithmetic in the point's own coordinate type and are exact for integer points: they reduce each cross product to a sign before combining, so no precision is lost and cross products are never multiplied together.

\begin{compactitem}
  \item \inlinecode{triangle\_area(a, b, c)} returns the area of the triangle with vertices \inlinecode{a}, \inlinecode{b}, and \inlinecode{c}.
  \item \inlinecode{triangle\_area\_sides(s1, s2, s3)} returns the area of a triangle with side lengths \inlinecode{s1}, \inlinecode{s2}, and \inlinecode{s3}. The given lengths must be nonnegative and form a valid triangle.
  \item \inlinecode{triangle\_area\_medians(m1, m2, m3)} returns the area of a triangle with medians of lengths \inlinecode{m1}, \inlinecode{m2}, and \inlinecode{m3}. The median of a triangle is a line segment joining a vertex to the midpoint of the opposing edge.
  \item \inlinecode{triangle\_area\_altitudes(h1, h2, h3)} returns the area of a triangle with altitudes \inlinecode{h1}, \inlinecode{h2}, and \inlinecode{h3}. An altitude of a triangle is the shortest line between a vertex and the infinite line that is extended from its opposite edge.
  \item \inlinecode{same\_side(p1, p2, a, b, include\_boundary = true)} returns whether points \inlinecode{p1} and \inlinecode{p2} lie on the same side of the line containing points \inlinecode{a} and \inlinecode{b}. If one or both points lie exactly on the line, then the result will depend on \inlinecode{include\_boundary}.
  \item \inlinecode{point\_in\_triangle(p, a, b, c, include\_boundary = true)} returns whether point \inlinecode{p} lies within the triangle with vertices \inlinecode{a}, \inlinecode{b}, and \inlinecode{c}. If the point lies on or close to an edge (by roughly \inlinecode{EPS}), then the result will depend on \inlinecode{include\_boundary}.
\end{compactitem}

Use an integral point type for integer coordinates, and a floating-point type for fractional coordinates.

Overflow warning: each individual cross product is still on the order of the squared coordinate magnitude, so for integer point types use a 64-bit coordinate type (e.g. \inlinecode{PointL} from \secref{7.1.1}) once coordinates reach a few tens of thousands.

\textbf{Time Complexity}: $O(1)$ per operation.

\textbf{Space Complexity}: $O(1)$ auxiliary for all operations.

\begin{lstlisting}[style=cpp]
#include <cmath>
#include <type_traits>

const double EPS = 1e-9;

template<typename T, typename U, typename C = std::common_type_t<T, U>>
bool EQ(T a, U b) {
  if constexpr (std::is_floating_point_v<C>) return C(a) == C(b) || std::fabs(C(a) - C(b)) <= EPS;
  return C(a) == C(b);
}

template<typename T, typename U, typename C = std::common_type_t<T, U>>
bool LT(T a, U b) {
  if constexpr (std::is_floating_point_v<C>) return C(a) < C(b) - EPS;
  return C(a) < C(b);
}

template<typename Pt>
double triangle_area(const Pt &a, const Pt &b, const Pt &c) {
  double acx = static_cast<double>(a.x) - c.x, acy = static_cast<double>(a.y) - c.y;
  double bcx = static_cast<double>(b.x) - c.x, bcy = static_cast<double>(b.y) - c.y;
  return fabs(acx * bcy - acy * bcx) / 2.0;
}

double triangle_area_sides(double s1, double s2, double s3) {
  double s = (s1 + s2 + s3) / 2.0;
  return std::sqrt(s * (s - s1) * (s - s2) * (s - s3));
}

double triangle_area_medians(double m1, double m2, double m3) {
  return 4.0 * triangle_area_sides(m1, m2, m3) / 3.0;
}

double triangle_area_altitudes(double h1, double h2, double h3) {
  if (EQ(h1, 0) || EQ(h2, 0) || EQ(h3, 0)) {
    return 0;
  }
  double x = h1 * h1, y = h2 * h2, z = h3 * h3;
  double v = 2.0 / (x * y) + 2.0 / (x * z) + 2.0 / (y * z);
  return 1.0 / std::sqrt(v - 1.0 / (x * x) - 1.0 / (y * y) - 1.0 / (z * z));
}

template<typename Pt>
bool same_side(const Pt &p1, const Pt &p2, const Pt &a, const Pt &b, bool include_boundary = true) {
  // Cross products in the point's native coordinate type (exact for integer points).
  auto c1 = (b.x - a.x) * (p1.y - a.y) - (b.y - a.y) * (p1.x - a.x);
  auto c2 = (b.x - a.x) * (p2.y - a.y) - (b.y - a.y) * (p2.x - a.x);
  // Compare the signs, not the product c1 * c2, to avoid integer overflow.
  int s1 = LT(0, c1) ? 1 : (LT(c1, 0) ? -1 : 0);
  int s2 = LT(0, c2) ? 1 : (LT(c2, 0) ? -1 : 0);
  return include_boundary ? (s1 * s2 >= 0) : (s1 * s2 > 0);
}

template<typename Pt>
bool point_in_triangle(
    const Pt &p, const Pt &a, const Pt &b, const Pt &c, bool include_boundary = true
) {
  return same_side(p, a, b, c, include_boundary) && same_side(p, b, a, c, include_boundary) &&
         same_side(p, c, a, b, include_boundary);
}
\end{lstlisting}
\Needspace*{4\baselineskip}
\textbf{Example Usage}\par
\begin{lstlisting}[style=cpp]
#include <cassert>

struct Point {
  double x, y;
  Point(double x = 0, double y = 0) : x(x), y(y) {}
};

struct PointI {
  int x, y;
  PointI(int x = 0, int y = 0) : x(x), y(y) {}
};

int main() {
  assert(EQ(6, triangle_area(Point(0, -1), Point(4, -1), Point(0, -4))));
  assert(EQ(6, triangle_area_sides(3, 4, 5)));
  assert(EQ(6, triangle_area_medians(3.605551275, 2.5, 4.272001873)));
  assert(EQ(6, triangle_area_altitudes(3, 4, 2.4)));

  // Fractional coordinates require a floating-point point type.
  assert(point_in_triangle(Point(0, 0), Point(-1, 0), Point(0, -2), Point(4, 0)));
  assert(!point_in_triangle(Point(0, 1), Point(-1, 0), Point(0, -2), Point(4, 0)));
  assert(point_in_triangle(Point(-2.44, 0.82), Point(-1, 0), Point(-3, 1), Point(4, 0)));
  assert(!point_in_triangle(Point(-2.44, 0.7), Point(-1, 0), Point(-3, 1), Point(4, 0)));

  // Integer coordinates: same_side / point_in_triangle are computed exactly in int.
  assert(EQ(8, triangle_area(PointI(0, 0), PointI(4, 0), PointI(0, 4))));
  assert(point_in_triangle(PointI(1, 1), PointI(0, 0), PointI(4, 0), PointI(0, 4)));
  assert(!point_in_triangle(PointI(5, 5), PointI(0, 0), PointI(4, 0), PointI(0, 4)));
  assert(point_in_triangle(PointI(2, 2), PointI(0, 0), PointI(4, 0), PointI(0, 4)));  // on edge
  assert(!point_in_triangle(PointI(2, 2), PointI(0, 0), PointI(4, 0), PointI(0, 4), false));
  return 0;
}
\end{lstlisting}
\subsection{Rectangle}
\label{sec:7.1.5}
Common axis-aligned rectangle calculations in two dimensions. The functions are templated on the point type \inlinecode{Pt}, which should accept \inlinecode{Point}, \inlinecode{PointD}, or \inlinecode{PointI} from \secref{7.1.1}, or any struct with numeric \inlinecode{.x} and \inlinecode{.y} fields.

\begin{compactitem}
  \item \inlinecode{point\_in\_rect(p, v, w, h, include\_boundary = true)} returns whether \inlinecode{p} is inside the axis-aligned rectangle whose opposite corners are \inlinecode{v} and (\inlinecode{v.x + w}, \inlinecode{v.y + h}). Negative widths and heights are supported. The \inlinecode{include\_boundary} flag controls whether boundary points count.
  \item \inlinecode{point\_in\_rect(p, a, b, include\_boundary = true)} returns whether \inlinecode{p} is inside the axis-aligned rectangle with opposite corners \inlinecode{a} and \inlinecode{b}.
  \item \inlinecode{rect\_intersection(a1, b1, a2, b2, \&p, \&q, include\_boundary = true)} finds the intersection of two axis-aligned rectangles, where \inlinecode{a1}/\inlinecode{b1} and \inlinecode{a2}/\inlinecode{b2} are opposite-corner pairs. Returns $-1$ if the rectangles are disjoint, $0$ if they partially intersect, $1$ if the first rectangle is completely inside the second, and $2$ if the second rectangle is completely inside the first. If an intersection exists, its lower-left and upper-right corners are stored in \inlinecode{p} and \inlinecode{q} when those pointers are non-null. The \inlinecode{include\_boundary} flag controls whether boundary-only contact counts as intersecting.
\end{compactitem}

Overflow warning: For integer-coordinate inputs, comparisons are exact as long as intermediate coordinate additions, subtractions, and min/max values do not overflow.

\textbf{Time Complexity}: $O(1)$ per operation.

\textbf{Space Complexity}: $O(1)$ auxiliary for all operations.

\begin{lstlisting}[style=cpp]
#include <algorithm>
#include <cmath>
#include <type_traits>

const double EPS = 1e-9;

template<typename T, typename U, typename C = std::common_type_t<T, U>>
bool EQ(T a, U b) {
  if constexpr (std::is_floating_point_v<C>) return C(a) == C(b) || std::fabs(C(a) - C(b)) <= EPS;
  return C(a) == C(b);
}

template<typename T, typename U, typename C = std::common_type_t<T, U>>
bool LT(T a, U b) {
  if constexpr (std::is_floating_point_v<C>) return C(a) < C(b) - EPS;
  return C(a) < C(b);
}

template<typename T, typename U> bool GT(T a, U b) { return LT(b, a); }
template<typename T, typename U> bool LE(T a, U b) { return !LT(b, a); }
template<typename T, typename U> bool GE(T a, U b) { return !LT(a, b); }

template<typename Pt, typename T>
bool point_in_rect(const Pt &p, const Pt &v, const T &w, const T &h, bool include_boundary = true) {
  if (w < 0) {
    return point_in_rect(p, Pt(v.x + w, v.y), -w, h, include_boundary);
  }
  if (h < 0) {
    return point_in_rect(p, Pt(v.x, v.y + h), w, -h, include_boundary);
  }
  return include_boundary ? (GE(p.x, v.x) && LE(p.x, v.x + w) && GE(p.y, v.y) && LE(p.y, v.y + h))
                          : (GT(p.x, v.x) && LT(p.x, v.x + w) && GT(p.y, v.y) && LT(p.y, v.y + h));
}

template<typename Pt>
bool point_in_rect(const Pt &p, const Pt &a, const Pt &b, bool include_boundary = true) {
  auto xl = std::min(a.x, b.x), yl = std::min(a.y, b.y);
  auto xh = std::max(a.x, b.x), yh = std::max(a.y, b.y);
  return point_in_rect(p, Pt(xl, yl), xh - xl, yh - yl, include_boundary);
}

template<typename Pt>
int rect_intersection(
    const Pt &a1, const Pt &b1, const Pt &a2, const Pt &b2, Pt *p = nullptr, Pt *q = nullptr,
    const bool include_boundary = true
) {
  // Normalize each rectangle to [xl, xh] x [yl, yh] with coordinate-wise min/max only.
  // (Using auto keeps the native coordinate type, so integer rectangles stay exact.)
  auto xl1 = std::min(a1.x, b1.x), xh1 = std::max(a1.x, b1.x);
  auto yl1 = std::min(a1.y, b1.y), yh1 = std::max(a1.y, b1.y);
  auto xl2 = std::min(a2.x, b2.x), xh2 = std::max(a2.x, b2.x);
  auto yl2 = std::min(a2.y, b2.y), yh2 = std::max(a2.y, b2.y);
  // The intersection is the overlap of the two coordinate intervals.
  auto ixl = std::max(xl1, xl2), ixh = std::min(xh1, xh2);
  auto iyl = std::max(yl1, yl2), iyh = std::min(yh1, yh2);
  if (include_boundary ? (GT(ixl, ixh) || GT(iyl, iyh)) : (GE(ixl, ixh) || GE(iyl, iyh))) {
    return -1;  // Completely disjoint.
  }
  // Opposing corners of the overlap: bottom-left in p, top-right in q.
  if (p != nullptr) {
    *p = Pt(ixl, iyl);
  }
  if (q != nullptr) {
    *q = Pt(ixh, iyh);
  }
  // The overlap equals a rectangle exactly when that rectangle is contained in the other.
  if (EQ(ixl, xl1) && EQ(ixh, xh1) && EQ(iyl, yl1) && EQ(iyh, yh1)) {
    return 1;  // Rectangle 1 completely inside 2.
  }
  if (EQ(ixl, xl2) && EQ(ixh, xh2) && EQ(iyl, yl2) && EQ(iyh, yh2)) {
    return 2;  // Rectangle 2 completely inside 1.
  }
  return 0;  // Partial intersection.
}
\end{lstlisting}
\Needspace*{4\baselineskip}
\textbf{Example Usage}\par
\begin{lstlisting}[style=cpp]
#include <cassert>

struct Point {
  double x, y;
  Point(double x = 0, double y = 0) : x(x), y(y) {}
  bool operator==(const Point &p) const { return x == p.x && y == p.y; }
};

int main() {
  assert(point_in_rect(Point(0, -1), Point(0, -3), 3, 2));
  assert(!point_in_rect(Point(0, -1), Point(0, -3), 3, 2, false));
  assert(point_in_rect(Point(2, -2), Point(3, -3), -3, 2));
  assert(!point_in_rect(Point(0, 0), Point(3, -1), -3, -2));
  assert(point_in_rect(Point(2, -2), Point(3, -3), Point(0, -1)));
  assert(!point_in_rect(Point(0, -1), Point(3, -3), Point(0, -1), false));
  assert(!point_in_rect(Point(-1, -2), Point(3, -3), Point(0, -1)));

  Point p, q;
  // Output coordinates come directly from input coordinates via min/max, so exact equality is safe.
  assert(-1 == rect_intersection(Point(0, 0), Point(1, 1), Point(2, 2), Point(3, 3)));
  assert(0 == rect_intersection(Point(1, 1), Point(7, 7), Point(5, 5), Point(0, 0), &p, &q));
  assert(p == Point(1, 1) && q == Point(5, 5));
  assert(1 == rect_intersection(Point(1, 1), Point(0, 0), Point(0, 0), Point(1, 10), &p, &q));
  assert(p == Point(0, 0) && q == Point(1, 1));
  assert(2 == rect_intersection(Point(0, 5), Point(5, 7), Point(1, 6), Point(2, 5), &p, &q));
  assert(p == Point(1, 5) && q == Point(2, 6));
  assert(
      -1 == rect_intersection(Point(0, 0), Point(1, 1), Point(1, 1), Point(2, 2), &p, &q, false)
  );
  return 0;
}
\end{lstlisting}

\section{\texorpdfstring{Angles, Distances, and Intersections}{7.2 Angles, Distances, and Intersections}}
\setcounter{section}{2}
\setcounter{subsection}{0}
\subsection{Angles}
\label{sec:7.2.1}
Angle calculations in two dimensions. All operations are inherently floating-point (\inlinecode{atan2}, \inlinecode{acos}, trigonometry), so these functions use the local double-coordinate \inlinecode{PointD} type. Convert integer points to \inlinecode{PointD} when asking for angles. The constants \inlinecode{DEG} and \inlinecode{RAD} may be used as multipliers to convert between degrees and radians. For example, if \inlinecode{t} is a value in degrees, then \inlinecode{t * DEG} is the equivalent angle in radians; if \inlinecode{t} is in radians, then \inlinecode{t * RAD} is the equivalent angle in degrees.

\begin{compactitem}
  \item \inlinecode{reduce\_deg(t)} takes an angle \inlinecode{t} degrees and returns an equivalent angle in the range $[0, 360)$ degrees (e.g. $-630$ becomes $90$).
  \item \inlinecode{reduce\_rad(t)} takes an angle \inlinecode{t} radians and returns an equivalent angle in the range $[0, 2\pi)$ radians (e.g. $720.5$ becomes $0.5$).
  \item \inlinecode{polar\_x(r, t)} and \inlinecode{polar\_y(r, t)} return the Cartesian coordinates of the point at radius \inlinecode{r} and angle \inlinecode{t} radians in polar coordinates (see \inlinecode{std::polar()}).
  \item \inlinecode{polar\_angle(p)} returns the angle in radians of the line segment from $(0, 0)$ to point \inlinecode{p}, relative counterclockwise to the positive $x$-axis.
  \item \inlinecode{angle(a, o, b)} returns the smallest angle in radians formed by the points \inlinecode{a}, \inlinecode{o}, \inlinecode{b} with vertex at point \inlinecode{o}. The points \inlinecode{a} and \inlinecode{b} must differ from \inlinecode{o}.
  \item \inlinecode{angle\_between(a, b)} returns the directed counter-clockwise angle in $[0, 2\pi)$ from vector \inlinecode{a} to vector \inlinecode{b}, treating both points as vectors from the origin. Both vectors must be nonzero.
  \item \inlinecode{angle\_between(a1, b1, a2, b2)} returns the smaller angle in radians between two lines whose normal vectors are (\inlinecode{a1}, \inlinecode{b1}) and (\inlinecode{a2}, \inlinecode{b2}), limited to $[0, \pi / 2]$. Both coefficient pairs must represent valid lines.
  \item \inlinecode{cross(a, b, o = Pt(0, 0))} returns the signed $z$-component of the three-dimensional cross product between points \inlinecode{a} and \inlinecode{b}, where the $z$-coordinates are implicitly zero and the origin is shifted to point \inlinecode{o}. This is also double the signed area of the triangle from these three points.
  \item \inlinecode{turn(a, o, b)} returns $1$ if the path \inlinecode{a} $\to$ \inlinecode{o} $\to$ \inlinecode{b} forms a left turn on the plane, $0$ if the path forms a straight line segment, or $-1$ if it forms a right turn.
\end{compactitem}

Overflow warning: \inlinecode{cross()} and \inlinecode{turn()} preserve the point's coordinate type, so their products may overflow for integral points. Use a 64-bit coordinate type for large integer coordinates.

\textbf{Time Complexity}: $O(1)$ per operation.

\textbf{Space Complexity}: $O(1)$ auxiliary for all operations.

\begin{lstlisting}[style=cpp]
#include <algorithm>
#include <cmath>
#include <type_traits>

const double EPS = 1e-9;
const double PI = std::acos(-1.0);
const double DEG = PI / 180;
const double RAD = 180 / PI;

template<typename T, typename U, typename C = std::common_type_t<T, U>>
bool EQ(T a, U b) {
  if constexpr (std::is_floating_point_v<C>) return C(a) == C(b) || std::fabs(C(a) - C(b)) <= EPS;
  return C(a) == C(b);
}

template<typename T, typename U, typename C = std::common_type_t<T, U>>
bool LT(T a, U b) {
  if constexpr (std::is_floating_point_v<C>) return C(a) < C(b) - EPS;
  return C(a) < C(b);
}

double reduce_deg(double t) {
  if (t < -360) {
    return reduce_deg(std::fmod(t, 360));
  }
  if (t < 0) {
    return t + 360;
  }
  return (t >= 360) ? std::fmod(t, 360) : t;
}

double reduce_rad(double t) {
  if (t < -2 * PI) {
    return reduce_rad(std::fmod(t, 2 * PI));
  }
  if (t < 0) {
    return t + 2 * PI;
  }
  return (t >= 2 * PI) ? std::fmod(t, 2 * PI) : t;
}

double polar_x(double r, double t) {
  return r * std::cos(t);
}

double polar_y(double r, double t) {
  return r * std::sin(t);
}

template<typename Pt>
double polar_angle(const Pt &p) {
  double t = std::atan2(static_cast<double>(p.y), static_cast<double>(p.x));
  return (t < 0) ? (t + 2 * PI) : t;
}

template<typename Pt>
double angle(const Pt &a, const Pt &o, const Pt &b) {
  double ux = o.x - a.x, uy = o.y - a.y;
  double vx = o.x - b.x, vy = o.y - b.y;
  double cosine = (ux * vx + uy * vy) / (std::hypot(ux, uy) * std::hypot(vx, vy));
  return std::acos(std::clamp(cosine, -1.0, 1.0));
}

template<typename Pt>
double angle_between(const Pt &a, const Pt &b) {
  double t = std::atan2(a.x * b.y - a.y * b.x, a.x * b.x + a.y * b.y);
  return (t < 0) ? (t + 2 * PI) : t;
}

double angle_between(double a1, double b1, double a2, double b2) {
  double t = std::atan2(a1 * b2 - a2 * b1, a1 * a2 + b1 * b2);
  if (t < 0) {
    t += PI;
  }
  return LT(PI / 2, t) ? (PI - t) : t;
}

template<typename Pt>
auto cross(const Pt &a, const Pt &b, const Pt &o = Pt(0, 0)) {
  return (a.x - o.x) * (b.y - o.y) - (a.y - o.y) * (b.x - o.x);
}

template<typename Pt>
int turn(const Pt &a, const Pt &o, const Pt &b) {
  auto c = cross(a, b, o);
  return LT(c, 0) ? 1 : (LT(0, c) ? -1 : 0);
}
\end{lstlisting}
\Needspace*{4\baselineskip}
\textbf{Example Usage}\par
\begin{lstlisting}[style=cpp]
struct Point {
  double x, y;
  Point(double x = 0, double y = 0) : x(x), y(y) {}
  bool operator==(const Point &p) const { return x == p.x && y == p.y; }
  friend bool EQ(const Point &a, const Point &b) { return EQ(a.x, b.x) && EQ(a.y, b.y); }
};

#include <cassert>

int main() {
  assert(EQ(123, reduce_deg(-8 * 360 + 123)));
  assert(EQ(1.2345, reduce_rad(2 * PI * 8 + 1.2345)));
  assert(EQ(-4, polar_x(4, PI)) && EQ(0, polar_y(4, PI)));
  assert(EQ(0, polar_x(4, -PI / 2)) && EQ(-4, polar_y(4, -PI / 2)));
  assert(EQ(45, polar_angle(Point(5, 5)) * RAD));
  assert(EQ(135 * DEG, polar_angle(Point(-4, 4))));
  assert(EQ(90 * DEG, angle(Point(5, 0), Point(0, 5), Point(-5, 0))));
  assert(EQ(225 * DEG, angle_between(Point(0, 5), Point(5, -5))));
  assert(-1 == cross(Point(0, 1), Point(1, 0), Point(0, 0)));
  assert(-1 == turn(Point(0, 1), Point(0, 0), Point(-5, -5)));
  return 0;
}
\end{lstlisting}
\subsection{Distances}
\label{sec:7.2.2}
Distance calculations in two dimensions for points, lines, and line segments. The functions are templated on the point type \inlinecode{Pt}: pass any struct with numeric \inlinecode{.x} and \inlinecode{.y} fields (for example \inlinecode{PointD}/\inlinecode{PointI} from \secref{7.1.1}, or the local example struct). \inlinecode{sqdist} preserves the coordinate type and is exact for integer points. Functions that return an actual distance use \inlinecode{double} for the final division/square root. \inlinecode{closest\_point} returns a point of the same type as its inputs, so use a floating-point \inlinecode{Pt} there for a meaningful (non-truncated) result.

\begin{compactitem}
  \item \inlinecode{dist(a, b)} and \inlinecode{sqdist(a, b)} respectively return the distance and squared distance between points \inlinecode{a} and \inlinecode{b}.
  \item \inlinecode{line\_dist(p, a, b, c)} returns the distance from point \inlinecode{p} to the valid line $ax + by + c = 0$.
  \item \inlinecode{line\_dist(p, a, b)} returns the distance from point \inlinecode{p} to the infinite line containing points \inlinecode{a} and \inlinecode{b}. If \inlinecode{a} = \inlinecode{b}, the distance from \inlinecode{p} to \inlinecode{a} is returned.
  \item \inlinecode{line\_dist(a1, b1, c1, a2, b2, c2)} returns the distance between two valid lines.
  \item \inlinecode{seg\_dist(p, a, b)} returns the distance from point \inlinecode{p} to the line segment \inlinecode{a}-\inlinecode{b}.
  \item \inlinecode{seg\_dist(a, b, c, d)} returns the minimum distance between line segments \inlinecode{a}-\inlinecode{b} and \inlinecode{c}-\inlinecode{d}, or $0$ if the segments touch or intersect.
  \item \inlinecode{closest\_point(a, b, p)} returns the point on segment \inlinecode{a}-\inlinecode{b} closest to point \inlinecode{p}, using the input point type for the result. Use a floating-point point type when the closest point may have fractional coordinates.
\end{compactitem}

Overflow warning: squared distances, dot products, and cross products are formed in the point's coordinate arithmetic type. For integer point types, use a 64-bit coordinate type when coordinates may exceed a few tens of thousands.

\textbf{Time Complexity}: $O(1)$ per operation.

\textbf{Space Complexity}: $O(1)$ auxiliary for all operations.

\begin{lstlisting}[style=cpp]
#include <algorithm>
#include <cmath>
#include <type_traits>

const double EPS = 1e-9;

template<typename T, typename U, typename C = std::common_type_t<T, U>>
bool EQ(T a, U b) {
  if constexpr (std::is_floating_point_v<C>) return C(a) == C(b) || std::fabs(C(a) - C(b)) <= EPS;
  return C(a) == C(b);
}

template<typename T, typename U, typename C = std::common_type_t<T, U>>
bool LT(T a, U b) {
  if constexpr (std::is_floating_point_v<C>) return C(a) < C(b) - EPS;
  return C(a) < C(b);
}

template<typename T, typename U> bool LE(T a, U b) { return !LT(b, a); }
template<typename T, typename U> bool GE(T a, U b) { return !LT(a, b); }

// sqdist returns the coordinate type (exact for integer points).
template<typename Pt>
auto sqdist(const Pt &a, const Pt &b) {
  auto dx = b.x - a.x, dy = b.y - a.y;
  return dx * dx + dy * dy;  // Overflow warning.
}

template<typename Pt>
double dist(const Pt &a, const Pt &b) {
  return std::sqrt(static_cast<double>(sqdist(a, b)));
}

template<typename Pt, typename T>
double line_dist(const Pt &p, const T &a, const T &b, const T &c) {
  return fabs(static_cast<double>(a) * p.x + static_cast<double>(b) * p.y + c) /
         std::sqrt(static_cast<double>(a) * a + static_cast<double>(b) * b);
}

template<typename Pt>
double line_dist(const Pt &p, const Pt &a, const Pt &b) {
  if (EQ(a.x, b.x) && EQ(a.y, b.y)) {
    return dist(p, a);
  }
  auto n = sqdist(a, b);
  auto d = (p.x - a.x) * (b.x - a.x) + (p.y - a.y) * (b.y - a.y);  // Overflow warning.
  double u = static_cast<double>(d) / n;
  double dx = a.x + u * (b.x - a.x) - p.x, dy = a.y + u * (b.y - a.y) - p.y;
  return std::sqrt(dx * dx + dy * dy);
}

template<typename T>
double line_dist(const T &a1, const T &b1, const T &c1, const T &a2, const T &b2, const T &c2) {
  if (EQ(static_cast<double>(a1) * b2, static_cast<double>(a2) * b1)) {
    double factor = EQ(b1, 0) ? (static_cast<double>(a1) / a2) : (static_cast<double>(b1) / b2);
    return EQ(c1, c2 * factor)
               ? 0
               : fabs(c2 * factor - c1) /
                     std::sqrt(static_cast<double>(a1) * a1 + static_cast<double>(b1) * b1);
  }
  return 0;
}

template<typename Pt>
double seg_dist(const Pt &p, const Pt &a, const Pt &b) {
  if (EQ(a.x, b.x) && EQ(a.y, b.y)) {
    return dist(p, a);
  }
  auto n = sqdist(a, b);
  auto d = (p.x - a.x) * (b.x - a.x) + (p.y - a.y) * (b.y - a.y);  // Overflow warning.
  if (LE(d, 0) || EQ(n, 0)) {
    return dist(p, a);
  }
  if (GE(d, n)) {
    return dist(p, b);
  }
  double t = static_cast<double>(d) / n;
  double dx = a.x + t * (b.x - a.x) - p.x, dy = a.y + t * (b.y - a.y) - p.y;
  return std::sqrt(dx * dx + dy * dy);
}

template<typename Pt>
double seg_dist(const Pt &a, const Pt &b, const Pt &c, const Pt &d) {
  if (EQ(a.x, b.x) && EQ(a.y, b.y)) {
    return seg_dist(a, c, d);
  }
  if (EQ(c.x, d.x) && EQ(c.y, d.y)) {
    return seg_dist(c, a, b);
  }
  auto ab_x = b.x - a.x, ab_y = b.y - a.y;
  auto ac_x = c.x - a.x, ac_y = c.y - a.y;
  auto cd_x = d.x - c.x, cd_y = d.y - c.y;
  auto c1 = ab_x * cd_y - ab_y * cd_x;  // Overflow warning.
  auto c2 = ac_x * ab_y - ac_y * ab_x;
  if (EQ(c1, 0) && EQ(c2, 0)) {
    auto less = [](const Pt &u, const Pt &v) { return u.x != v.x ? u.x < v.x : u.y < v.y; };
    auto minp = [&](const Pt &u, const Pt &v) -> const Pt & { return less(v, u) ? v : u; };
    auto maxp = [&](const Pt &u, const Pt &v) -> const Pt & { return less(u, v) ? v : u; };
    const Pt &res1 = maxp(minp(a, b), minp(c, d));
    const Pt &res2 = minp(maxp(a, b), maxp(c, d));
    if (!less(res2, res1)) {
      return 0;
    }
  } else if (!EQ(c1, 0)) {
    auto t_num = ac_x * cd_y - ac_y * cd_x;
    bool c1_pos = c1 > 0;
    bool t_between_01 = c1_pos ? (LE(0, t_num) && LE(t_num, c1)) : (LE(t_num, 0) && LE(c1, t_num));
    bool u_between_01 = c1_pos ? (LE(0, c2) && LE(c2, c1)) : (LE(c2, 0) && LE(c1, c2));
    if (t_between_01 && u_between_01) {
      return 0;
    }
  }
  return std::min(
      std::min(seg_dist(a, c, d), seg_dist(b, c, d)), std::min(seg_dist(c, a, b), seg_dist(d, a, b))
  );
}

template<typename Pt>
Pt closest_point(const Pt &a, const Pt &b, const Pt &p) {
  Pt res{};
  if (EQ(a.x, b.x) && EQ(a.y, b.y)) {
    res.x = a.x;
    res.y = a.y;
    return res;
  }
  auto n = sqdist(a, b);
  auto d = (p.x - a.x) * (b.x - a.x) + (p.y - a.y) * (b.y - a.y);  // Overflow warning.
  double t = static_cast<double>(d) / n;
  if (t <= 0) {
    res.x = a.x;
    res.y = a.y;
  } else if (t >= 1) {
    res.x = b.x;
    res.y = b.y;
  } else {
    res.x = a.x + t * (b.x - a.x);
    res.y = a.y + t * (b.y - a.y);
  }
  return res;
}
\end{lstlisting}
\Needspace*{4\baselineskip}
\textbf{Example Usage}\par
\begin{lstlisting}[style=cpp]
#include <cassert>

struct Point {
  double x, y;
  Point(double x = 0, double y = 0) : x(x), y(y) {}
};

struct PointI {
  int x, y;
  PointI(int x = 0, int y = 0) : x(x), y(y) {}
};

int main() {
  assert(EQ(5, dist(Point(-1, -1), Point(2, 3))));
  assert(EQ(25, sqdist(Point(-1, -1), Point(2, 3))));
  assert(EQ(1.2, line_dist(Point(2, 1), -4, 3, -1)));
  assert(EQ(0.8, line_dist(Point(3, 3), Point(-1, -1), Point(2, 3))));
  assert(EQ(1.2, line_dist(Point(2, 1), Point(-1, -1), Point(2, 3))));
  assert(EQ(0, line_dist(-4, 3, -1, 8, 6, 2)));
  assert(EQ(0.8, line_dist(-4, 3, -1, -8, 6, -10)));
  assert(EQ(1.0, seg_dist(Point(3, 3), Point(-1, -1), Point(2, 3))));
  assert(EQ(1.2, seg_dist(Point(2, 1), Point(-1, -1), Point(2, 3))));
  assert(EQ(0, seg_dist(Point(0, 2), Point(3, 3), Point(-1, -1), Point(2, 3))));
  assert(EQ(0.6, seg_dist(Point(-1, 0), Point(-2, 2), Point(-1, -1), Point(2, 3))));

  // Integer points: sqdist is exact, dist returns double.
  PointI ia{-1, -1}, ib{2, 3};
  assert(sqdist(ia, ib) == 25);   // int result
  assert(EQ(dist(ia, ib), 5.0));  // double result
  return 0;
}
\end{lstlisting}
\subsection{Line Intersection}
\label{sec:7.2.3}
Intersection calculations in two dimensions for straight lines and line segments. The functions are templated on the point type \inlinecode{Pt}: pass any struct with numeric \inlinecode{.x} and \inlinecode{.y} fields (for example \inlinecode{PointD}/\inlinecode{PointI} from \secref{7.1.1}, or the local example struct) for which \inlinecode{operator\textless{}} orders points lexicographically using exact coordinate comparisons. Point output types are deduced from the caller-supplied arguments, so they may differ from the input type, e.g. integer endpoints with a floating-point output point.

\begin{compactitem}
  \item \inlinecode{line\_intersection(a1, b1, c1, a2, b2, c2, \&p)} intersects lines $a_1 x + b_1 y + c_1 = 0$ and $a_2 x + b_2 y + c_2 = 0$, returning $-1$ (parallel), $0$ (one point, stored into \inlinecode{p} if it is not \inlinecode{nullptr}), or $1$ (identical). Both coefficient triples must represent valid lines.
  \item \inlinecode{line\_intersection(p1, p2, p3, p4, \&p)} intersects the infinite lines determined by points \inlinecode{p1}, \inlinecode{p2}, \inlinecode{p3}, and \inlinecode{p4}, returning $-1$ (parallel), $0$ (one point, stored into \inlinecode{p} if it is not \inlinecode{nullptr}), or $1$ (identical). The points in each pair must differ.
  \item \inlinecode{seg\_intersection(a, b, c, d, \&p, \&q, include\_boundary = true)} intersects segments \inlinecode{a}-\inlinecode{b} and \inlinecode{c}-\inlinecode{d}, returning $-1$ (none), $0$ (one point), or $1$ (overlapping segment). The \inlinecode{include\_boundary} flag controls whether segments that meet only at an endpoint count as intersecting. If the intersection is a point (and \inlinecode{p} is not \inlinecode{nullptr}), it will be stored into \inlinecode{p}. If the intersection is an overlapping segment (and \inlinecode{p} and \inlinecode{q} are not \inlinecode{nullptr}) then their endpoints will be stored into pointers \inlinecode{p} and \inlinecode{q}, which must reference a floating-point point type.
  \item \inlinecode{seg\_intersection(a, b, c, d, include\_boundary = true)} is a simplified, detection-only version of the above. Both versions are exact for detection if the input point type is integral.
  \item \inlinecode{closest\_point(a, b, c, p)} returns the closest point on the valid line $ax + by + c = 0$ to point \inlinecode{p}.
\end{compactitem}

Overflow warning: the exact integer paths multiply coordinate differences (cross products and squared lengths), which grow like the squared coordinate magnitude. With 32-bit \inlinecode{int} coordinates these overflow once coordinates exceed roughly a few tens of thousands, so for larger integer inputs use a point type with 64-bit (\inlinecode{int64\_t}) coordinates.

\textbf{Time Complexity}: $O(1)$ per operation.

\textbf{Space Complexity}: $O(1)$ auxiliary for all operations.

\begin{lstlisting}[style=cpp]
#include <algorithm>
#include <cmath>
#include <type_traits>

const double EPS = 1e-9;

template<typename T, typename U, typename C = std::common_type_t<T, U>>
bool EQ(T a, U b) {
  if constexpr (std::is_floating_point_v<C>) return C(a) == C(b) || std::fabs(C(a) - C(b)) <= EPS;
  return C(a) == C(b);
}

template<typename T, typename U, typename C = std::common_type_t<T, U>>
bool LT(T a, U b) {
  if constexpr (std::is_floating_point_v<C>) return C(a) < C(b) - EPS;
  return C(a) < C(b);
}

template<typename T, typename U>
bool LE(T a, U b) {
  return !LT(b, a);
}

// Guard to check outputs are real coordinates, since intersections are generally not integral.
template<typename OutPt>
constexpr bool is_float_pt =
    std::is_floating_point_v<decltype(OutPt::x)> && std::is_floating_point_v<decltype(OutPt::y)>;

template<typename OutPt>
int line_intersection(double a1, double b1, double c1, double a2, double b2, double c2, OutPt *p) {
  static_assert(is_float_pt<OutPt>, "output point coordinates must be floating-point");
  double det = a1 * b2 - a2 * b1;  // Cramer's rule for a1*x + b1*y = -c1, a2*x + b2*y = -c2.
  if (EQ(det, 0)) {                // Parallel (lines need not be given in normalized form).
    // Identical iff the other two 2x2 determinants also vanish.
    return (EQ(a1 * c2 - a2 * c1, 0) && EQ(b1 * c2 - b2 * c1, 0)) ? 1 : -1;
  }
  if (p != nullptr) {
    double x = (b1 * c2 - b2 * c1) / det;  // Local x to ensure y is derived from full-precision x.
    p->x = x;
    p->y = !EQ(b1, 0) ? -(a1 * x + c1) / b1 : -(a2 * x + c2) / b2;
  }
  return 0;
}

template<typename Pt, typename OutPt>
int line_intersection(const Pt &p1, const Pt &p2, const Pt &p3, const Pt &p4, OutPt *p) {
  static_assert(is_float_pt<OutPt>, "output point coordinates must be floating-point");
  auto a1 = p2.y - p1.y, b1 = p1.x - p2.x;
  auto c1 = -(p1.x * p2.y - p2.x * p1.y);
  auto a2 = p4.y - p3.y, b2 = p3.x - p4.x;
  auto c2 = -(p3.x * p4.y - p4.x * p3.y);
  auto x = -(c1 * b2 - c2 * b1), y = -(a1 * c2 - a2 * c1);
  auto det = a1 * b2 - a2 * b1;
  if (EQ(det, 0)) {
    return (EQ(x, 0) && EQ(y, 0)) ? 1 : -1;
  }
  if (p != nullptr) {
    p->x = static_cast<double>(x) / det;
    p->y = static_cast<double>(y) / det;
  }
  return 0;
}

template<typename Pt>
bool point_on_segment(const Pt &p, const Pt &a, const Pt &b) {
  // Overflow risk for integer Pt: these products are ~O(max_coord^2); use int64_t if necessary.
  return EQ((p.x - a.x) * (b.y - a.y) - (p.y - a.y) * (b.x - a.x), 0) &&
         LE(std::min(a.x, b.x), p.x) && LE(p.x, std::max(a.x, b.x)) &&
         LE(std::min(a.y, b.y), p.y) && LE(p.y, std::max(a.y, b.y));
}

// Detection is exact for integer Pt. Intersection point output may use a separate float OutPt.
template<typename Pt, typename OutPt>
int seg_intersection(
    const Pt &a, const Pt &b, const Pt &c, const Pt &d, OutPt *p = nullptr, OutPt *q = nullptr,
    const bool include_boundary = true
) {
  static_assert(is_float_pt<OutPt>, "output point coordinates must be floating-point");
  auto set_out = [](auto *out, auto x, auto y) {
    if (out != nullptr) {
      out->x = x;
      out->y = y;
    }
  };
  auto ab_x = b.x - a.x, ab_y = b.y - a.y;
  auto ac_x = c.x - a.x, ac_y = c.y - a.y;
  auto cd_x = d.x - c.x, cd_y = d.y - c.y;
  auto ab2 = ab_x * ab_x + ab_y * ab_y;
  auto cd2 = cd_x * cd_x + cd_y * cd_y;
  if (EQ(ab2, 0)) {
    if (include_boundary && point_on_segment(a, c, d)) {
      set_out(p, a.x, a.y);
      return 0;  // Degenerate: first segment is a point intersecting the second segment.
    }
    return -1;  // Degenerate: first segment is a point not intersecting the second segment.
  }
  if (EQ(cd2, 0)) {
    if (include_boundary && point_on_segment(c, a, b)) {
      set_out(p, c.x, c.y);
      return 0;  // Degenerate: second segment is a point intersecting the first segment.
    }
    return -1;  // Degenerate: second segment is a point not intersecting the first segment.
  }
  auto c1 = ab_x * cd_y - ab_y * cd_x;
  auto c2 = ac_x * ab_y - ac_y * ab_x;
  if (EQ(c1, 0) && EQ(c2, 0)) {  // Collinear.
    Pt res1 = std::max(std::min(a, b), std::min(c, d));
    Pt res2 = std::min(std::max(a, b), std::max(c, d));
    if (include_boundary ? !(res2 < res1) : (res1 < res2)) {
      if (res1 == res2) {
        set_out(p, res1.x, res1.y);
        return 0;  // Collinear and overlapping: touch at one endpoint.
      }
      if (p != nullptr && q != nullptr) {  // Both ends are reported, or neither is.
        set_out(p, res1.x, res1.y);
        set_out(q, res2.x, res2.y);
      }
      return 1;  // Collinear and overlapping: overlap is a segment.
    }
    return -1;  // Collinear and disjoint.
  }
  if (EQ(c1, 0)) {
    return -1;  // Parallel and disjoint.
  }
  auto t_num = ac_x * cd_y - ac_y * cd_x;
  bool c1_pos = c1 > 0;
  bool t_ok =
      c1_pos
          ? (include_boundary ? (LE(0, t_num) && LE(t_num, c1)) : (LT(0, t_num) && LT(t_num, c1)))
          : (include_boundary ? (LE(c1, t_num) && LE(t_num, 0)) : (LT(c1, t_num) && LT(t_num, 0)));
  bool u_ok = c1_pos ? (include_boundary ? (LE(0, c2) && LE(c2, c1)) : (LT(0, c2) && LT(c2, c1)))
                     : (include_boundary ? (LE(c1, c2) && LE(c2, 0)) : (LT(c1, c2) && LT(c2, 0)));
  if (t_ok && u_ok) {
    double t = static_cast<double>(t_num) / c1;
    set_out(p, a.x + t * ab_x, a.y + t * ab_y);
    return 0;  // Non-parallel with one intersection.
  }
  return -1;  // Non-parallel with no intersection.
}

// Simplified detection-only version (no output points needed); exact for integer Pt.
// Alternatively, just call the version above and pass static_cast<Pt *>(nullptr) for p and q.
template<typename Pt>
int seg_intersection(
    const Pt &a, const Pt &b, const Pt &c, const Pt &d, bool include_boundary = true
) {
  auto ab_x = b.x - a.x, ab_y = b.y - a.y;
  auto ac_x = c.x - a.x, ac_y = c.y - a.y;
  auto cd_x = d.x - c.x, cd_y = d.y - c.y;
  auto ab2 = ab_x * ab_x + ab_y * ab_y;
  auto cd2 = cd_x * cd_x + cd_y * cd_y;
  if (EQ(ab2, 0)) {
    return (include_boundary && point_on_segment(a, c, d)) ? 0 : -1;
  }
  if (EQ(cd2, 0)) {
    return (include_boundary && point_on_segment(c, a, b)) ? 0 : -1;
  }
  auto c1 = ab_x * cd_y - ab_y * cd_x;
  auto c2 = ac_x * ab_y - ac_y * ab_x;
  if (EQ(c1, 0) && EQ(c2, 0)) {  // Collinear.
    Pt res1 = std::max(std::min(a, b), std::min(c, d));
    Pt res2 = std::min(std::max(a, b), std::max(c, d));
    if (include_boundary ? !(res2 < res1) : (res1 < res2)) {
      return (res1 == res2) ? 0 : 1;
    }
    return -1;
  }
  if (EQ(c1, 0)) {
    return -1;
  }
  auto t_num = ac_x * cd_y - ac_y * cd_x;
  bool c1_pos = c1 > 0;
  bool t_ok =
      c1_pos
          ? (include_boundary ? (LE(0, t_num) && LE(t_num, c1)) : (LT(0, t_num) && LT(t_num, c1)))
          : (include_boundary ? (LE(c1, t_num) && LE(t_num, 0)) : (LT(c1, t_num) && LT(t_num, 0)));
  bool u_ok = c1_pos ? (include_boundary ? (LE(0, c2) && LE(c2, c1)) : (LT(0, c2) && LT(c2, c1)))
                     : (include_boundary ? (LE(c1, c2) && LE(c2, 0)) : (LT(c1, c2) && LT(c2, 0)));
  return (t_ok && u_ok) ? 0 : -1;
}

template<typename Pt>
Pt closest_point(double a, double b, double c, const Pt &p) {
  Pt res{};
  if (EQ(a, 0)) {
    res.x = p.x;
    res.y = -c / b;
  } else if (EQ(b, 0)) {
    res.x = -c / a;
    res.y = p.y;
  } else {
    line_intersection(a, b, c, -b, a, b * p.x - a * p.y, &res);
  }
  return res;
}
\end{lstlisting}
\Needspace*{4\baselineskip}
\textbf{Example Usage}\par
\begin{lstlisting}[style=cpp]
#include <cassert>

struct Point {
  double x, y;
  Point(double x = 0, double y = 0) : x(x), y(y) {}
  bool operator==(const Point &p) const { return x == p.x && y == p.y; }
  friend bool EQ(const Point &a, const Point &b) { return EQ(a.x, b.x) && EQ(a.y, b.y); }
  bool operator!=(const Point &p) const { return !(*this == p); }
  bool operator<(const Point &p) const { return x != p.x ? x < p.x : y < p.y; }
  bool operator>(const Point &p) const { return p < *this; }
};

struct PointI {
  int x, y;
  PointI(int x = 0, int y = 0) : x(x), y(y) {}
  bool operator==(const PointI &p) const { return x == p.x && y == p.y; }
  bool operator!=(const PointI &p) const { return !(*this == p); }
  bool operator<(const PointI &p) const { return x != p.x ? x < p.x : y < p.y; }
  bool operator>(const PointI &p) const { return p < *this; }
};

int main() {
  Point p, q;
  assert(line_intersection(-1.0, 1.0, 0.0, 1.0, 1.0, -3.0, &p) == 0);
  assert(EQ(p, Point(1.5, 1.5)));
  assert(line_intersection(Point(0, 0), Point(1, 1), Point(0, 4), Point(4, 0), &p) == 0);
  assert(EQ(p, Point(2, 2)));

  auto test = [&](double a, double b, double c, double d, double e, double f, double g, double h) {
    return seg_intersection(Point(a, b), Point(c, d), Point(e, f), Point(g, h), &p, &q);
  };
  assert(0 == test(-4, 0, 4, 0, 0, -4, 0, 4) && EQ(p, Point(0, 0)));
  assert(0 == test(0, 0, 10, 10, 2, 2, 16, 4) && EQ(p, Point(2, 2)));
  assert(0 == test(-2, 2, -2, -2, -2, 0, 0, 0) && EQ(p, Point(-2, 0)));
  assert(0 == test(0, 4, 4, 4, 4, 0, 4, 8) && EQ(p, Point(4, 4)));
  assert(1 == test(10, 10, 0, 0, 2, 2, 6, 6));
  assert(EQ(p, Point(2, 2)) && EQ(q, Point(6, 6)));
  assert(-1 == test(6, 8, 8, 10, 12, 12, 4, 4));
  assert(-1 == test(-4, 2, -8, 8, 0, 0, -4, 6));

  assert(EQ(Point(2.5, 2.5), closest_point(-1.0, -1.0, 5.0, Point(0, 0))));
  assert(EQ(Point(3, 0), closest_point(1.0, 0.0, -3.0, Point(0, 0))));

  // Detection-only with integer coordinates (exact, no float needed).
  assert(seg_intersection(PointI(0, 0), PointI(4, 4), PointI(0, 4), PointI(4, 0)) == 0);
  assert(seg_intersection(PointI(0, 0), PointI(1, 0), PointI(2, 0), PointI(3, 0)) == -1);
  assert(seg_intersection(PointI(0, 0), PointI(4, 4), PointI(0, 4), PointI(4, 0), &p) == 0);
  assert(EQ(p, Point(2, 2)));

  // include_boundary = false treats endpoint-only contact as non-intersecting. The T-junction
  // and shared-endpoint cases below intersect by default but are rejected when the flag is off.
  assert(0 == seg_intersection(Point(0, 4), Point(4, 4), Point(4, 0), Point(4, 8), &p, &q));
  assert(-1 == seg_intersection(Point(0, 4), Point(4, 4), Point(4, 0), Point(4, 8), &p, &q, false));
  assert(-1 == seg_intersection(PointI(0, 0), PointI(2, 2), PointI(2, 2), PointI(6, 6), false));
  return 0;
}
\end{lstlisting}
\subsection{Circle Intersection}
\label{sec:7.2.4}
Circle tangent and intersection calculations in two dimensions. Such calculations are floating-point by nature, so this section mostly computes in \inlinecode{double}. The input point \inlinecode{p} to \inlinecode{tangent()} can be any type with numeric \inlinecode{.x} and \inlinecode{.y} fields, and intersection point outputs are written through a caller-supplied point type.

Circle radii are normalized to be nonnegative.

\begin{compactitem}
  \item \inlinecode{tangent(c, p, \&l1, \&l2)} determines the line(s) tangent to circle \inlinecode{c} that pass through point \inlinecode{p}, returning $-1$ if there is no tangent line because \inlinecode{p} is strictly inside \inlinecode{c}, $0$ if there is exactly one tangent line because \inlinecode{p} is on the boundary of \inlinecode{c} (in which case the line will be stored into pointer \inlinecode{l1} if it's not \inlinecode{nullptr}), or $1$ if there are two tangent lines because \inlinecode{p} is strictly outside of \inlinecode{c} (in which case the lines will be stored into pointers \inlinecode{l1} and \inlinecode{l2} if they are not \inlinecode{nullptr}). The circle must have positive radius.
  \item \inlinecode{intersection(c, l, \&p, \&q)} determines the intersection between the circle \inlinecode{c} and line \inlinecode{l}, returning $-1$ if there is no intersection, $0$ if the line has one intersection point because the line is tangent (in which case it will be stored into pointer \inlinecode{p} if it's not \inlinecode{nullptr}), or $1$ if there are two intersection points because the line crosses through the circle (stored independently into \inlinecode{p} and \inlinecode{q} when those pointers are non-null). The line must be valid.
  \item \inlinecode{intersection(c1, c2, \&p, \&q)} determines the intersection points between two circles \inlinecode{c1} and \inlinecode{c2}, returning $-2$ if circle \inlinecode{c2} completely encloses circle \inlinecode{c1}, $-1$ if circle \inlinecode{c1} completely encloses circle \inlinecode{c2}, $0$ if the circles are completely disjoint, $1$ if the circles are tangent with one intersection (stored in \inlinecode{p} if it's not \inlinecode{nullptr}), $2$ if the circles intersect at two points (stored independently in \inlinecode{p} and \inlinecode{q} when those pointers are non-null), or $3$ if the circles are equal and intersect at infinite points.
  \item \inlinecode{intersection\_area(c1, c2)} returns the intersection area of circles \inlinecode{c1} and \inlinecode{c2}.
\end{compactitem}

\textbf{Time Complexity}: $O(1)$ per operation.

\textbf{Space Complexity}: $O(1)$ auxiliary for all operations.

\begin{lstlisting}[style=cpp]
#include <algorithm>
#include <cmath>
#include <type_traits>

const double EPS = 1e-9, PI = std::acos(-1.0);

template<typename T, typename U, typename C = std::common_type_t<T, U>>
bool EQ(T a, U b) {
  if constexpr (std::is_floating_point_v<C>) return C(a) == C(b) || std::fabs(C(a) - C(b)) <= EPS;
  return C(a) == C(b);
}

template<typename T, typename U, typename C = std::common_type_t<T, U>>
bool LT(T a, U b) {
  if constexpr (std::is_floating_point_v<C>) return C(a) < C(b) - EPS;
  return C(a) < C(b);
}

template<typename T, typename U>
bool LE(T a, U b) {
  return !LT(b, a);
}

struct Circle {
  double h, k, r;

  Circle(double h, double k, double r) : h(h), k(k), r(fabs(r)) {}
};

struct Line {
  double a, b, c;

  Line() : a(0), b(0), c(0) {}

  Line(double a, double b, double c) {
    if (!EQ(b, 0)) {
      this->a = a / b;
      this->c = c / b;
      this->b = 1;
    } else {
      this->c = c / a;
      this->a = 1;
      this->b = 0;
    }
  }

  Line(double px, double py, double qx, double qy) : a(0), b(0), c(0) {
    if (EQ(px, qx)) {
      if (!EQ(py, qy)) {  // Vertical line.
        a = 1;
        b = 0;
        c = -px;
      }  // Else, invalid line.
    } else {
      a = -static_cast<double>(py - qy) / (px - qx);
      b = 1;
      c = -(a * px) - (b * py);
    }
  }

  bool operator==(const Line &l) const { return a == l.a && b == l.b && c == l.c; }

  friend bool EQ(const Line &l1, const Line &l2) {
    return EQ(l1.a, l2.a) && EQ(l1.b, l2.b) && EQ(l1.c, l2.c);
  }
};

template<typename Pt>
int tangent(const Circle &c, const Pt &p, Line *l1 = nullptr, Line *l2 = nullptr) {
  auto sqnorm = [](double x, double y) { return x * x + y * y; };
  double vopx = p.x - c.h, vopy = p.y - c.k;
  if (EQ(sqnorm(vopx, vopy), c.r * c.r)) {  // Point on an edge.
    if (l1 != nullptr) {                    // Tangent through p is perpendicular to the radius cp.
      *l1 = Line(vopx, vopy, -(vopx * p.x + vopy * p.y));
    }
    return 0;
  }
  if (LE(sqnorm(vopx, vopy), c.r * c.r)) {
    return -1;  // Point inside Circle, no intersection.
  }
  double qx = vopx / c.r, qy = vopy / c.r;
  double n = sqnorm(qx, qy), d = qy * std::sqrt(n - 1.0);
  double t1x = (qx - d) / n, t1y = c.k;
  double t2x = (qx + d) / n, t2y = c.k;
  if (!EQ(qy, 0)) {  // Common case.
    t1y += c.r * (1.0 - t1x * qx) / qy;
    t2y += c.r * (1.0 - t2x * qx) / qy;
  } else {  // Point at center horizontal, y = 0.
    d = c.r * std::sqrt(1.0 - t1x * t1x);
    t1y += d;
    t2y -= d;
  }
  t1x = t1x * c.r + c.h;
  t2x = t2x * c.r + c.h;
  // Note: here, t1 and t2 are the two points of tangencies
  if (l1 != nullptr) {
    *l1 = Line(p.x, p.y, t1x, t1y);
  }
  if (l2 != nullptr) {
    *l2 = Line(p.x, p.y, t2x, t2y);
  }
  return 1;
}

// Guard to check outputs are real coordinates, since intersections are generally not integral.
template<typename OutPt>
constexpr bool is_float_pt =
    std::is_floating_point_v<decltype(OutPt::x)> && std::is_floating_point_v<decltype(OutPt::y)>;

template<typename OutPt>
int intersection(const Circle &c, const Line &l, OutPt *p = nullptr, OutPt *q = nullptr) {
  static_assert(is_float_pt<OutPt>, "output point coordinates must be floating-point");
  auto set_out = [](auto *out, auto x, auto y) {
    if (out != nullptr) {
      out->x = x;
      out->y = y;
    }
  };
  double v = c.h * l.a + c.k * l.b + l.c;
  double aabb = l.a * l.a + l.b * l.b;
  double disc = v * v / aabb - c.r * c.r;
  if (disc > EPS) {
    return -1;
  }
  double x0 = -l.a * v / aabb, y0 = -l.b * v / aabb;
  if (disc > -EPS) {
    set_out(p, x0 + c.h, y0 + c.k);
    return 0;
  }
  double k = std::sqrt(std::max(0.0, disc / -aabb));
  set_out(p, x0 + k * l.b + c.h, y0 - k * l.a + c.k);
  set_out(q, x0 - k * l.b + c.h, y0 + k * l.a + c.k);
  return 1;
}

template<typename OutPt>
int intersection(const Circle &c1, const Circle &c2, OutPt *p = nullptr, OutPt *q = nullptr) {
  static_assert(is_float_pt<OutPt>, "output point coordinates must be floating-point");
  auto set_out = [](auto *out, auto x, auto y) {
    if (out != nullptr) {
      out->x = x;
      out->y = y;
    }
  };
  if (EQ(c1.h, c2.h) && EQ(c1.k, c2.k)) {
    return EQ(c1.r, c2.r) ? 3 : (c1.r > c2.r ? -1 : -2);
  }
  double d12x = c2.h - c1.h, d12y = c2.k - c1.k;
  double d = std::hypot(d12x, d12y);
  if (LT(c1.r + c2.r, d)) {
    return 0;
  }
  if (LT(d, fabs(c1.r - c2.r))) {
    return c1.r > c2.r ? -1 : -2;
  }
  double a = (c1.r * c1.r - c2.r * c2.r + d * d) / (2 * d);
  double x0 = c1.h + (d12x * a / d), y0 = c1.k + (d12y * a / d);
  double s = std::sqrt(std::max(0.0, c1.r * c1.r - a * a));
  double rx = -d12y * s / d, ry = d12x * s / d;
  if (EQ(rx, 0) && EQ(ry, 0)) {
    set_out(p, x0, y0);
    return 1;
  }
  set_out(p, x0 - rx, y0 - ry);
  set_out(q, x0 + rx, y0 + ry);
  return 2;
}

double intersection_area(const Circle &c1, const Circle &c2) {
  double r = std::min(c1.r, c2.r), R = std::max(c1.r, c2.r);
  double d = std::hypot(c2.h - c1.h, c2.k - c1.k);
  if (LE(d, R - r)) {
    return PI * r * r;
  }
  if (LE(R + r, d)) {
    return 0;
  }
  double alpha = std::clamp((d * d + r * r - R * R) / 2 / d / r, -1.0, 1.0);
  double beta = std::clamp((d * d + R * R - r * r) / 2 / d / R, -1.0, 1.0);
  return r * r * std::acos(alpha) + R * R * std::acos(beta) -
         0.5 * std::sqrt((-d + r + R) * (d + r - R) * (d - r + R) * (d + r + R));
}
\end{lstlisting}
\Needspace*{4\baselineskip}
\textbf{Example Usage}\par
\begin{lstlisting}[style=cpp]
#include <cassert>

struct Point {
  double x, y;
  Point(double x = 0, double y = 0) : x(x), y(y) {}
  bool operator==(const Point &p) const { return x == p.x && y == p.y; }
  friend bool EQ(const Point &a, const Point &b) { return EQ(a.x, b.x) && EQ(a.y, b.y); }
};

int main() {
  Line l1, l2;
  assert(-1 == tangent(Circle(0, 0, 4), Point(1, 1), &l1, &l2));
  assert(0 == tangent(Circle(0, 0, sqrt(2)), Point(1, 1), &l1, &l2));
  assert(EQ(l1, Line(-1, -1, 2)));
  assert(1 == tangent(Circle(0, 0, 2), Point(2, 2), &l1, &l2));
  assert(EQ(l1, Line(0, -2, 4)));
  assert(EQ(l2, Line(2, 0, -4)));

  Point p, q;
  assert(-1 == intersection(Circle(1, 1, 3), Line(5, 3, -30), &p, &q));
  assert(0 == intersection(Circle(1, 1, 3), Line(0, 1, -4), &p, &q));
  assert(EQ(p, Point(1, 4)));
  assert(1 == intersection(Circle(1, 1, 3), Line(0, 1, -1), &p, &q));
  assert(EQ(p, Point(4, 1)));
  assert(EQ(q, Point(-2, 1)));
  assert(1 == intersection(Circle(1, 1, 3), Line(0, 1, -1), &p));
  assert(EQ(p, Point(4, 1)));

  assert(-2 == intersection(Circle(1, 1, 1), Circle(0, 0, 3), &p, &q));
  assert(-1 == intersection(Circle(0, 0, 3), Circle(1, 1, 1), &p, &q));
  assert(0 == intersection(Circle(5, 0, 4), Circle(-5, 0, 4), &p, &q));
  assert(1 == intersection(Circle(-5, 0, 5), Circle(5, 0, 5), &p, &q));
  assert(EQ(p, Point(0, 0)));
  assert(2 == intersection(Circle(-0.5, 0, 1), Circle(0.5, 0, 1), &p, &q));
  assert(EQ(p, Point(0, -sqrt(3) / 2)));
  assert(EQ(q, Point(0, sqrt(3) / 2)));
  assert(
      2 == intersection(Circle(-0.5, 0, 1), Circle(0.5, 0, 1), static_cast<Point *>(nullptr), &q)
  );
  assert(EQ(q, Point(0, sqrt(3) / 2)));

  // Each circle passes through the other's center.
  double r = 3, a = intersection_area(Circle(-r / 2, 0, r), Circle(r / 2, 0, r));
  assert(EQ(a, r * r * (2 * PI / 3 - sqrt(3) / 2)));
  return 0;
}
\end{lstlisting}
\subsection{Circle Tangents}
\label{sec:7.2.5}
Given two circles, compute their common tangent lines by returning the two points of tangency on each circle. External tangents touch both circles on the same side of the line joining their centers; internal tangents cross between the circles. A point can be treated as a zero-radius circle, so this also covers tangents from a point to a circle.

\begin{compactitem}
  \item \inlinecode{circle\_tangents(c1, r1, c2, r2, internal = false)} returns $0$, $1$, or $2$ tangent point pairs. Each pair $(p, q)$ means the tangent line touches the first circle at $p$ and the second circle at $q$. Set \inlinecode{internal = true} to get internal tangents; otherwise external tangents are returned.
  \item \inlinecode{all\_circle\_tangents(c1, r1, c2, r2)} returns all external and internal tangents.
\end{compactitem}

If the two circles are identical, there are infinitely many tangents and the functions return an empty vector. If one circle strictly contains the other, there are no external tangents; if the circles overlap, there are no internal tangents.

\textbf{Time Complexity}: $O(1)$ per call.

\textbf{Space Complexity}: $O(1)$ auxiliary, excluding the returned vector.

\begin{lstlisting}[style=cpp]
#include <algorithm>
#include <cmath>
#include <type_traits>
#include <utility>
#include <vector>

const double EPS = 1e-9;

template<typename T, typename U, typename C = std::common_type_t<T, U>>
bool EQ(T a, U b) {
  if constexpr (std::is_floating_point_v<C>) return C(a) == C(b) || std::fabs(C(a) - C(b)) <= EPS;
  return C(a) == C(b);
}

template<typename T, typename U, typename C = std::common_type_t<T, U>>
bool LT(T a, U b) {
  if constexpr (std::is_floating_point_v<C>) return C(a) < C(b) - EPS;
  return C(a) < C(b);
}

struct Point {
  double x, y;
  Point(double x = 0, double y = 0) : x(x), y(y) {}
  bool operator==(const Point &p) const { return x == p.x && y == p.y; }
  friend bool EQ(const Point &a, const Point &b) { return EQ(a.x, b.x) && EQ(a.y, b.y); }
  Point operator+(const Point &p) const { return {x + p.x, y + p.y}; }
  Point operator-(const Point &p) const { return {x - p.x, y - p.y}; }
  Point operator*(double k) const { return {x * k, y * k}; }
  Point operator/(double k) const { return {x / k, y / k}; }
  double sqnorm() const { return x * x + y * y; }
  Point rotate90() const { return {-y, x}; }
};

template<typename PtA, typename PtB>
std::vector<std::pair<Point, Point>> circle_tangents(
    const PtA &c1, double r1, const PtB &c2, double r2, bool internal = false
) {
  r1 = fabs(r1);
  r2 = fabs(r2);
  if (internal) {
    r2 = -r2;
  }
  Point d(c2.x - c1.x, c2.y - c1.y);
  double dr = r1 - r2, d2 = d.sqnorm(), h2 = d2 - dr * dr;
  if (EQ(d2, 0) || LT(h2, 0)) {
    return {};
  }
  std::vector<std::pair<Point, Point>> res;
  for (double s : {-1.0, 1.0}) {
    Point v = (d * dr + d.rotate90() * (std::sqrt(std::max(0.0, h2)) * s)) / d2;
    res.emplace_back(Point(c1.x, c1.y) + v * r1, Point(c2.x, c2.y) + v * r2);
  }
  if (EQ(h2, 0)) {
    res.pop_back();
  }
  return res;
}

template<typename PtA, typename PtB>
std::vector<std::pair<Point, Point>> all_circle_tangents(
    const PtA &c1, double r1, const PtB &c2, double r2
) {
  auto res = circle_tangents(c1, r1, c2, r2);
  auto internal = circle_tangents(c1, r1, c2, r2, true);
  res.insert(res.end(), internal.begin(), internal.end());
  return res;
}
\end{lstlisting}
\Needspace*{4\baselineskip}
\textbf{Example Usage}\par
\begin{lstlisting}[style=cpp]
#include <cassert>
using namespace std;

int main() {
  Point a(0, 0), b(4, 0);
  auto ext = circle_tangents(a, 1, b, 1);
  assert(ext.size() == 2);
  assert(EQ(ext[0].first, Point(0, -1)) && EQ(ext[0].second, Point(4, -1)));
  assert(EQ(ext[1].first, Point(0, 1)) && EQ(ext[1].second, Point(4, 1)));

  auto in = circle_tangents(a, 1, b, 1, true);
  assert(in.size() == 2);
  assert(
      EQ(in[0].first, Point(0.5, -sqrt(3.0) / 2)) && EQ(in[0].second, Point(3.5, sqrt(3.0) / 2))
  );
  assert(
      EQ(in[1].first, Point(0.5, sqrt(3.0) / 2)) && EQ(in[1].second, Point(3.5, -sqrt(3.0) / 2))
  );
  auto all = ext;
  all.insert(all.end(), in.begin(), in.end());
  assert(all_circle_tangents(a, 1, b, 1) == all);
  auto point_to_circle = circle_tangents(Point(0, 0), 0, Point(2, 0), 1);
  assert(point_to_circle.size() == 2);
  assert(
      EQ(point_to_circle[0].first, Point(0, 0)) &&
      EQ(point_to_circle[0].second, Point(1.5, -sqrt(3.0) / 2))
  );
  assert(
      EQ(point_to_circle[1].first, Point(0, 0)) &&
      EQ(point_to_circle[1].second, Point(1.5, sqrt(3.0) / 2))
  );
  assert(circle_tangents(Point(0, 0), 5, Point(1, 0), 1).empty());        // contained
  assert(circle_tangents(Point(0, 0), 2, Point(3, 0), 2, true).empty());  // overlapping
  return 0;
}
\end{lstlisting}

\section{\texorpdfstring{Polygons and Point Sets}{7.3 Polygons and Point Sets}}
\setcounter{section}{3}
\setcounter{subsection}{0}
\subsection{Polygon Sorting and Area}
\label{sec:7.3.1}
Given a list of distinct points in two dimensions, order them into a valid polygon and determine the area. The sorting comparator orders two points by the sign of their cross product around a chosen center, that is, by angle. To form a simple polygon, choose a center inside the points' convex hull that differs from every input point; the arithmetic mean used in the example is usually convenient. The area functions use the shoelace formula, summing the cross products of consecutive vertex pairs. All functions accept \inlinecode{Point}, \inlinecode{PointI}, or \inlinecode{PointL} from \secref{7.1.1}, or any compatible point type with numeric \inlinecode{.x} and \inlinecode{.y} fields and the comparison operators used by the requested operation.

\begin{compactitem}
  \item \inlinecode{cw\_comp(a, b, c)} returns whether point \inlinecode{a} compares clockwise before \inlinecode{b} about \inlinecode{c}.
  \item \inlinecode{polygon\_area\_2x(lo, hi)} returns exactly double the area of the polygon with vertices specified by the range $[{\inlinecodemath{lo}}, {\inlinecodemath{hi}})$ of points in either clockwise or counter-clockwise order. The return value is integral or floating-point, depending on the input point type. For integer vertices, divide by $2$ in the caller if the exact area is needed.
  \item \inlinecode{polygon\_area(lo, hi)} returns the area as \inlinecode{double}.
  \item \inlinecode{polygon\_centroid(lo, hi)} returns the centroid (center of mass) of a non-degenerate simple polygon with vertices in boundary order as an \inlinecode{std::pair\textless{}double, double\textgreater{}}. It uses signed shoelace weights, so either clockwise or counter-clockwise inputs will work.
\end{compactitem}

Overflow warning: \inlinecode{cw\_comp} and \inlinecode{polygon\_area\_2x} form cross products that grow like the squared coordinate magnitude (and the shoelace sum accumulates over all vertices). For integer point types use a 64-bit coordinate type (e.g. \inlinecode{PointL} from \secref{7.1.1}) for large or numerous coordinates.

\Needspace*{3\baselineskip}
\textbf{Time Complexity}:
\begin{compactitem}
  \item $O(n)$ per call to \inlinecode{polygon\_area\_2x()}, \inlinecode{polygon\_area()}, and \inlinecode{polygon\_centroid()}, where $n$ is the number of points.
  \item $O(1)$ per call to \inlinecode{cw\_comp()} and the comparators.
\end{compactitem}

\textbf{Space Complexity}: $O(1)$ auxiliary for all operations.

\begin{lstlisting}[style=cpp]
#include <cassert>
#include <cmath>
#include <stdexcept>
#include <type_traits>
#include <utility>

const double EPS = 1e-9;

template<typename T, typename U, typename C = std::common_type_t<T, U>>
bool EQ(T a, U b) {
  if constexpr (std::is_floating_point_v<C>) return C(a) == C(b) || std::fabs(C(a) - C(b)) <= EPS;
  return C(a) == C(b);
}

// Comparator (clockwise angular order about c). Exact: comparisons are pure sign tests on
// coordinate differences and the cross product, so it is a valid strict weak ordering and is
// exact for integer-coordinate points.
template<typename Pt>
bool cw_comp(const Pt &a, const Pt &b, const Pt &c) {
  if (a.x - c.x >= 0 && b.x - c.x < 0) {
    return true;
  }
  if (a.x - c.x < 0 && b.x - c.x >= 0) {
    return false;
  }
  if (a.x - c.x == 0 && b.x - c.x == 0) {
    if (a.y - c.y >= 0 || b.y - c.y >= 0) {
      return a.y > b.y;
    }
    return b.y > a.y;
  }
  auto acx = a.x - c.x, acy = a.y - c.y;
  auto bcx = b.x - c.x, bcy = b.y - c.y;
  auto det = acx * bcy - acy * bcx;  // Overflow warning.
  if (det == 0) {
    auto acnorm = acx * acx + acy * acy;  // Overflow warning.
    auto bcnorm = bcx * bcx + bcy * bcy;
    return acnorm > bcnorm;
  }
  return det < 0;
}

// Returns 2 * area. Result is exact (integer) for integer-coordinate points.
template<typename It>
auto polygon_area_2x(It lo, It hi) {
  using T = decltype(lo->x * lo->y);
  if (lo == hi) {
    return T{0};
  }
  T area = 0;
  for (It i = lo, j = hi - 1; i != hi; j = i++) {
    area += static_cast<T>(j->x - i->x) * static_cast<T>(j->y + i->y);  // Overflow warning.
  }
  return area < T{0} ? -area : area;
}

template<typename It>
double polygon_area(It lo, It hi) {
  return static_cast<double>(polygon_area_2x(lo, hi)) / 2.0;
}

template<typename It>
std::pair<double, double> polygon_centroid(It lo, It hi) {
  assert(hi - lo >= 3);
  double cx = 0, cy = 0, area2 = 0;
  for (It i = lo, j = hi - 1; i != hi; j = i++) {
    double cross = static_cast<double>(j->x) * i->y - static_cast<double>(i->x) * j->y;
    cx += (static_cast<double>(j->x) + i->x) * cross;
    cy += (static_cast<double>(j->y) + i->y) * cross;
    area2 += cross;
  }
  if (EQ(area2, 0)) {
    throw std::runtime_error("Cannot compute centroid of zero-area polygon.");
  }
  return {cx / (3 * area2), cy / (3 * area2)};
}
\end{lstlisting}
\Needspace*{4\baselineskip}
\textbf{Example Usage}\par
\begin{lstlisting}[style=cpp]
#include <algorithm>
#include <cassert>
#include <random>
#include <vector>
using namespace std;

struct Point {
  double x, y;
  Point(double x = 0, double y = 0) : x(x), y(y) {}
  bool operator==(const Point &p) const { return x == p.x && y == p.y; }
  friend bool EQ(const Point &a, const Point &b) { return EQ(a.x, b.x) && EQ(a.y, b.y); }
  bool operator!=(const Point &p) const { return !(*this == p); }
  bool operator<(const Point &p) const { return x != p.x ? x < p.x : y < p.y; }
  bool operator>(const Point &p) const { return p < *this; }
};

struct PointI {
  int x, y;
  PointI(int x = 0, int y = 0) : x(x), y(y) {}
  bool operator==(const PointI &p) const { return x == p.x && y == p.y; }
  bool operator!=(const PointI &p) const { return !(*this == p); }
  bool operator<(const PointI &p) const { return x != p.x ? x < p.x : y < p.y; }
  bool operator>(const PointI &p) const { return p < *this; }
};

template<typename It>
Point mean_center(It lo, It hi) {
  assert(lo != hi);
  double x_sum = 0, y_sum = 0, n = hi - lo;
  for (It it = lo; it != hi; ++it) {
    x_sum += it->x;
    y_sum += it->y;
  }
  return Point(x_sum / n, y_sum / n);
}

int main() {
  vector<Point> points{Point(1, 3), Point(1, 2), Point(2, 1), Point(0, 0), Point(-1, 3)};
  vector<Point> v(points);
  mt19937 rng(1234567);  // Fixed seed for reproducibility.
  shuffle(v.begin(), v.end(), rng);
  Point c = mean_center(v.begin(), v.end());
  assert(EQ(c, Point(0.6, 1.8)));
  sort(v.begin(), v.end(), [c](const Point &a, const Point &b) { return cw_comp(a, b, c); });
  assert((v == vector<Point>{{1, 3}, {1, 2}, {2, 1}, {0, 0}, {-1, 3}}));
  assert(EQ(polygon_area(v.begin(), v.end()), 5));

  sort(v.begin(), v.end(), [c](const Point &a, const Point &b) { return cw_comp(b, a, c); });
  assert((v == vector<Point>{{-1, 3}, {0, 0}, {2, 1}, {1, 2}, {1, 3}}));
  assert(EQ(polygon_area(v.begin(), v.end()), 5));

  // Integer points: polygon_area_2x is exact (no float arithmetic).
  vector<PointI> iv{{0, 0}, {4, 0}, {0, 3}};            // right triangle, area = 6
  assert(polygon_area_2x(iv.begin(), iv.end()) == 12);  // exact int
  assert(EQ(polygon_area(iv.begin(), iv.end()), 6.0));
  auto centroid = polygon_centroid(iv.begin(), iv.end());
  assert(EQ(centroid.first, 4.0 / 3.0) && EQ(centroid.second, 1.0));
  return 0;
}
\end{lstlisting}
\subsection{Point-in-Polygon}
\label{sec:7.3.2}
Given a point \inlinecode{p} and a polygon, determines whether \inlinecode{p} lies inside using ray casting. A ray cast from \inlinecode{p} must cross the polygon's boundary an odd number of times exactly when \inlinecode{p} is inside, so the function counts crossings between a horizontal ray and each polygon edge. The function is templated on the point type \inlinecode{Pt} and accepts \inlinecode{Point}, \inlinecode{PointI}, or \inlinecode{PointL} from \secref{7.1.1}, or any compatible type with numeric \inlinecode{.x} and \inlinecode{.y} fields. All comparisons are exact (no epsilon): the ray-crossing parity needs a consistent comparison to be correct, and the result is exact for integer-coordinate points.

\begin{compactitem}
  \item \inlinecode{point\_in\_polygon(p, lo, hi, include\_boundary = true)} returns whether \inlinecode{p} lies within the polygon with vertices specified by the range $[{\inlinecodemath{lo}}, {\inlinecodemath{hi}})$ of points in either clockwise or counter-clockwise order. The \inlinecode{include\_boundary} flag controls whether points exactly on an edge or vertex count as inside.
\end{compactitem}

Overflow warning: the edge orientation test forms a cross product that grows like the squared coordinate magnitude. For integer point types use a 64-bit coordinate type (e.g. \inlinecode{PointL} from \secref{7.1.1}) once coordinates exceed roughly a few tens of thousands.

\textbf{Time Complexity}: $O(n)$ per call, where $n$ is the distance between \inlinecode{lo} and \inlinecode{hi}.

\textbf{Space Complexity}: $O(1)$ auxiliary.

\begin{lstlisting}[style=cpp]
template<typename Pt>
auto cross(const Pt &a, const Pt &b, const Pt &o) {
  // Overflow risk for integer Pt: these products are ~O(max_coord^2); use int64_t if necessary.
  return (a.x - o.x) * (b.y - o.y) - (a.y - o.y) * (b.x - o.x);
}

// Detection is exact for integer-coordinate Pt.
template<typename Pt, typename It>
bool point_in_polygon(const Pt &p, It lo, It hi, bool include_boundary = true) {
  if (lo == hi) {
    return false;
  }
  bool res = false;
  for (It i = lo, j = hi - 1; i != hi; j = i++) {
    if (i->y == p.y && (i->x == p.x || (j->y == p.y && (i->x <= p.x || j->x <= p.x)))) {
      return include_boundary;
    }
    if ((p.y < i->y) != (p.y < j->y)) {
      auto det = cross(*i, *j, p);
      if (det == 0) {
        return include_boundary;
      }
      if ((0 < det) != (i->y < j->y)) {
        res = !res;
      }
    }
  }
  return res;
}
\end{lstlisting}
\Needspace*{4\baselineskip}
\textbf{Example Usage}\par
\begin{lstlisting}[style=cpp]
#include <cassert>
#include <vector>
using namespace std;

struct Point {
  double x, y;
  Point(double x = 0, double y = 0) : x(x), y(y) {}
};

struct PointI {
  int x, y;
  PointI(int x = 0, int y = 0) : x(x), y(y) {}
};

int main() {
  vector<Point> poly{Point(-1, 3), Point(1, 3), Point(2, 1), Point(0, 0)};
  assert(point_in_polygon(Point(1, 2), poly.begin(), poly.end()));
  assert(point_in_polygon(Point(0, 3), poly.begin(), poly.end()));
  assert(!point_in_polygon(Point(0, 3), poly.begin(), poly.end(), false));
  assert(!point_in_polygon(Point(0, 3.01), poly.begin(), poly.end()));
  assert(!point_in_polygon(Point(2, 2), poly.begin(), poly.end()));

  // Integer-coordinate points: exact detection.
  vector<PointI> ipoly{{0, 0}, {4, 0}, {4, 4}, {0, 4}};
  assert(point_in_polygon(PointI(2, 2), ipoly.begin(), ipoly.end()));
  assert(!point_in_polygon(PointI(4, 2), ipoly.begin(), ipoly.end(), false));
  assert(!point_in_polygon(PointI(5, 2), ipoly.begin(), ipoly.end()));
  return 0;
}
\end{lstlisting}
\subsection{Convex Hull and Diametral Pair}
\label{sec:7.3.3}
The convex hull of points in two dimensions is the smallest convex set containing them. A diametral pair is a pair of input points at maximum distance. Monotone chain computes the hull by sorting the points lexicographically and building the lower and upper boundaries in one pass each, popping any point that would fail to create a counter-clockwise turn. Rotating calipers then walks two antipodal pointers around the hull, advancing whichever increases the separation and visiting every candidate diametral pair in linear time. Both functions accept either floating-point or integral coordinates, and use exact comparisons for integral points.

\begin{compactitem}
  \item \inlinecode{convex\_hull(p)} returns the convex hull of points \inlinecode{p} in counter-clockwise order. Duplicate input points and collinear points along hull edges are omitted. To instead return the hull points in clockwise order, replace the cross product comparisons \inlinecode{\textless{}= 0} with \inlinecode{\textgreater{}= 0}.
  \item \inlinecode{diametral\_pair(p)} returns the maximum-distance pair of points in \inlinecode{p}.
\end{compactitem}

Overflow warning: \inlinecode{cross()} and the squared distance in \inlinecode{diametral\_pair()} grow like the squared coordinate magnitude. With 32-bit \inlinecode{int} coordinates they overflow once coordinates exceed a few tens of thousands; use a 64-bit (\inlinecode{int64\_t}) coordinate type for larger integer inputs.

\textbf{Time Complexity}: $O(n \log n)$ per call, where $n$ is the number of points.

\textbf{Space Complexity}: $O(n)$ auxiliary for storage of the convex hull.

\begin{lstlisting}[style=cpp]
#include <algorithm>
#include <cmath>
#include <utility>
#include <vector>

template<typename Pt>
auto cross(const Pt &a, const Pt &b, const Pt &o) {
  // Overflow risk for integer Pt: these products are ~O(max_coord^2); use int64_t if necessary.
  return (a.x - o.x) * (b.y - o.y) - (a.y - o.y) * (b.x - o.x);
}

// Convex hull: exact for integer-coordinate points.
template<typename Pt>
std::vector<Pt> convex_hull(std::vector<Pt> p) {
  std::sort(p.begin(), p.end());
  auto same_point = [](const Pt &a, const Pt &b) { return !(a < b) && !(b < a); };
  p.erase(std::unique(p.begin(), p.end(), same_point), p.end());
  int n = static_cast<int>(p.size());
  if (n <= 1) {
    return p;
  }
  int k = 0;
  std::vector<Pt> res(2 * n);
  for (const Pt &q : p) {
    while (k >= 2 && cross(res[k - 1], q, res[k - 2]) <= 0) {
      k--;
    }
    res[k++] = q;
  }
  int t = k + 1;
  for (int i = n - 2; i >= 0; i--) {
    while (k >= t && cross(res[k - 1], p[i], res[k - 2]) <= 0) {
      k--;
    }
    res[k++] = p[i];
  }
  res.resize(k - 1);
  return res;
}

// Diametral pair: squared-distance comparisons are exact for integer-coordinate points.
template<typename Pt>
std::pair<Pt, Pt> diametral_pair(const std::vector<Pt> &p) {
  auto h = convex_hull(p);
  int m = static_cast<int>(h.size());
  if (m == 0) {
    return std::pair<Pt, Pt>{};
  }
  if (m == 1) {
    return std::pair<Pt, Pt>{h[0], h[0]};
  }
  if (m == 2) {
    return std::pair<Pt, Pt>{h[0], h[1]};
  }
  int k = 1;
  while (std::abs(cross(h[0], h[(k + 1) % m], h[m - 1])) > std::abs(cross(h[0], h[k], h[m - 1]))) {
    k++;
  }
  auto sqdist = [](const Pt &a, const Pt &b) {
    // Overflow risk for integer Pt: ~O(max_coord^2); use int64_t if necessary.
    auto dx = a.x - b.x, dy = a.y - b.y;
    return dx * dx + dy * dy;
  };
  auto maxsq = sqdist(h[0], h[0]);
  std::pair<Pt, Pt> res{h[0], h[0]};
  for (int i = 0, j = k; i <= k && j < m; i++) {
    auto d = sqdist(h[i], h[j]);
    if (d > maxsq) {
      maxsq = d;
      res = {h[i], h[j]};
    }
    while (j < m && std::abs(cross(h[(i + 1) % m], h[(j + 1) % m], h[i])) >
                        std::abs(cross(h[(i + 1) % m], h[j], h[i]))) {
      d = sqdist(h[i], h[(j + 1) % m]);
      if (d > maxsq) {
        maxsq = d;
        res = {h[i], h[(j + 1) % m]};
      }
      j++;
    }
  }
  return res;
}
\end{lstlisting}
\Needspace*{4\baselineskip}
\textbf{Example Usage}\par
\begin{lstlisting}[style=cpp]
#include <cassert>
#include <random>
#include <vector>
using namespace std;

struct PointI {
  int x, y;
  PointI(int x = 0, int y = 0) : x(x), y(y) {}
  bool operator==(const PointI &o) const { return x == o.x && y == o.y; }
  bool operator!=(const PointI &o) const { return !(*this == o); }
  bool operator<(const PointI &o) const { return x != o.x ? x < o.x : y < o.y; }
  bool operator>(const PointI &o) const { return o < *this; }
};

int main() {
  {
    vector<PointI> v{{1, 3}, {1, 2}, {2, 1}, {0, 0}, {-1, 3}};
    mt19937 rng(1234567);  // Fixed seed for reproducibility.
    shuffle(v.begin(), v.end(), rng);
    vector<PointI> h{{-1, 3}, {0, 0}, {2, 1}, {1, 3}};
    assert(convex_hull(v) == h);
  }
  {
    auto [p1, p2] = diametral_pair(vector<PointI>{{0, 0}, {3, 0}, {0, 3}, {1, 1}, {4, 4}});
    assert(p1 == PointI(0, 0) && p2 == PointI(4, 4));
  }
  {
    vector<PointI> v{{0, 0}, {4, 0}, {4, 4}, {0, 4}, {2, 2}};
    auto h = convex_hull(v);
    assert(h.size() == 4);  // interior point (2,2) excluded
    auto [p1, p2] = diametral_pair(v);
    // diametral pair is exact
    assert(
        (p1 == PointI(0, 0) && p2 == PointI(4, 4)) || (p1 == PointI(4, 4) && p2 == PointI(0, 0)) ||
        (p1 == PointI(4, 0) && p2 == PointI(0, 4)) || (p1 == PointI(0, 4) && p2 == PointI(4, 0))
    );
  }
  {
    assert((convex_hull(vector<PointI>{{2, 3}, {2, 3}, {2, 3}}) == vector<PointI>{{2, 3}}));
  }
  return 0;
}
\end{lstlisting}
\subsection{Convex Polygon Queries}
\label{sec:7.3.4}
Given a convex polygon in counter-clockwise order, answer containment and tangent/intersection queries. Point containment is logarithmic by splitting the polygon into triangles sharing \inlinecode{poly[0]}. The line and tangent helpers below are written as simple linear scans; this is shorter and less assumption-heavy than the classic logarithmic versions, while still documenting the exact query semantics in one place.

\begin{compactitem}
  \item \inlinecode{point\_in\_convex\_polygon(poly, p, include\_boundary = true)} returns whether point \inlinecode{p} lies inside the convex polygon \inlinecode{poly}. The polygon must be in counter-clockwise order, with no repeated final vertex. Points on the boundary count as inside unless \inlinecode{include\_boundary} is set to \inlinecode{false}.
  \item \inlinecode{line\_convex\_polygon\_intersection(poly, a, b)} describes where the infinite directed line through \inlinecode{a} $\to$ \inlinecode{b} hits the polygon: $(-1, -1)$ for no intersection, $(i, -1)$ for touching vertex $i$, $(i, i)$ for containing side $i \to i+1$, or $(i, j)$ for crossing sides $i \to i+1$ and $j \to j+1$, with side indices taken modulo the number of vertices. The polygon must have at least three vertices, and \inlinecode{a} and \inlinecode{b} must differ. When the line crosses two sides, their order is determined using a floating-point line parameter.
  \item \inlinecode{convex\_polygon\_tangents(poly, p)} returns the two tangent vertex indices from an outside point \inlinecode{p} as (\inlinecode{left}, \inlinecode{right}). The left tangent has all polygon vertices on or to the left of the directed line \inlinecode{p} $\to$ \inlinecode{poly[left]}; the right tangent is analogous for the right side.
\end{compactitem}

Overflow warning: For integer-coordinate inputs, all orientation tests are exact provided the cross products do not overflow.

\Needspace*{3\baselineskip}
\textbf{Time Complexity}:
\begin{compactitem}
  \item $O(\log n)$ per call to \inlinecode{point\_in\_convex\_polygon()}, where $n$ is the number of polygon vertices.
  \item $O(n)$ per call to \inlinecode{line\_convex\_polygon\_intersection} and \inlinecode{convex\_polygon\_tangents}.
\end{compactitem}

\textbf{Space Complexity}: $O(1)$ auxiliary.

\begin{lstlisting}[style=cpp]
#include <algorithm>
#include <cassert>
#include <utility>
#include <vector>

template<typename Pt>
auto cross(const Pt &a, const Pt &b, const Pt &o) {
  // Overflow risk for integer Pt: these products are ~O(max_coord^2); use int64_t if necessary.
  return (a.x - o.x) * (b.y - o.y) - (a.y - o.y) * (b.x - o.x);
}

template<typename T>
int sgn(const T &x) {
  return (T{0} < x) - (x < T{0});
}

template<typename Pt>
bool point_on_segment(const Pt &p, const Pt &a, const Pt &b) {
  return cross(a, b, p) == 0 && std::min(a.x, b.x) <= p.x && p.x <= std::max(a.x, b.x) &&
         std::min(a.y, b.y) <= p.y && p.y <= std::max(a.y, b.y);
}

template<typename Pt>
bool point_in_convex_polygon(
    const std::vector<Pt> &poly, const Pt &p, bool include_boundary = true
) {
  int n = static_cast<int>(poly.size());
  if (n == 0) {
    return false;
  }
  if (n == 1) {
    return include_boundary && p.x == poly[0].x && p.y == poly[0].y;
  }
  if (n == 2) {
    return include_boundary && point_on_segment(p, poly[0], poly[1]);
  }
  int left = sgn(cross(poly[1], p, poly[0]));
  int right = sgn(cross(poly[n - 1], p, poly[0]));
  if (left < 0 || right > 0) {
    return false;
  }
  if (left == 0) {
    return include_boundary && point_on_segment(p, poly[0], poly[1]);
  }
  if (right == 0) {
    return include_boundary && point_on_segment(p, poly[0], poly[n - 1]);
  }
  int lo = 1, hi = n - 1;
  while (hi - lo > 1) {
    int mid = lo + (hi - lo) / 2;
    if (cross(poly[mid], p, poly[0]) >= 0) {
      lo = mid;
    } else {
      hi = mid;
    }
  }
  int s = sgn(cross(poly[lo + 1], p, poly[lo]));
  return include_boundary ? s >= 0 : s > 0;
}

template<typename Pt>
std::pair<int, int> line_convex_polygon_intersection(
    const std::vector<Pt> &poly, const Pt &a, const Pt &b
) {
  int n = static_cast<int>(poly.size());
  assert(n >= 3 && (a.x != b.x || a.y != b.y));
  std::vector<int> on, cross_edges;
  bool has_pos = false, has_neg = false;
  for (int i = 0; i < n; i++) {
    int s = sgn(cross(b, poly[i], a));
    has_pos |= s > 0;
    has_neg |= s < 0;
    if (s == 0) {
      on.push_back(i);
    }
  }
  if (!has_pos && !has_neg) {
    return n >= 2 ? std::make_pair(0, 0) : std::make_pair(-1, -1);
  }
  if (!(has_pos && has_neg)) {
    for (int i = 0; i < static_cast<int>(on.size()); i++) {
      int j = (i + 1) % static_cast<int>(on.size());
      if ((on[i] + 1) % n == on[j]) {
        return {on[i], on[i]};
      }
    }
    return on.empty() ? std::make_pair(-1, -1) : std::make_pair(on[0], -1);
  }
  for (int i = 0; i < n; i++) {
    int j = (i + 1) % n;
    int si = sgn(cross(b, poly[i], a)), sj = sgn(cross(b, poly[j], a));
    if (si == 0 || si * sj < 0) {
      cross_edges.push_back(i);
    }
  }
  if (cross_edges.size() == 1) {
    return {cross_edges[0], -1};
  }
  auto line_parameter = [&](int i) {
    int j = (i + 1) % n;
    auto ex = poly[j].x - poly[i].x, ey = poly[j].y - poly[i].y;
    auto ax = a.x, ay = a.y, dx = b.x - a.x, dy = b.y - a.y;
    auto det = dx * ey - dy * ex;
    return static_cast<double>((poly[i].x - ax) * ey - (poly[i].y - ay) * ex) / det;
  };
  if (line_parameter(cross_edges[1]) < line_parameter(cross_edges[0])) {
    std::swap(cross_edges[0], cross_edges[1]);
  }
  return {cross_edges[0], cross_edges[1]};
}

template<typename Pt>
std::pair<int, int> convex_polygon_tangents(const std::vector<Pt> &poly, const Pt &p) {
  int n = static_cast<int>(poly.size());
  std::pair<int, int> res{-1, -1};
  for (int i = 0; i < n; i++) {
    int prev = (i + n - 1) % n, next = (i + 1) % n;
    int s1 = sgn(cross(poly[i], poly[prev], p));
    int s2 = sgn(cross(poly[i], poly[next], p));
    bool left = s1 >= 0 && s2 >= 0;
    bool right = s1 <= 0 && s2 <= 0;
    if (left) {
      res.first = i;
    }
    if (right) {
      res.second = i;
    }
  }
  return res;
}
\end{lstlisting}
\Needspace*{4\baselineskip}
\textbf{Example Usage}\par
\begin{lstlisting}[style=cpp]
#include <algorithm>
#include <cstdint>
using namespace std;

struct PointI {
  int64_t x, y;
  PointI(int64_t x = 0, int64_t y = 0) : x(x), y(y) {}
};

int main() {
  vector<PointI> square{{0, 0}, {4, 0}, {4, 4}, {0, 4}};
  assert(point_in_convex_polygon(square, PointI(2, 2)));
  assert(point_in_convex_polygon(square, PointI(4, 2)));
  assert(!point_in_convex_polygon(square, PointI(4, 2), false));
  assert(!point_in_convex_polygon(square, PointI(5, 2)));

  vector<PointI> pentagon{{0, 0}, {3, 0}, {5, 2}, {2, 5}, {-1, 3}};
  assert(point_in_convex_polygon(pentagon, PointI(2, 3)));
  assert(point_in_convex_polygon(pentagon, PointI(0, 0)));
  assert(!point_in_convex_polygon(pentagon, PointI(0, 0), false));
  assert(!point_in_convex_polygon(pentagon, PointI(5, 4)));

  assert(line_convex_polygon_intersection(square, PointI(-1, 2), PointI(5, 2)) == make_pair(3, 1));
  assert(line_convex_polygon_intersection(square, PointI(0, 0), PointI(4, 0)) == make_pair(0, 0));
  assert(line_convex_polygon_intersection(square, PointI(5, 0), PointI(5, 4)) == make_pair(-1, -1));

  // From the point to the right of the square, the upper-right and lower-right vertices are
  // tangent.
  assert(convex_polygon_tangents(square, PointI(6, 2)) == make_pair(2, 1));
  return 0;
}
\end{lstlisting}
\subsection{Minimum Enclosing Circle}
\label{sec:7.3.5}
Given a set of points in two dimensions, finds the unique circle of minimum radius (equivalently, minimum area) containing all points.

This implementation uses Welzl's randomized incremental algorithm. The points are shuffled and processed one at a time while maintaining the current minimum enclosing circle. Whenever a point lies outside the current circle, the circle is rebuilt with that point constrained to lie on the boundary. The final circle is determined by at most three boundary points.

\begin{compactitem}
  \item \inlinecode{min\_enclosing\_circle(p)} returns the minimum enclosing circle for the points in \inlinecode{p}. The point type may be any type exposing numeric \inlinecode{.x} and \inlinecode{.y} members. The returned circle always uses \inlinecode{double} coordinates, so integer-coordinate inputs are accepted.
\end{compactitem}

\textbf{Time Complexity}: $O(n)$ expected time per call, where $n$ is the number of points, or $O(n^{3})$ in the worst case for a particular shuffle order.

\textbf{Space Complexity}: $O(n)$ auxiliary for the shuffled copy of the points.

\begin{lstlisting}[style=cpp]
#include <algorithm>
#include <chrono>
#include <cmath>
#include <random>
#include <stdexcept>
#include <type_traits>
#include <vector>

const double EPS = 1e-9;

template<typename T, typename U, typename C = std::common_type_t<T, U>>
bool EQ(T a, U b) {
  if constexpr (std::is_floating_point_v<C>) return C(a) == C(b) || std::fabs(C(a) - C(b)) <= EPS;
  return C(a) == C(b);
}

template<typename T, typename U, typename C = std::common_type_t<T, U>>
bool LT(T a, U b) {
  if constexpr (std::is_floating_point_v<C>) return C(a) < C(b) - EPS;
  return C(a) < C(b);
}

template<typename T, typename U>
bool LE(T a, U b) {
  return !LT(b, a);
}

double sqnorm(double x, double y) {
  return x * x + y * y;
}

// SFINAE guard: valid only when Pt exposes numeric .x/.y members. This keeps the templated point
// constructors from hijacking calls like Circle(double, double, double).
template<typename Pt>
using if_point = decltype(std::declval<const Pt &>().x, std::declval<const Pt &>().y, void());

struct Circle {
  double h, k, r;

  Circle() : h(0), k(0), r(0) {}
  Circle(double h, double k, double r) : h(h), k(k), r(fabs(r)) {}

  template<typename Pt, typename = if_point<Pt>>
  Circle(const Pt &a, const Pt &b) {
    h = (a.x + b.x) / 2.0;
    k = (a.y + b.y) / 2.0;
    r = std::hypot(a.x - h, a.y - k);
  }

  template<typename Pt, typename = if_point<Pt>>
  Circle(const Pt &a, const Pt &b, const Pt &c) {
    double an = sqnorm(b.x - c.x, b.y - c.y);
    double bn = sqnorm(a.x - c.x, a.y - c.y);
    double cn = sqnorm(a.x - b.x, a.y - b.y);
    double wa = an * (bn + cn - an), wb = bn * (an + cn - bn), wc = cn * (an + bn - cn);
    double w = wa + wb + wc;
    if (EQ(w, 0)) {
      throw std::runtime_error("No circumcircle from collinear points.");
    }
    h = (wa * a.x + wb * b.x + wc * c.x) / w;
    k = (wa * a.y + wb * b.y + wc * c.y) / w;
    r = std::hypot(a.x - h, a.y - k);
  }

  template<typename Pt, typename = if_point<Pt>>
  bool contains(const Pt &p) const {
    return LE(sqnorm(p.x - h, p.y - k), r * r);
  }
};

// Input points can be any type with .x and .y fields; Circle output is always double.
template<typename Pt>
Circle min_enclosing_circle(std::vector<Pt> p) {
  if (p.empty()) {
    return Circle(0, 0, 0);
  }
  if (p.size() == 1) {
    return Circle(p[0].x, p[0].y, 0);
  }
  static std::mt19937 rng(std::chrono::steady_clock::now().time_since_epoch().count());
  std::shuffle(p.begin(), p.end(), rng);
  Circle res(p[0], p[1]);
  for (int i = 2; i < static_cast<int>(p.size()); i++) {
    if (res.contains(p[i])) {
      continue;
    }
    res = Circle(p[0], p[i]);
    for (int j = 1; j < i; j++) {
      if (res.contains(p[j])) {
        continue;
      }
      res = Circle(p[i], p[j]);
      for (int k = 0; k < j; k++) {
        if (!res.contains(p[k])) {
          res = Circle(p[i], p[j], p[k]);
        }
      }
    }
  }
  return res;
}
\end{lstlisting}
\Needspace*{4\baselineskip}
\textbf{Example Usage}\par
\begin{lstlisting}[style=cpp]
#include <cassert>
using namespace std;

struct Point {
  double x, y;
  Point(double x = 0, double y = 0) : x(x), y(y) {}
  bool operator==(const Point &p) const { return x == p.x && y == p.y; }
};

struct PointI {
  int x, y;
  PointI(int x = 0, int y = 0) : x(x), y(y) {}
};

int main() {
  vector<Point> empty;
  Circle none = min_enclosing_circle(empty);
  assert(EQ(none.h, 0) && EQ(none.k, 0) && EQ(none.r, 0));

  vector<Point> one{{2, 3}};
  Circle single = min_enclosing_circle(one);
  assert(EQ(single.h, 2) && EQ(single.k, 3) && EQ(single.r, 0));

  vector<Point> two{{0, 0}, {4, 0}};
  Circle diameter = min_enclosing_circle(two);
  assert(EQ(diameter.h, 2) && EQ(diameter.k, 0) && EQ(diameter.r, 2));

  vector<Point> v{{0, 0}, {0, 1}, {1, 0}, {1, 1}};
  Circle res = min_enclosing_circle(v);
  assert(EQ(res.h, 0.5) && EQ(res.k, 0.5) && EQ(res.r, 1 / sqrt(2)));

  // Integer-coordinate input: Circle output is always double.
  vector<PointI> iv{{0, 0}, {0, 2}, {2, 0}, {2, 2}};
  Circle ic = min_enclosing_circle(iv);
  assert(EQ(ic.h, 1.0) && EQ(ic.k, 1.0));
  return 0;
}
\end{lstlisting}
\subsection{Closest Pair}
\label{sec:7.3.6}
Given a list of points in two dimensions, finds the closest pair using a divide-and-conquer algorithm. The points are split in half by $x$-coordinate and each half is solved recursively; the combining step then only needs to examine points within the best distance so far of the dividing line, where each point in this strip is compared to a constant number of $y$-ordered neighbors.

\begin{compactitem}
  \item \inlinecode{closest\_pair(p, \&res)} returns the minimum squared distance between any two points in \inlinecode{p}. If \inlinecode{res} is non-null, one closest pair is stored there in lexicographic order. With fewer than two points, the maximum value of the squared-distance type is returned and \inlinecode{res} is unchanged. The point type may be any type exposing numeric \inlinecode{.x} and \inlinecode{.y} members and a lexicographic \inlinecode{operator\textless{}}. The distance preserves the coordinate arithmetic type, and the returned pair preserves the point type. For integer coordinates, the distance is exact provided intermediate products do not overflow.
\end{compactitem}

Overflow warning: squared distances grow like the square of the coordinate magnitude. For integer point types, use a 64-bit coordinate type (e.g. \inlinecode{PointL} from \secref{7.1.1}) when coordinates may exceed a few tens of thousands.

\textbf{Time Complexity}: $O(n \log n)$ per call, where $n$ is the number of points.

\Needspace*{3\baselineskip}
\textbf{Space Complexity}:
\begin{compactitem}
  \item $O(n)$ auxiliary heap space.
  \item $O(\log n)$ auxiliary stack space.
\end{compactitem}

\begin{lstlisting}[style=cpp]
#include <algorithm>
#include <iterator>
#include <limits>
#include <utility>
#include <vector>

template<typename Pt>
auto sqdist(const Pt &a, const Pt &b) {
  auto dx = b.x - a.x, dy = b.y - a.y;
  return dx * dx + dy * dy;  // Overflow warning.
}

template<typename It, typename T, typename Pt = typename std::iterator_traits<It>::value_type>
void closest_pair_rec(It lo, It hi, std::vector<Pt> &tmp, T &best, std::pair<Pt, Pt> *res) {
  int n = hi - lo;
  auto by_y = [](const Pt &a, const Pt &b) { return a.y != b.y ? a.y < b.y : a.x < b.x; };
  if (n <= 3) {
    for (It i = lo; i != hi; ++i) {
      for (It j = i + 1; j != hi; ++j) {
        T d = sqdist(*i, *j);
        if (d < best) {
          best = d;
          if (res != nullptr) {
            *res = std::minmax(*i, *j);
          }
        }
      }
    }
    std::sort(lo, hi, by_y);
    return;
  }
  It mid = lo + n / 2;
  auto midx = mid->x;
  closest_pair_rec(lo, mid, tmp, best, res);
  closest_pair_rec(mid, hi, tmp, best, res);
  // Each half is now y-sorted, so merge them before examining the center strip.
  tmp.clear();
  std::merge(lo, mid, mid, hi, std::back_inserter(tmp), by_y);
  std::move(tmp.begin(), tmp.end(), lo);
  tmp.clear();
  for (It it = lo; it != hi; ++it) {
    auto dx = it->x - midx;
    if (dx * dx < best) {
      tmp.push_back(*it);
    }
  }
  for (int i = 0; i < static_cast<int>(tmp.size()); i++) {
    for (int j = i + 1; j < static_cast<int>(tmp.size()); j++) {
      auto dy = tmp[j].y - tmp[i].y;
      if (dy * dy >= best) {
        break;
      }
      T d = sqdist(tmp[i], tmp[j]);
      if (d < best) {
        best = d;
        if (res != nullptr) {
          *res = std::minmax(tmp[i], tmp[j]);
        }
      }
    }
  }
  tmp.clear();
}

template<typename Pt>
auto closest_pair(std::vector<Pt> p, std::pair<Pt, Pt> *res = nullptr) {
  using T = decltype(sqdist(p[0], p[0]));
  T best = std::numeric_limits<T>::max();
  std::sort(p.begin(), p.end(), [](const Pt &a, const Pt &b) {
    return a.x != b.x ? a.x < b.x : a.y < b.y;
  });
  std::vector<Pt> tmp;
  tmp.reserve(p.size());
  closest_pair_rec(p.begin(), p.end(), tmp, best, res);
  return best;
}
\end{lstlisting}
\Needspace*{4\baselineskip}
\textbf{Example Usage}\par
\begin{lstlisting}[style=cpp]
#include <cassert>
#include <cmath>
#include <vector>
using namespace std;

bool EQ(double a, double b) {
  return fabs(a - b) < 1e-9;
}

struct Point {
  double x, y;
  Point(double x = 0, double y = 0) : x(x), y(y) {}
  bool operator==(const Point &p) const { return x == p.x && y == p.y; }
  bool operator<(const Point &p) const { return x != p.x ? x < p.x : y < p.y; }
};

struct PointI {
  int x, y;
  PointI(int x = 0, int y = 0) : x(x), y(y) {}
  bool operator==(const PointI &p) const { return x == p.x && y == p.y; }
  bool operator<(const PointI &o) const { return x != o.x ? x < o.x : y < o.y; }
};

int main() {
  vector<Point> v{{2, 3}, {12, 30}, {40, 50}, {5, 1}, {12, 10}, {3, 4}};
  pair<Point, Point> res;
  assert(EQ(closest_pair(v, &res), 2));
  auto [p1, p2] = res;
  assert(p1 == Point(2, 3) && p2 == Point(3, 4));

  // Integer-coordinate input: exact pair selection, int squared distance returned.
  vector<PointI> iv{{0, 0}, {10, 10}, {3, 4}};
  pair<PointI, PointI> ires;
  assert(closest_pair(iv, &ires) == 25);
  auto [i1, i2] = ires;
  assert(i1 == PointI(0, 0) && i2 == PointI(3, 4));
  return 0;
}
\end{lstlisting}
\subsection{Maximum Collinear Points}
\label{sec:7.3.7}
Given a set of points with integer coordinates, compute the maximum number of points that lie on one line. For every anchor point, normalize the direction vector to every other point by dividing by \inlinecode{gcd(abs(dx), abs(dy))} and choosing a canonical sign; equal normalized directions from the same anchor are collinear with it.

\begin{compactitem}
  \item \inlinecode{max\_collinear\_points(p)} returns the largest number of collinear points. Duplicate points are counted as separate input points.
\end{compactitem}

Overflow warning: coordinate differences and their absolute values must fit in \inlinecode{int64\_t}.

\textbf{Time Complexity}: $O(n^{2} \log n)$ per call, where $n$ is the number of points.

\textbf{Space Complexity}: $O(n)$ auxiliary.

\begin{lstlisting}[style=cpp]
#include <algorithm>
#include <cmath>
#include <cstdint>
#include <map>
#include <numeric>
#include <utility>
#include <vector>

struct Point {
  int64_t x, y;
  Point(int64_t x = 0, int64_t y = 0) : x(x), y(y) {}
  bool operator==(const Point &p) const { return x == p.x && y == p.y; }
};

int max_collinear_points(const std::vector<Point> &p) {
  int n = static_cast<int>(p.size()), ans = std::min(n, 1);
  for (int i = 0; i < n; i++) {
    std::map<std::pair<int64_t, int64_t>, int> cnt;
    int duplicates = 0, best = 0;
    for (int j = 0; j < n; j++) {
      if (i == j) {
        continue;
      }
      int64_t dx = p[j].x - p[i].x, dy = p[j].y - p[i].y;  // Overflow warning.
      if (dx == 0 && dy == 0) {
        duplicates++;
        continue;
      }
      int64_t g = std::gcd(std::abs(dx), std::abs(dy));
      dx /= g;
      dy /= g;
      if (dx < 0 || (dx == 0 && dy < 0)) {
        dx = -dx;
        dy = -dy;
      }
      best = std::max(best, ++cnt[{dx, dy}]);
    }
    ans = std::max(ans, best + duplicates + 1);
  }
  return ans;
}
\end{lstlisting}
\Needspace*{4\baselineskip}
\textbf{Example Usage}\par
\begin{lstlisting}[style=cpp]
#include <cassert>
using namespace std;

int main() {
  vector<Point> p{{0, 0}, {1, 1}, {2, 2}, {3, 5}, {4, 6}, {3, 3}};
  assert(max_collinear_points(p) == 4);
  p.emplace_back(1, 1);
  assert(max_collinear_points(p) == 5);
  assert(max_collinear_points({}) == 0);
  assert(max_collinear_points({Point(2, 3)}) == 1);
  return 0;
}
\end{lstlisting}
\subsection{Rectangle Coverage}
\label{sec:7.3.8}
Given weighted axis-aligned rectangles, iterate over every elementary cell induced by their coordinate-compressed boundaries. Within each such cell, the accumulated weight is constant, so the callback can compute rectangle union area, exactly-$k$ coverage area, threshold area, maximum overlap, weighted integrals, coverage histograms, or other offline coverage statistics.

Rectangles are half-open, $[x_1, x_2) \times [y_1, y_2)$, so adjacent rectangles sharing only an edge do not double-count any area. This also matches the difference-array update: add at the lower boundaries and subtract at the first coordinate after the rectangle.

\begin{compactitem}
  \item \inlinecode{for\_each\_rect\_cell(rects, visit)} calls \inlinecode{visit(x1, y1, x2, y2, weight)} once for each positive-area compressed cell inside the bounding box of the rectangles, where \inlinecode{weight} is the sum of weights covering that cell. Cells are visited in increasing \inlinecode{x1}, then increasing \inlinecode{y1}. Each rectangle must satisfy ${\inlinecodemath{x1}} < {\inlinecodemath{x2}}$ and ${\inlinecodemath{y1}} < {\inlinecodemath{y2}}$.
\end{compactitem}

Enumerating cells costs $O(n^{2})$ time and memory, which is the price of answering any statistic at all rather than one chosen in advance. When only the union area or perimeter is wanted, the sweep of section \secref{7.3.9} computes it in $O(n \log n)$ time and $O(n)$ memory, which is the only option once the rectangle count reaches the thousands.

Overflow warning: Coordinate differences, cell areas, and accumulated weights must fit in \inlinecode{int64\_t}.

\textbf{Time Complexity}: $O(n^{2} + n \log n)$ per call, where $n$ is the number of rectangles.

\textbf{Space Complexity}: $O(n^{2})$ auxiliary.

\begin{lstlisting}[style=cpp]
#include <algorithm>
#include <cassert>
#include <cstdint>
#include <vector>

struct Rect {
  int64_t x1, y1, x2, y2, w = 1;
};

template<typename Fn>
void for_each_rect_cell(const std::vector<Rect> &rects, Fn visit) {
  assert(std::all_of(rects.begin(), rects.end(), [](const Rect &r) {
    return r.x1 < r.x2 && r.y1 < r.y2;
  }));
  std::vector<int64_t> xs, ys;
  for (const Rect &r : rects) {
    xs.push_back(r.x1);
    xs.push_back(r.x2);
    ys.push_back(r.y1);
    ys.push_back(r.y2);
  }
  std::sort(xs.begin(), xs.end());
  std::sort(ys.begin(), ys.end());
  xs.erase(std::unique(xs.begin(), xs.end()), xs.end());
  ys.erase(std::unique(ys.begin(), ys.end()), ys.end());
  int X = static_cast<int>(xs.size()), Y = static_cast<int>(ys.size());
  std::vector<std::vector<int64_t>> diff(X + 1, std::vector<int64_t>(Y + 1));
  auto index = [](const std::vector<int64_t> &v, int64_t x) {
    return static_cast<int>(std::lower_bound(v.begin(), v.end(), x) - v.begin());
  };
  for (const Rect &r : rects) {
    int x1 = index(xs, r.x1), x2 = index(xs, r.x2);
    int y1 = index(ys, r.y1), y2 = index(ys, r.y2);
    diff[x1][y1] += r.w;  // Overflow warning.
    diff[x2][y1] -= r.w;
    diff[x1][y2] -= r.w;
    diff[x2][y2] += r.w;
  }
  for (int x = 0; x < X; x++) {
    for (int y = 0; y < Y; y++) {
      diff[x][y] += (x > 0) ? diff[x - 1][y] : 0;
      diff[x][y] += (y > 0) ? diff[x][y - 1] : 0;
      diff[x][y] -= (x > 0 && y > 0) ? diff[x - 1][y - 1] : 0;
      if (x + 1 < X && y + 1 < Y) {
        visit(xs[x], ys[y], xs[x + 1], ys[y + 1], diff[x][y]);
      }
    }
  }
}
\end{lstlisting}
\Needspace*{4\baselineskip}
\textbf{Example Usage}\par
\begin{lstlisting}[style=cpp]
#include <cassert>
using namespace std;

int main() {
  auto rect_area_if = [](const vector<Rect> &rects, auto pred) {
    int64_t area = 0;
    for_each_rect_cell(rects, [&](int64_t x1, int64_t y1, int64_t x2, int64_t y2, int64_t w) {
      if (pred(w)) {
        area += (x2 - x1) * (y2 - y1);  // Overflow warning.
      }
    });
    return area;
  };

  // Rectangles A=[0,2)x[0,2), B=[1,3)x[1,3), and C=[2,4)x[0,1) all with default weight 1.
  //
  // y=3 +---+---+---+---+
  //     |   | B | B |   |
  // y=2 +---+---+---+---+
  //     | A |A/B| B |   |
  // y=1 +---+---+---+---+
  //     | A | A | C | C |
  // y=0 +---+---+---+---+
  //   x=0   1   2   3   4
  vector<Rect> rects{{0, 0, 2, 2}, {1, 1, 3, 3}, {2, 0, 4, 1}};
  assert(rect_area_if(rects, [](int64_t w) { return w > 0; }) == 9);
  assert(rect_area_if(rects, [](int64_t w) { return w == 2; }) == 1);

  // Rectangles A=[0,4)x[0,3) of w=2, B=[2,6)x[1,4) of w=3, and C=[1,5)x[2,5) of w=1.
  //
  // y=5 +---+---+---+---+---+---+
  //     |   | 1 | 1 | 1 | 1 |   |
  // y=4 +---+---+---+---+---+---+
  //     |   | 1 | 4 | 4 | 4 | 3 |
  // y=3 +---+---+---+---+---+---+
  //     | 2 | 3 | 6 | 6 | 4 | 3 |
  // y=2 +---+---+---+---+---+---+
  //     | 2 | 2 | 5 | 5 | 3 | 3 |
  // y=1 +---+---+---+---+---+---+
  //     | 2 | 2 | 2 | 2 |   |   |
  // y=0 +---+---+---+---+---+---+
  //   x=0   1   2   3   4   5   6
  vector<Rect> weighted{{0, 0, 4, 3, 2}, {2, 1, 6, 4, 3}, {1, 2, 5, 5, 1}};
  assert(rect_area_if(weighted, [](int64_t w) { return w >= 3; }) == 13);
  assert(rect_area_if(weighted, [](int64_t w) { return w >= 5; }) == 4);

  int64_t max_weight = 0, integral = 0;
  for_each_rect_cell(weighted, [&](int64_t x1, int64_t y1, int64_t x2, int64_t y2, int64_t w) {
    max_weight = max(max_weight, w);
    integral += w * (x2 - x1) * (y2 - y1);  // Overflow warning.
  });
  assert(max_weight == 6);
  assert(integral == 12 * 2 + 12 * 3 + 12);
  return 0;
}
\end{lstlisting}
\subsection{Rectangle Union}
\label{sec:7.3.9}
Given axis-aligned rectangles that may overlap arbitrarily, measure their union. Klee's algorithm sweeps a vertical line left to right, maintaining how much of that line the rectangles currently cover. That covered length is constant between consecutive events, so each strip contributes the length times its width, and every other measure below is another aggregate of the same cross section.

A segment tree over the compressed $y$ coordinates holds the state, storing at each node the number of intervals covering its whole span and the length of that span which is covered. It never pushes anything down: a node's count refers only to intervals covering it entirely, so a node with a positive count reports its full span while one with a zero count defers to its children. Every addition is matched by a later removal, so counts never go negative and no lazy tag is needed.

Three further aggregates ride along on the same nodes. Doubly covered length shifts the recurrence by one, a node covered twice reporting its whole span and a node covered once reporting whatever its children cover at all. The weighted maximum is the largest root-to-leaf sum of weights, so each node reports its own weight plus the larger of its children's; left at the default weight of one per rectangle, that is simply the greatest number of rectangles covering a point. The perimeter needs the number of maximal covered runs, plus whether a node's first and last units are covered so that a run crossing a boundary is not counted twice: vertical edges are then the change in covered length at each event, and horizontal edges are twice the run count times the strip width.

Rectangles are half-open, $[x_1, x_2) \times [y_1, y_2)$, so adjacent rectangles sharing only an edge neither double-count any area nor leave a seam.

\begin{compactitem}
  \item \inlinecode{rect\_union(rects)} returns a \inlinecode{UnionMeasure} holding the \inlinecode{area} covered at least once, the \inlinecode{overlap\_area} covered at least twice, the \inlinecode{perimeter} of the outline including any enclosed hole, and \inlinecode{max\_weight}, the greatest weighted sum at any point. Subtracting the second from the first leaves the area covered exactly once. Each rectangle must satisfy ${\inlinecodemath{x1}} < {\inlinecodemath{x2}}$ and ${\inlinecodemath{y1}} < {\inlinecodemath{y2}}$.
\end{compactitem}

Weights change only the maximum, since the three geometric measures count coverage rather than payload, and weighted area needs no sweep at all: it is $\sum w_i$ times area, one loop over the input. Area covered at least $k$ times instead needs a vector of $k$ lengths per node; for that or any aggregate that does not decompose this way, the compression of section \secref{7.3.8} enumerates every elementary cell at $O(n^{2})$.

Overflow warning: Coordinate differences, areas, perimeters, and weight sums must fit in \inlinecode{int64\_t}.

\textbf{Time Complexity}: $O(n \log n)$ per call, where $n$ is the number of rectangles.

\textbf{Space Complexity}: $O(n)$ auxiliary.

\begin{lstlisting}[style=cpp]
#include <algorithm>
#include <cassert>
#include <cstdint>
#include <cstdlib>
#include <tuple>
#include <utility>
#include <vector>

struct Rect {
  int64_t x1, y1, x2, y2, w = 1;
};

struct UnionMeasure {
  int64_t area = 0, overlap_area = 0, perimeter = 0, max_weight = 0;
};

// The cross section of the sweep: over the currently covered y intervals, the length covered at
// least once and at least twice, the number of maximal covered runs, and the greatest weight at any
// point. Each is its own recurrence, so each is pulled up separately.
class CoverTree {
  static int checked_len(const std::vector<int64_t> &ys) {
    assert(ys.size() >= 2);
    return static_cast<int>(ys.size()) - 1;
  }

  std::vector<int64_t> ys;  // Compressed coordinates; leaf i spans [ys[i], ys[i + 1]).
  int len;
  std::vector<int> count, runs;
  std::vector<int64_t> covered, covered_twice, weight, heaviest;
  std::vector<char> covers_lo, covers_hi;

  // Length covered at least once and at least twice. The second reads the first one level down, so
  // covering at least k times extends this chain by one array per level.
  void pull_coverage(int i, int lo, int hi) {
    int left = 2 * i + 1, right = left + 1;
    if (count[i] > 0) {
      covered[i] = ys[hi + 1] - ys[lo];
      // One cover here turns anything already covered below into a double cover.
      covered_twice[i] =
          count[i] > 1 ? covered[i] : (lo == hi ? 0 : covered[left] + covered[right]);
    } else if (lo == hi) {
      covered[i] = covered_twice[i] = 0;
    } else {
      covered[i] = covered[left] + covered[right];
      covered_twice[i] = covered_twice[left] + covered_twice[right];
    }
  }

  // Number of maximal covered runs, with the flags that keep a run spanning a boundary from being
  // counted in both children.
  void pull_runs(int i, int lo, int hi) {
    int left = 2 * i + 1, right = left + 1;
    if (count[i] > 0) {
      runs[i] = covers_lo[i] = covers_hi[i] = 1;
    } else if (lo == hi) {
      runs[i] = covers_lo[i] = covers_hi[i] = 0;
    } else {
      runs[i] = runs[left] + runs[right] - (covers_hi[left] && covers_lo[right] ? 1 : 0);
      covers_lo[i] = covers_lo[left];
      covers_hi[i] = covers_hi[right];
    }
  }

  // The weight over an interval is the sum of the weights above it, so the heaviest point is the
  // largest root-to-leaf sum. This one ignores the cover count entirely.
  void pull_weight(int i, int lo, int hi) {
    heaviest[i] = weight[i] + (lo == hi ? 0 : std::max(heaviest[2 * i + 1], heaviest[2 * i + 2]));
  }

 public:
  explicit CoverTree(std::vector<int64_t> coords)
      : ys(std::move(coords)),
        len(checked_len(ys)),
        count(4 * len),
        runs(4 * len),
        covered(4 * len),
        covered_twice(4 * len),
        weight(4 * len),
        heaviest(4 * len),
        covers_lo(4 * len),
        covers_hi(4 * len) {}

  // Adds delta covers of weight w to every elementary interval in [lo, hi]; a removal negates both.
  void update(int lo, int hi, int delta, int64_t w) {
    auto rec = [&](auto &&rec, int i, int l, int h) {
      if (hi < l || h < lo) {
        return;
      }
      if (lo <= l && h <= hi) {
        count[i] += delta;
        weight[i] += delta * w;
      } else {
        int mid = l + (h - l) / 2;
        rec(rec, 2 * i + 1, l, mid);
        rec(rec, 2 * i + 2, mid + 1, h);
      }
      pull_coverage(i, l, h);
      pull_runs(i, l, h);
      pull_weight(i, l, h);
    };
    rec(rec, 0, 0, len - 1);
  }

  int64_t covered_length() const { return covered[0]; }
  int64_t twice_covered_length() const { return covered_twice[0]; }
  int run_count() const { return runs[0]; }
  int64_t heaviest_weight() const { return heaviest[0]; }
};

UnionMeasure rect_union(const std::vector<Rect> &rects) {
  if (rects.empty()) {
    return {};
  }
  assert(std::all_of(rects.begin(), rects.end(), [](const Rect &r) {
    return r.x1 < r.x2 && r.y1 < r.y2;
  }));
  std::vector<int64_t> ys;
  for (const Rect &r : rects) {
    ys.push_back(r.y1);
    ys.push_back(r.y2);
  }
  std::sort(ys.begin(), ys.end());
  ys.erase(std::unique(ys.begin(), ys.end()), ys.end());
  auto index = [&](int64_t y) {
    return static_cast<int>(std::lower_bound(ys.begin(), ys.end(), y) - ys.begin());
  };
  std::vector<std::tuple<int64_t, int, int, int, int64_t>> events;  // (x, lo, hi, delta, weight)
  for (const Rect &r : rects) {
    events.emplace_back(r.x1, index(r.y1), index(r.y2) - 1, 1, r.w);
    events.emplace_back(r.x2, index(r.y1), index(r.y2) - 1, -1, r.w);
  }
  std::sort(events.begin(), events.end());
  CoverTree tree(std::move(ys));
  UnionMeasure res;
  int64_t width = 0, doubled = 0, strips = 0, prev_x = 0;
  for (int i = 0, nevents = static_cast<int>(events.size()); i < nevents;) {
    int64_t x = std::get<0>(events[i]);
    if (i > 0) {  // The strip since the previous event has a constant cross-section.
      res.area += width * (x - prev_x);
      res.overlap_area += doubled * (x - prev_x);
      res.perimeter += 2 * strips * (x - prev_x);
    }
    int end = i;
    while (end < nevents && std::get<0>(events[end]) == x) {
      end++;
    }
    // Add before removing at the same x. The sum of the resulting length changes is the symmetric
    // difference of the old and new cross-sections: shared edges vanish, while disjoint entering
    // and leaving intervals are both counted.
    for (int wanted_delta : {1, -1}) {
      int64_t before = tree.covered_length();
      for (int j = i; j < end; j++) {
        auto [_, tgt_lo, tgt_hi, delta, w] = events[j];
        if (delta == wanted_delta) {
          tree.update(tgt_lo, tgt_hi, delta, w);
        }
      }
      res.perimeter += std::llabs(tree.covered_length() - before);
    }
    i = end;
    res.max_weight = std::max(res.max_weight, tree.heaviest_weight());
    width = tree.covered_length();
    doubled = tree.twice_covered_length();
    strips = tree.run_count();
    prev_x = x;
  }
  return res;
}
\end{lstlisting}
\Needspace*{4\baselineskip}
\textbf{Example Usage}\par
\begin{lstlisting}[style=cpp]
#include <cassert>
using namespace std;

int main() {
  assert(rect_union({}).area == 0 && rect_union({}).perimeter == 0);

  UnionMeasure one = rect_union({{0, 0, 2, 3}});
  assert(one.area == 6 && one.perimeter == 10);
  assert(one.overlap_area == 0 && one.max_weight == 1);

  // Disjoint rectangles simply add up, and nothing is covered twice.
  UnionMeasure apart = rect_union({{0, 0, 1, 1}, {2, 0, 3, 1}});
  assert(apart.area == 2 && apart.perimeter == 8);
  assert(apart.overlap_area == 0 && apart.max_weight == 1);

  // Disjoint intervals can leave and enter at the same x; both vertical edges still count.
  UnionMeasure crossing = rect_union({{0, 0, 1, 1}, {1, 2, 2, 3}});
  assert(crossing.area == 2 && crossing.perimeter == 8);

  // Two squares meeting in a unit square, whose union outlines an L-shaped hexagon.
  UnionMeasure pair = rect_union({{0, 0, 2, 2}, {1, 1, 3, 3}});
  assert(pair.area == 7 && pair.perimeter == 12);
  assert(pair.overlap_area == 1 && pair.max_weight == 2);
  assert(pair.area - pair.overlap_area == 6);  // Area belonging to exactly one square.

  // A rectangle nested inside another adds no area, but doubles the cover where it sits.
  UnionMeasure nested = rect_union({{0, 0, 10, 10}, {2, 2, 4, 4}});
  assert(nested.area == 100 && nested.perimeter == 40);
  assert(nested.overlap_area == 4 && nested.max_weight == 2);

  // Four rectangles enclosing a hole, overlapping only at the four corners.
  UnionMeasure ring = rect_union({{0, 0, 3, 1}, {0, 2, 3, 3}, {0, 0, 1, 3}, {2, 0, 3, 3}});
  assert(ring.area == 8);
  assert(ring.perimeter == 16);  // 12 around the outside and 4 around the unit hole.
  assert(ring.overlap_area == 4 && ring.max_weight == 2);

  // Three rectangles stacked over one cell.
  UnionMeasure stack = rect_union({{0, 0, 2, 2}, {1, 0, 3, 2}, {1, 1, 3, 3}});
  assert(stack.max_weight == 3 && stack.overlap_area == 3);
  return 0;
}
\end{lstlisting}
\subsection{Multiple Segment Intersection}
\label{sec:7.3.10}
Given a list of line segments in two dimensions, determine whether any pair of segments intersect using a sweep line algorithm. Endpoint events are processed from left to right while an ordered set maintains the segments currently crossing the sweep line, sorted by $y$; only segments that become adjacent in this set need to be tested, since any leftmost intersection must involve a pair that is adjacent just before it occurs. The input type \inlinecode{Segment\textless{}Pt\textgreater{}} is templated on the point type, so endpoints may be integer (\inlinecode{PointI}) or floating-point (\inlinecode{Point} or \inlinecode{PointD}), but must support \inlinecode{operator\textless{}} which orders points lexicographically. The cross-product sign tests in \inlinecode{seg\_intersection} are exact for integer endpoints, so intersection detection is exact.

This implementation counts endpoint contacts as intersections. See \secref{7.2.3} for pairwise intersection predicates when that distinction matters.

\begin{compactitem}
  \item \inlinecode{find\_intersecting\_pair(segs)} returns the indices of one pair of intersecting segments, or \inlinecode{std::nullopt} if none exist. Constructing a \inlinecode{Segment} places its lexicographically smaller endpoint first, as required by the sweep.
\end{compactitem}

Overflow warning: \inlinecode{seg\_intersection} forms the usual quadratic cross products, but the $y$-ordering cross-multiplication (\inlinecode{ay * bdx}) is cubic in the coordinate magnitude. Use \inlinecode{int64\_t} coordinates for modest integer inputs and a wider intermediate type when these cubic products may exceed it.

\textbf{Time Complexity}: $O(n \log n)$ per call, where $n$ is the number of segments.

\textbf{Space Complexity}: $O(n)$ auxiliary, where $n$ is the number of segments.

\begin{lstlisting}[style=cpp]
#include <algorithm>
#include <optional>
#include <set>
#include <tuple>
#include <utility>
#include <vector>

// Simplified detection-only version of seg_intersection() from 7.2.3 (exact for integral Pt).
template<typename Pt>
bool seg_intersection(const Pt &a, const Pt &b, const Pt &c, const Pt &d) {
  auto cross = [](const Pt &p, const Pt &q, const Pt &r) {
    return (q.x - p.x) * (r.y - p.y) - (q.y - p.y) * (r.x - p.x);
  };
  auto c1 = cross(a, b, c), c2 = cross(a, b, d), c3 = cross(c, d, a), c4 = cross(c, d, b);
  if (c1 == 0 && c2 == 0 && c3 == 0 && c4 == 0) {
    Pt p = std::max(std::min(a, b), std::min(c, d));
    Pt q = std::min(std::max(a, b), std::max(c, d));
    return !(q < p);
  }
  auto straddles = [](auto x, auto y) { return (x <= 0 && 0 <= y) || (y <= 0 && 0 <= x); };
  return straddles(c1, c2) && straddles(c3, c4);
}

// The point type Pt must support operator< (used to canonicalize endpoint order).
template<typename Pt>
struct Segment {
  Pt p, q;
  Segment(const Pt &p, const Pt &q) : p(std::min(p, q)), q(std::max(p, q)) {}
};

template<typename Pt>
bool intersects(const Segment<Pt> &s1, const Segment<Pt> &s2) {
  return seg_intersection(s1.p, s1.q, s2.p, s2.q);
}

template<typename Pt>
std::optional<std::pair<int, int>> find_intersecting_pair(const std::vector<Segment<Pt>> &segs) {
  struct Event {
    Pt p;
    bool is_end;
    int seg;
  };
  std::vector<Event> e;
  for (int i = 0; i < static_cast<int>(segs.size()); i++) {
    e.push_back(Event{segs[i].p, false, i});
    e.push_back(Event{segs[i].q, true, i});
  }
  std::sort(e.begin(), e.end(), [](const Event &a, const Event &b) {
    return std::tie(a.p.x, a.is_end, a.p.y, a.seg) < std::tie(b.p.x, b.is_end, b.p.y, b.seg);
  });
  // Compare y-values at sweep coordinate x without division: y = (p.y*dx + dy*(x - p.x)) / dx.
  // Vertical segments use their lower endpoint.
  auto ycmp = [](const auto &a, const auto &b, auto x) {
    // Overflow risk for integer Pt: the final cross-multiply is ~O(max_coord^3).
    auto adx = a.q.x - a.p.x, bdx = b.q.x - b.p.x;
    auto ay = a.p.y * adx + (a.q.y - a.p.y) * (x - a.p.x);
    auto by = b.p.y * bdx + (b.q.y - b.p.y) * (x - b.p.x);
    if (adx == 0) {
      ay = a.p.y;
      adx = 1;
    }
    if (bdx == 0) {
      by = b.p.y;
      bdx = 1;
    }
    auto lhs = ay * bdx, rhs = by * adx;
    return (lhs < rhs) ? -1 : ((rhs < lhs) ? 1 : 0);
  };
  auto cmp = [&segs, ycmp](int a, int b) {
    if (a == b) {
      return false;
    }
    auto x = std::max(segs[a].p.x, segs[b].p.x);
    int order = ycmp(segs[a], segs[b], x);
    return (order == 0) ? (a < b) : (order < 0);
  };
  using ActiveSet = std::set<int, decltype(cmp)>;
  ActiveSet s(cmp);
  std::vector<typename ActiveSet::iterator> position(segs.size());
  auto result = [](int i, int j) { return std::pair{std::min(i, j), std::max(i, j)}; };
  for (const auto &ev : e) {
    int seg = ev.seg;
    if (!ev.is_end) {
      auto it = s.insert(seg).first;
      position[seg] = it;
      auto next = it;
      if (++next != s.end() && intersects(segs[*next], segs[seg])) {
        return result(*next, seg);
      }
      if (it != s.begin() && intersects(segs[*--it], segs[seg])) {
        return result(*it, seg);
      }
    } else {
      auto it = position[seg];
      auto next = it;
      if (++next != s.end() && it != s.begin()) {
        auto prev = it;
        if (intersects(segs[*next], segs[*--prev])) {
          return result(*next, *prev);
        }
      }
      s.erase(it);
    }
  }
  return std::nullopt;
}
\end{lstlisting}
\Needspace*{4\baselineskip}
\textbf{Example Usage}\par
\begin{lstlisting}[style=cpp]
#include <cassert>
#include <vector>
using namespace std;

struct Point {
  double x, y;
  Point(double x = 0, double y = 0) : x(x), y(y) {}
  bool operator<(const Point &p) const { return x != p.x ? x < p.x : y < p.y; }
};

struct PointI {
  int x, y;
  PointI(int x = 0, int y = 0) : x(x), y(y) {}
  bool operator<(const PointI &p) const { return x != p.x ? x < p.x : y < p.y; }
};

int main() {
  {  // Double-coordinate segments.
    vector<Segment<Point>> v{
        Segment<Point>(Point(0, 0), Point(2, 2)), Segment<Point>(Point(3, 0), Point(0, -1)),
        Segment<Point>(Point(0, 2), Point(2, -2)), Segment<Point>(Point(0, 3), Point(9, 0))
    };
    assert((find_intersecting_pair(v) == pair{0, 2}));
  }
  {  // Integer-coordinate segments: detection is exact.
    vector<Segment<PointI>> v{
        Segment<PointI>({0, 0}, {2, 2}), Segment<PointI>({3, 0}, {0, -1}),
        Segment<PointI>({0, 2}, {2, -2}), Segment<PointI>({0, 3}, {9, 0})
    };
    assert((find_intersecting_pair(v) == pair{0, 2}));

    const vector<Segment<PointI>> disjoint{
        Segment<PointI>({0, 0}, {1, 0}), Segment<PointI>({0, 5}, {1, 5})
    };
    assert(!find_intersecting_pair(disjoint));
    const vector<Segment<PointI>> duplicate{
        Segment<PointI>({0, 0}, {2, 2}), Segment<PointI>({0, 0}, {2, 2})
    };
    assert((find_intersecting_pair(duplicate) == pair{0, 1}));
  }
  {  // Shared endpoints count as intersections.
    vector<Segment<PointI>> shared{
        Segment<PointI>({0, 0}, {2, 2}), Segment<PointI>({2, 2}, {4, 0})
    };
    assert((find_intersecting_pair(shared) == pair{0, 1}));
  }
  return 0;
}
\end{lstlisting}
\subsection{Manhattan MST}
\label{sec:7.3.11}
Given points in the plane, produce a small set of candidate edges guaranteed to contain a minimum spanning tree under Manhattan distance $|x_1 - x_2| + |y_1 - y_2|$. The plane is swept in four rotations/reflections; in each sweep, dominance by $x + y$ identifies the only nearby candidates that can matter. Run Kruskal on the returned edges to obtain the MST.

\begin{compactitem}
  \item \inlinecode{manhattan\_mst\_candidates(p)} returns $O(n)$ candidate edges as tuples (\inlinecode{weight}, \inlinecode{u}, \inlinecode{v}), where $n = {\inlinecodemath{p.size()}}$ and \inlinecode{u} and \inlinecode{v} are indices into \inlinecode{p}.
  \item \inlinecode{manhattan\_mst\_weight(p)} returns the MST weight by running Kruskal's algorithm on those candidate edges.
\end{compactitem}

Overflow warning: the sweeps form coordinate sums, negations, and differences such as \inlinecode{x + y} and \inlinecode{dx + dy}. With \inlinecode{int64\_t} coordinates, those intermediate values must fit in \inlinecode{int64\_t}.

\textbf{Time Complexity}: $O(n \log n)$ per call, where $n$ is the number of points.

\textbf{Space Complexity}: $O(n)$ auxiliary.

\begin{lstlisting}[style=cpp]
#include <algorithm>
#include <cstdint>
#include <cstdlib>
#include <map>
#include <numeric>
#include <tuple>
#include <utility>
#include <vector>

struct Point {
  int64_t x, y;
  Point(int64_t x = 0, int64_t y = 0) : x(x), y(y) {}
};

std::vector<std::tuple<int64_t, int, int>> manhattan_mst_candidates(std::vector<Point> p) {
  int n = static_cast<int>(p.size());
  std::vector<int> id(n);
  std::iota(id.begin(), id.end(), 0);
  std::vector<std::tuple<int64_t, int, int>> edges;
  for (int rot = 0; rot < 4; rot++) {
    std::sort(id.begin(), id.end(), [&](int i, int j) {
      return p[i].x + p[i].y < p[j].x + p[j].y;
    });
    std::map<int64_t, int> sweep;
    for (int i : id) {
      for (auto it = sweep.lower_bound(-p[i].y); it != sweep.end();) {
        int j = it->second;
        int64_t dx = p[i].x - p[j].x, dy = p[i].y - p[j].y;
        if (dy > dx) {
          break;
        }
        edges.emplace_back(dx + dy, i, j);
        it = sweep.erase(it);
      }
      sweep[-p[i].y] = i;
    }
    for (Point &q : p) {
      if (rot & 1) {
        q.x = -q.x;
      } else {
        std::swap(q.x, q.y);
      }
    }
  }
  return edges;
}

struct DSU {
  std::vector<int> root, rank;

  explicit DSU(int n) : root(n), rank(n, 0) { std::iota(root.begin(), root.end(), 0); }
  int find(int u) { return root[u] == u ? u : root[u] = find(root[u]); }

  bool unite(int u, int v) {
    u = find(u);
    v = find(v);
    if (u == v) {
      return false;
    }
    if (rank[u] < rank[v]) {
      std::swap(u, v);
    }
    root[v] = u;
    if (rank[u] == rank[v]) {
      rank[u]++;
    }
    return true;
  }
};

int64_t manhattan_mst_weight(const std::vector<Point> &p) {
  auto edges = manhattan_mst_candidates(p);
  std::sort(edges.begin(), edges.end());
  DSU dsu(static_cast<int>(p.size()));
  int64_t res = 0;
  for (auto [w, u, v] : edges) {
    if (dsu.unite(u, v)) {
      res += w;
    }
  }
  return res;
}
\end{lstlisting}
\Needspace*{4\baselineskip}
\textbf{Example Usage}\par
\begin{lstlisting}[style=cpp]
#include <cassert>
using namespace std;

int main() {
  assert(manhattan_mst_weight({Point(5, -7)}) == 0);
  assert(manhattan_mst_weight({Point(0, 0), Point(3, -4)}) == 7);

  vector<Point> p{{0, 0}, {2, 1}, {3, 4}, {-1, 2}, {5, 0}};
  assert(manhattan_mst_weight(p) == 14);
  auto edges = manhattan_mst_candidates(p);
  assert(edges.size() <= 4 * p.size());
  for (auto [w, u, v] : edges) {
    assert(0 <= u && u < static_cast<int>(p.size()));
    assert(0 <= v && v < static_cast<int>(p.size()));
    assert(w == abs(p[u].x - p[v].x) + abs(p[u].y - p[v].y));
  }
  return 0;
}
\end{lstlisting}

\section{\texorpdfstring{Advanced Planar Geometry}{7.4 Advanced Planar Geometry}}
\setcounter{section}{4}
\setcounter{subsection}{0}
\subsection{Convex Polygon Cut}
\label{sec:7.4.1}
Given a convex polygon and a directed line from \inlinecode{p} to \inlinecode{q}, clips the polygon to the closed left half-plane of that line. The polygon is walked edge by edge: vertices on the left side of the line are kept, and whenever an edge crosses the line, the intersection point is appended to the output.

\begin{compactitem}
  \item \inlinecode{convex\_cut(lo, hi, p, q)} returns the portion of the polygon lying on or to the left of the directed line from \inlinecode{p} to \inlinecode{q}. The input range $[{\inlinecodemath{lo}}, {\inlinecodemath{hi}})$ must contain the vertices of a convex polygon in boundary order, either clockwise or counterclockwise. The returned polygon preserves that boundary order and uses \inlinecode{Point} with \inlinecode{double} coordinates because edge-line intersections are computed in floating point. The points \inlinecode{p} and \inlinecode{q} must differ.
\end{compactitem}

Overflow warning: For integer-coordinate inputs, classification is exact only if the intermediate products do not overflow.

\textbf{Time Complexity}: $O(n)$ per call, where $n$ is the distance between \inlinecode{lo} and \inlinecode{hi}.

\textbf{Space Complexity}: $O(n)$ for the returned polygon and $O(1)$ auxiliary.

\begin{lstlisting}[style=cpp]
#include <cassert>
#include <cmath>
#include <type_traits>
#include <vector>

const double EPS = 1e-9;

template<typename T, typename U, typename C = std::common_type_t<T, U>>
bool EQ(T a, U b) {
  if constexpr (std::is_floating_point_v<C>) return C(a) == C(b) || std::fabs(C(a) - C(b)) <= EPS;
  return C(a) == C(b);
}

template<typename T, typename U, typename C = std::common_type_t<T, U>>
bool LT(T a, U b) {
  if constexpr (std::is_floating_point_v<C>) return C(a) < C(b) - EPS;
  return C(a) < C(b);
}

struct Point {
  double x, y;
  Point(double x = 0, double y = 0) : x(x), y(y) {}
  bool operator==(const Point &p) const { return x == p.x && y == p.y; }
  friend bool EQ(const Point &a, const Point &b) { return EQ(a.x, b.x) && EQ(a.y, b.y); }
};

template<typename PtA, typename PtB, typename PtO>
auto cross(const PtA &a, const PtB &b, const PtO &o) {
  return (a.x - o.x) * (b.y - o.y) - (a.y - o.y) * (b.x - o.x);
}

template<typename PtA, typename PtO, typename PtB>
int turn(const PtA &a, const PtO &o, const PtB &b) {
  auto c = cross(a, b, o);
  return LT(c, 0) ? 1 : (LT(0, c) ? -1 : 0);
}

template<typename PtA, typename PtB>
int line_intersection(
    const PtA &p1, const PtA &p2, const PtB &p3, const PtB &p4, Point *p = nullptr
) {
  double a1 = static_cast<double>(p2.y) - p1.y, b1 = static_cast<double>(p1.x) - p2.x;
  double c1 = -(static_cast<double>(p1.x) * p2.y - static_cast<double>(p2.x) * p1.y);
  double a2 = static_cast<double>(p4.y) - p3.y, b2 = static_cast<double>(p3.x) - p4.x;
  double c2 = -(static_cast<double>(p3.x) * p4.y - static_cast<double>(p4.x) * p3.y);
  double x = -(c1 * b2 - c2 * b1), y = -(a1 * c2 - a2 * c1);
  double det = a1 * b2 - a2 * b1;
  if (EQ(det, 0)) {
    return (EQ(x, 0) && EQ(y, 0)) ? 1 : -1;
  }
  if (p != nullptr) {
    *p = Point(x / det, y / det);
  }
  return 0;
}

template<typename It, typename Pt>
std::vector<Point> convex_cut(It lo, It hi, const Pt &p, const Pt &q) {
  assert(!EQ(p.x, q.x) || !EQ(p.y, q.y));
  if (lo == hi) {
    return {};
  }
  std::vector<Point> res;
  for (It i = lo, j = hi - 1; i != hi; j = i++) {
    Point pj(static_cast<double>(j->x), static_cast<double>(j->y));
    Point pi(static_cast<double>(i->x), static_cast<double>(i->y));
    int d1 = turn(q, p, *j), d2 = turn(q, p, *i);
    if (d1 <= 0) {
      res.push_back(pj);
    }
    if (d1 * d2 < 0) {
      Point r;
      line_intersection(p, q, *j, *i, &r);
      res.push_back(r);
    }
  }
  return res;
}
\end{lstlisting}
\Needspace*{4\baselineskip}
\textbf{Example Usage}\par
\begin{lstlisting}[style=cpp]
using namespace std;

struct PointI {
  int x, y;
  PointI(int x = 0, int y = 0) : x(x), y(y) {}
};

bool EQ(const vector<Point> &a, const vector<Point> &b) {
  return a.size() == b.size() &&
         equal(a.begin(), a.end(), b.begin(), [](const Point &p, const Point &q) {
           return EQ(p, q);
         });
}

int main() {
  {
    vector<Point> v{{1, 3}, {2, 2}, {2, 1}, {0, 0}, {-1, 3}};
    // Cut using the vertical line through (0, 0).
    vector<Point> c{{-1, 3}, {0, 3}, {0, 0}};
    assert(EQ(convex_cut(v.begin(), v.end(), Point(0, 0), Point(0, 1)), c));
  }
  {  // On a non-convex input, the result may be multiple disjoint polygons!
    vector<Point> v{{0, 0}, {2, 2}, {0, 4}, {3, 4}, {3, 0}};
    vector<Point> c{{1, 0}, {0, 0}, {1, 1}, {1, 3}, {0, 4}, {1, 4}};
    assert(EQ(convex_cut(v.begin(), v.end(), Point(1, 0), Point(1, 4)), c));
  }
  {
    vector<PointI> v{{0, 0}, {4, 0}, {4, 4}, {0, 4}};
    vector<Point> c{{0, 4}, {0, 0}, {2, 0}, {2, 4}};
    assert(EQ(convex_cut(v.begin(), v.end(), PointI(2, 0), PointI(2, 4)), c));
  }
  return 0;
}
\end{lstlisting}
\subsection{Half-Plane Intersection}
\label{sec:7.4.2}
Given a list of directed lines, compute the convex polygon contained in every left half-plane. Half-plane intersection sorts the lines by direction, keeps only the tightest representative among parallel lines, and maintains a deque of lines whose pairwise intersections form the current boundary. Directions are ordered by half-plane and cross-product signs, without epsilon comparisons, to satisfy \inlinecode{std::sort}'s strict ordering requirement. \inlinecode{EPS} is then used outside the comparator to group parallel directions and retain only the tightest half-plane in each group.

\begin{compactitem}
  \item \inlinecode{half\_plane\_intersection(planes)} returns the polygon cut out by the half-planes in counter-clockwise order. Each half-plane is represented by \inlinecode{HalfPlane(p, q)}, meaning the closed region to the left of the directed line \inlinecode{p} $\to$ \inlinecode{q}. This implementation is intended for intersections that form a bounded polygon with positive area; empty, unbounded, or degenerate inputs may return an empty vector. Add explicit bounding-box half-planes when a bounded polygon is required. The two points defining each half-plane must differ.
  \item \inlinecode{out(q)} returns whether point \inlinecode{q} lies strictly outside this half-plane, that is, to the right of its directed line.
  \item \inlinecode{intersect(h)} returns the point where the boundary lines of this half-plane and \inlinecode{h} meet, which must not be parallel.
\end{compactitem}

\textbf{Time Complexity}: $O(n \log n)$ per call, where $n$ is the number of half-planes.

\textbf{Space Complexity}: $O(n)$ auxiliary.

\begin{lstlisting}[style=cpp]
#include <algorithm>
#include <cmath>
#include <deque>
#include <type_traits>
#include <vector>

const double EPS = 1e-9;

template<typename T, typename U, typename C = std::common_type_t<T, U>>
bool EQ(T a, U b) {
  if constexpr (std::is_floating_point_v<C>) return C(a) == C(b) || std::fabs(C(a) - C(b)) <= EPS;
  return C(a) == C(b);
}

template<typename T, typename U, typename C = std::common_type_t<T, U>>
bool LT(T a, U b) {
  if constexpr (std::is_floating_point_v<C>) return C(a) < C(b) - EPS;
  return C(a) < C(b);
}

struct Point {
  double x, y;
  Point(double x = 0, double y = 0) : x(x), y(y) {}
  bool operator==(const Point &p) const { return x == p.x && y == p.y; }
  friend bool EQ(const Point &a, const Point &b) { return EQ(a.x, b.x) && EQ(a.y, b.y); }
  Point operator+(const Point &p) const { return {x + p.x, y + p.y}; }
  Point operator-(const Point &p) const { return {x - p.x, y - p.y}; }
  Point operator*(double k) const { return {x * k, y * k}; }
  double cross(const Point &p) const { return x * p.y - y * p.x; }
};

struct HalfPlane {
  Point p, dir;

  HalfPlane() = default;
  HalfPlane(Point p, Point q) : p(p), dir(q - p) {}

  bool out(const Point &q) const { return LT(dir.cross(q - p), 0); }

  Point intersect(const HalfPlane &h) const {
    double t = h.dir.cross(p - h.p) / dir.cross(h.dir);
    return p + dir * t;
  }

  bool operator<(const HalfPlane &h) const {
    auto half = [](const Point &d) { return d.y < 0 || (d.y == 0 && d.x < 0); };
    if (half(dir) != half(h.dir)) {
      return half(dir) < half(h.dir);
    }
    return dir.cross(h.dir) > 0;
  }
};

std::vector<Point> half_plane_intersection(std::vector<HalfPlane> planes) {
  std::sort(planes.begin(), planes.end());
  std::vector<HalfPlane> unique;
  for (const HalfPlane &h : planes) {
    if (!unique.empty() && EQ(unique.back().dir.cross(h.dir), 0)) {
      if (unique.back().dir.cross(h.p - unique.back().p) > 0) {
        unique.back() = h;
      }
      continue;
    }
    unique.push_back(h);
  }
  std::deque<HalfPlane> dq;
  auto bad_back = [&](const HalfPlane &h) {
    return dq.size() >= 2 && h.out(dq[dq.size() - 2].intersect(dq.back()));
  };
  auto bad_front = [&](const HalfPlane &h) {
    return dq.size() >= 2 && h.out(dq[0].intersect(dq[1]));
  };
  for (const HalfPlane &h : unique) {
    while (bad_back(h)) {
      dq.pop_back();
    }
    while (bad_front(h)) {
      dq.pop_front();
    }
    dq.push_back(h);
  }
  while (dq.size() >= 3 && dq.front().out(dq[dq.size() - 2].intersect(dq.back()))) {
    dq.pop_back();
  }
  while (dq.size() >= 3 && dq.back().out(dq[0].intersect(dq[1]))) {
    dq.pop_front();
  }
  if (dq.size() < 3) {
    return {};
  }
  std::vector<Point> res;
  for (int i = 0; i < static_cast<int>(dq.size()); i++) {
    if (EQ(dq[i].dir.cross(dq[(i + 1) % dq.size()].dir), 0)) {
      return {};
    }
    res.push_back(dq[i].intersect(dq[(i + 1) % dq.size()]));
  }
  return res;
}
\end{lstlisting}
\Needspace*{4\baselineskip}
\textbf{Example Usage}\par
\begin{lstlisting}[style=cpp]
#include <cassert>
using namespace std;

double polygon_area(const vector<Point> &p) {
  double area = 0;
  for (int i = 0, j = static_cast<int>(p.size()) - 1; i < static_cast<int>(p.size()); j = i++) {
    area += p[j].cross(p[i]);
  }
  return fabs(area) / 2.0;
}

int main() {
  vector<HalfPlane> box{
      HalfPlane(Point(0, 0), Point(4, 0)),
      HalfPlane(Point(4, 0), Point(4, 4)),
      HalfPlane(Point(4, 4), Point(0, 4)),
      HalfPlane(Point(0, 4), Point(0, 0)),
  };
  auto square = half_plane_intersection(box);
  assert(square.size() == 4 && EQ(polygon_area(square), 16));

  auto tighter = box;
  tighter.emplace_back(Point(0, 1), Point(4, 1));  // y >= 1
  assert(EQ(polygon_area(half_plane_intersection(tighter)), 12));

  box.emplace_back(Point(2, 0), Point(2, 4));  // x <= 2
  auto rect = half_plane_intersection(box);
  assert(rect.size() == 4 && EQ(polygon_area(rect), 8));

  vector<HalfPlane> empty{
      HalfPlane(Point(0, 0), Point(1, 0)),  // y >= 0
      HalfPlane(Point(1, 1), Point(0, 1)),  // y <= 1
      HalfPlane(Point(0, 2), Point(1, 2)),  // y >= 2
  };
  assert(half_plane_intersection(empty).empty());
  return 0;
}
\end{lstlisting}
\subsection{Minkowski Sum}
\label{sec:7.4.3}
Given two convex polygons, compute their Minkowski sum: the set of all points $a + b$ where $a$ comes from the first polygon and $b$ comes from the second. After rotating both polygons to start at their lowest-leftmost vertex, the boundary edges of the sum are obtained by merging the two cyclic edge-angle lists. Minkowski sums are useful for collision/translation problems, convex polygon distance checks, and diameter computations via $P + (-P)$.

\begin{compactitem}
  \item \inlinecode{minkowski\_sum(a, b)} returns the convex polygon representing the sum. Inputs must be convex polygons in boundary order, clockwise or counter-clockwise, with no repeated final vertex. Collinear consecutive output vertices are removed.
\end{compactitem}

\textbf{Time Complexity}: $O(n + m)$ per call, where $n$ and $m$ are the polygon sizes.

\textbf{Space Complexity}: $O(n + m)$ auxiliary.

\begin{lstlisting}[style=cpp]
#include <algorithm>
#include <cmath>
#include <type_traits>
#include <utility>
#include <vector>

const double EPS = 1e-9;

template<typename T, typename U, typename C = std::common_type_t<T, U>>
bool EQ(T a, U b) {
  if constexpr (std::is_floating_point_v<C>) return C(a) == C(b) || std::fabs(C(a) - C(b)) <= EPS;
  return C(a) == C(b);
}

struct Point {
  double x, y;
  Point(double x = 0, double y = 0) : x(x), y(y) {}
  bool operator==(const Point &p) const { return x == p.x && y == p.y; }
  friend bool EQ(const Point &a, const Point &b) { return EQ(a.x, b.x) && EQ(a.y, b.y); }
  Point operator+(const Point &p) const { return {x + p.x, y + p.y}; }
  Point operator-(const Point &p) const { return {x - p.x, y - p.y}; }
  double cross(const Point &p) const { return x * p.y - y * p.x; }
};

std::vector<Point> normalize(std::vector<Point> p) {
  while (p.size() > 1 && EQ(p.front(), p.back())) {
    p.pop_back();
  }
  if (p.size() <= 1) {
    return p;
  }
  double area = 0;
  for (int i = 0, j = static_cast<int>(p.size()) - 1; i < static_cast<int>(p.size()); j = i++) {
    area += p[j].cross(p[i]);
  }
  if (area < 0) {
    std::reverse(p.begin(), p.end());
  }
  int start = 0;
  for (int i = 1; i < static_cast<int>(p.size()); i++) {
    if (p[i].y < p[start].y || (EQ(p[i].y, p[start].y) && p[i].x < p[start].x)) {
      start = i;
    }
  }
  std::rotate(p.begin(), p.begin() + start, p.end());
  return p;
}

std::vector<Point> remove_collinear(const std::vector<Point> &p) {
  std::vector<Point> res;
  for (const Point &q : p) {
    while (res.size() >= 2 && EQ((res.back() - res[res.size() - 2]).cross(q - res.back()), 0)) {
      res.pop_back();
    }
    if (res.empty() || !EQ(q, res.back())) {
      res.push_back(q);
    }
  }
  while (res.size() >= 3 && EQ((res.back() - res[res.size() - 2]).cross(res[0] - res.back()), 0)) {
    res.pop_back();
  }
  while (res.size() >= 3 && EQ((res[0] - res.back()).cross(res[1] - res[0]), 0)) {
    res.erase(res.begin());
  }
  return res;
}

std::vector<Point> minkowski_sum(std::vector<Point> a, std::vector<Point> b) {
  a = normalize(std::move(a));
  b = normalize(std::move(b));
  if (a.empty() || b.empty()) {
    return {};
  }
  if (a.size() == 1) {
    for (Point &p : b) {
      p = p + a[0];
    }
    return b;
  }
  if (b.size() == 1) {
    for (Point &p : a) {
      p = p + b[0];
    }
    return a;
  }
  int n = static_cast<int>(a.size()), m = static_cast<int>(b.size());
  a.push_back(a[0]);
  a.push_back(a[1]);
  b.push_back(b[0]);
  b.push_back(b[1]);
  std::vector<Point> res;
  int i = 0, j = 0;
  while (i < n || j < m) {
    res.push_back(a[i] + b[j]);
    if (i == n) {
      j++;
      continue;
    }
    if (j == m) {
      i++;
      continue;
    }
    double cr = (a[i + 1] - a[i]).cross(b[j + 1] - b[j]);
    if (cr > -EPS) {
      i++;
    }
    if (cr < EPS) {
      j++;
    }
  }
  return remove_collinear(res);
}
\end{lstlisting}
\Needspace*{4\baselineskip}
\textbf{Example Usage}\par
\begin{lstlisting}[style=cpp]
#include <cassert>
using namespace std;

int main() {
  vector<Point> empty;
  assert(minkowski_sum(empty, {Point(1, 2)}).empty());

  vector<Point> square{{0, 0}, {1, 0}, {1, 1}, {0, 1}};
  vector<Point> tri{{0, 0}, {2, 0}, {0, 2}};
  // The square thickens the triangle by one unit in the positive x/y directions.
  assert((minkowski_sum(square, tri) == vector<Point>{{0, 0}, {3, 0}, {3, 1}, {1, 3}, {0, 3}}));

  // Adding a single point translates the polygon.
  assert((minkowski_sum(square, {Point(2, 3)}) == vector<Point>{{2, 3}, {3, 3}, {3, 4}, {2, 4}}));
  assert((minkowski_sum({Point(1, 2)}, {Point(3, 4)}) == vector<Point>{{4, 6}}));
  return 0;
}
\end{lstlisting}
\subsection{Polygon Intersection and Union}
\label{sec:7.4.4}
Given two simple (non-self-intersecting) polygons, determine the areas of their intersection and union using a sweep line algorithm and the inclusion-exclusion principle. Every vertex and every pairwise edge intersection contributes its $x$-coordinate as a sweep boundary, so within each vertical slab between consecutive boundaries the two borders cannot cross, and the slab's overlap is accumulated exactly from trapezoid areas. The union then follows by inclusion-exclusion as the sum of the individual areas minus the intersection.

\begin{compactitem}
  \item \inlinecode{intersection\_area(a, b)} returns the intersection area of polygons \inlinecode{a} and \inlinecode{b}, whose vertices are listed in boundary order.
  \item \inlinecode{union\_area(a, b)} returns the union area of polygons \inlinecode{a} and \inlinecode{b}.
\end{compactitem}

Overflow warning: the segment-intersection tests form cross products in the input point's coordinate type, which grow like the squared coordinate magnitude. For integral point types, use 64-bit coordinates (e.g. \inlinecode{PointL} from \secref{7.1.1}) if necessary.

\textbf{Time Complexity}: $O(n \cdot m \cdot (n + m) \cdot \log\left(n + m\right))$ per call to \inlinecode{intersection\_area()} and \inlinecode{union\_area()}, where $n$ is the number of vertices in the first polygon and $m$ is the number of vertices in the second polygon (or $O(N^{3} \log N)$ for $N = n + m$).

\textbf{Space Complexity}: $O(n \cdot m)$ auxiliary for \inlinecode{intersection\_area()} and \inlinecode{union\_area()}, where $n$ and $m$ are the respective polygon sizes.

\begin{lstlisting}[style=cpp]
#include <algorithm>
#include <cmath>
#include <type_traits>
#include <utility>
#include <vector>

const double EPS = 1e-9;

template<typename T, typename U, typename C = std::common_type_t<T, U>>
bool EQ(T a, U b) {
  if constexpr (std::is_floating_point_v<C>) return C(a) == C(b) || std::fabs(C(a) - C(b)) <= EPS;
  return C(a) == C(b);
}

template<typename T, typename U, typename C = std::common_type_t<T, U>>
bool LT(T a, U b) {
  if constexpr (std::is_floating_point_v<C>) return C(a) < C(b) - EPS;
  return C(a) < C(b);
}

template<typename T, typename U>
bool LE(T a, U b) {
  return !LT(b, a);
}

template<typename Pt>
bool point_on_segment(const Pt &p, const Pt &a, const Pt &b) {
  // Overflow risk for integer Pt: these products are ~O(max_coord^2); use int64_t if necessary.
  return EQ((p.x - a.x) * (b.y - a.y) - (p.y - a.y) * (b.x - a.x), 0) &&
         LE(std::min(a.x, b.x), p.x) && LE(p.x, std::max(a.x, b.x)) &&
         LE(std::min(a.y, b.y), p.y) && LE(p.y, std::max(a.y, b.y));
}

// Specialized version of seg_intersection() from 7.2.3, simplified for include_boundary = true,
// returning -1 for no intersection, 0 for one intersection point, 1 for positive length overlap.
template<typename Pt>
int seg_intersection1(
    const Pt &a, const Pt &b, const Pt &c, const Pt &d, double *outx, double *outy
) {
  auto ab_x = b.x - a.x, ab_y = b.y - a.y;
  auto ac_x = c.x - a.x, ac_y = c.y - a.y;
  auto cd_x = d.x - c.x, cd_y = d.y - c.y;
  auto ab2 = ab_x * ab_x + ab_y * ab_y;
  auto cd2 = cd_x * cd_x + cd_y * cd_y;
  if (EQ(ab2, 0)) {
    if (point_on_segment(a, c, d)) {
      *outx = a.x;
      *outy = a.y;
      return 0;
    }
    return -1;
  }
  if (EQ(cd2, 0)) {
    if (point_on_segment(c, a, b)) {
      *outx = c.x;
      *outy = c.y;
      return 0;
    }
    return -1;
  }
  auto c1 = ab_x * cd_y - ab_y * cd_x;
  auto c2 = ac_x * ab_y - ac_y * ab_x;
  if (EQ(c1, 0) && EQ(c2, 0)) {  // Collinear.
    Pt res1 = std::max(std::min(a, b), std::min(c, d));
    Pt res2 = std::min(std::max(a, b), std::max(c, d));
    if (!(res2 < res1)) {
      if (res1 == res2) {
        *outx = static_cast<double>(res1.x);
        *outy = static_cast<double>(res1.y);
        return 0;
      }
      return 1;
    }
    return -1;
  }
  if (EQ(c1, 0)) {
    return -1;
  }
  auto t_num = ac_x * cd_y - ac_y * cd_x;
  bool c1_pos = c1 > 0;
  bool t_ok = c1_pos ? (LE(0, t_num) && LE(t_num, c1)) : (LE(c1, t_num) && LE(t_num, 0));
  bool u_ok = c1_pos ? (LE(0, c2) && LE(c2, c1)) : (LE(c1, c2) && LE(c2, 0));
  if (t_ok && u_ok) {
    double t = static_cast<double>(t_num) / c1;
    *outx = static_cast<double>(a.x + t * ab_x);
    *outy = static_cast<double>(a.y + t * ab_y);
    return 0;
  }
  return -1;
}

// The two segments may use different point types (e.g. polygon edge vs. sweep line).
template<typename PtA, typename PtB>
int line_intersection1(
    const PtA &p1, const PtA &p2, const PtB &p3, const PtB &p4, double *outx, double *outy
) {
  double a1 = static_cast<double>(p2.y) - p1.y, b1 = static_cast<double>(p1.x) - p2.x;
  double c1 = -(a1 * p1.x + b1 * p1.y);
  double a2 = static_cast<double>(p4.y) - p3.y, b2 = static_cast<double>(p3.x) - p4.x;
  double c2 = -(a2 * p3.x + b2 * p3.y);
  double x = -(c1 * b2 - c2 * b1), y = -(a1 * c2 - a2 * c1);
  double det = a1 * b2 - a2 * b1;
  if (EQ(det, 0)) {
    return (EQ(x, 0) && EQ(y, 0)) ? 1 : -1;
  }
  if (outx != nullptr && outy != nullptr) {
    *outx = x / det;
    *outy = y / det;
  }
  return 0;
}

template<typename Pt>
double intersection_area(const std::vector<Pt> &a, const std::vector<Pt> &b) {
  if (a.empty() || b.empty()) {
    return 0;
  }
  struct SweepPoint {
    double x, y;
  };  // For line intersection with the sweep line.
  const std::vector<Pt> *polys[2] = {&a, &b};
  int orientation[2];
  for (int id = 0; id < 2; id++) {
    const auto &p = *polys[id];
    int n = static_cast<int>(p.size());
    double area = 0;
    for (int i = 0, j = n - 1; i < n; j = i++) {
      area += (static_cast<double>(p[j].x) - p[i].x) * (static_cast<double>(p[j].y) + p[i].y);
    }
    orientation[id] = (area < 0 ? 1 : (area > 0 ? -1 : 0));
  }
  std::vector<double> x_coords;
  x_coords.reserve(a.size() + b.size());
  for (const Pt &p : a) {
    x_coords.push_back(p.x);
  }
  for (const Pt &p : b) {
    x_coords.push_back(p.x);
  }
  int n = static_cast<int>(a.size()), m = static_cast<int>(b.size());
  for (int i1 = 0, j1 = n - 1; i1 < n; j1 = i1++) {
    for (int i2 = 0, j2 = m - 1; i2 < m; j2 = i2++) {
      double outx, outy;
      if (seg_intersection1(a[i1], a[j1], b[i2], b[j2], &outx, &outy) == 0) {
        x_coords.push_back(outx);
      }
    }
  }
  std::sort(x_coords.begin(), x_coords.end());
  x_coords.erase(
      std::unique(x_coords.begin(), x_coords.end(), [](double a, double b) { return EQ(a, b); }),
      x_coords.end()
  );
  double res = 0;
  for (int k = 0; k + 1 < static_cast<int>(x_coords.size()); k++) {
    double x = (x_coords[k] + x_coords[k + 1]) / 2;
    SweepPoint sweep0{x, 0}, sweep1{x, 1};
    std::vector<std::pair<double, int>> events;  // (y, mask delta)
    for (int id = 0; id < 2; id++) {
      const auto &p = *polys[id];
      int size = static_cast<int>(p.size());
      for (int j = 0, i = size - 1; j < size; i = j++) {
        double px, py;
        if (line_intersection1(p[j], p[i], sweep0, sweep1, &px, &py) == 0) {
          double y = py, x0 = p[i].x, x1 = p[j].x;
          if (x0 < x && x1 > x) {
            events.emplace_back(y, orientation[id] * (1 << id));
          } else if (x0 > x && x1 < x) {
            events.emplace_back(y, -orientation[id] * (1 << id));
          }
        }
      }
    }
    std::sort(events.begin(), events.end());
    double height = 0;
    int coverage[2] = {0, 0};
    for (int j = 0; j < static_cast<int>(events.size()); j++) {
      auto [y, mask_delta] = events[j];
      if (coverage[0] != 0 && coverage[1] != 0) {
        height += y - events[j - 1].first;
      }
      int bit = std::abs(mask_delta);
      int poly = (bit == 1 ? 0 : 1);
      coverage[poly] += mask_delta / bit;
    }
    res += height * (x_coords[k + 1] - x_coords[k]);
  }
  return res;
}

template<typename Pt>
double polygon_area(const std::vector<Pt> &p) {
  if (p.empty()) {
    return 0;
  }
  int n = static_cast<int>(p.size());
  double area = 0;
  for (int i = 0, j = n - 1; i < n; j = i++) {
    area += (static_cast<double>(p[j].x) - p[i].x) * (static_cast<double>(p[j].y) + p[i].y);
  }
  return fabs(area / 2.0);
}

template<typename Pt>
double union_area(const std::vector<Pt> &a, const std::vector<Pt> &b) {
  return polygon_area(a) + polygon_area(b) - intersection_area(a, b);
}
\end{lstlisting}
\Needspace*{4\baselineskip}
\textbf{Example Usage}\par
\begin{lstlisting}[style=cpp]
#include <cassert>
using namespace std;

struct Point {
  double x, y;
  Point(double x = 0, double y = 0) : x(x), y(y) {}
  bool operator==(const Point &p) const { return x == p.x && y == p.y; }
  bool operator!=(const Point &p) const { return !(*this == p); }
  bool operator<(const Point &p) const { return x != p.x ? x < p.x : y < p.y; }
};

struct PointI {
  int x, y;
  PointI(int x = 0, int y = 0) : x(x), y(y) {}
  bool operator==(const PointI &o) const { return x == o.x && y == o.y; }
  bool operator!=(const PointI &p) const { return !(*this == p); }
  bool operator<(const PointI &p) const { return x != p.x ? x < p.x : y < p.y; }
};

int main() {
  vector<Point> p, s;
  // Irregular pentagon overlapping the square below in a triangle of area 1.5.
  p.emplace_back(1, 3);
  p.emplace_back(1, 2);
  p.emplace_back(2, 1);
  p.emplace_back(0, 0);
  p.emplace_back(-1, 3);
  // Square of area 12.5 in quadrant 2.
  s.emplace_back(0, 0);
  s.emplace_back(0, 3);
  s.emplace_back(-3, 3);
  s.emplace_back(-3, 0);
  assert(EQ(1.5, intersection_area(p, s)));
  assert(EQ(12.5, union_area(p, s)));

  // Integer-coordinate polygons are accepted (computation proceeds in double).
  vector<PointI> ip{{1, 3}, {1, 2}, {2, 1}, {0, 0}, {-1, 3}};
  vector<PointI> is{{0, 0}, {0, 3}, {-3, 3}, {-3, 0}};
  assert(EQ(1.5, intersection_area(ip, is)));
  assert(EQ(12.5, union_area(ip, is)));
  return 0;
}
\end{lstlisting}
\subsection{Delaunay Triangulation (Simple)}
\label{sec:7.4.5}
Given a set $P$ of distinct two-dimensional points, a Delaunay triangulation of $P$ is a triangulation of the convex hull of $P$ such that no point of $P$ lies strictly inside the circumcircle of any triangle in the triangulation. The Delaunay triangulation is not necessarily unique when four or more points are cocircular.

This implementation produces one valid triangulation using a simple brute-force algorithm. Each candidate triangle is tested against the empty-circumcircle condition, and candidates whose edges would properly cross already accepted triangles are rejected.

All arithmetic uses the point's own coordinate type, so integer inputs yield an exact triangulation.

\begin{compactitem}
  \item \inlinecode{delaunay\_triangulation(p)} returns a Delaunay triangulation of the distinct points in \inlinecode{p} as a vector of clockwise-oriented vertex triples, or an empty vector if no triangulation exists (e.g. fewer than three points, or all points collinear).
\end{compactitem}

Overflow warning: the lifted-paraboloid test grows like the fourth power of the coordinate magnitude. With 64-bit integer coordinates, it overflows once coordinates exceed roughly $17600$. For floating-point coordinates the test is subject to the usual rounding error.

\textbf{Time Complexity}: $O(n^{6})$ per call, where $n$ is the number of points. There are $O(n^{3})$ candidate triangles; each scans all points and, conservatively, up to $O(n^{3})$ previously accepted candidates.

\textbf{Space Complexity}: $O(n)$ auxiliary, excluding the returned triangulation.

\begin{lstlisting}[style=cpp]
#include <cmath>
#include <tuple>
#include <type_traits>
#include <vector>

const double EPS = 1e-9;

template<typename T, typename U, typename C = std::common_type_t<T, U>>
bool EQ(T a, U b) {
  if constexpr (std::is_floating_point_v<C>) return C(a) == C(b) || std::fabs(C(a) - C(b)) <= EPS;
  return C(a) == C(b);
}

template<typename T, typename U, typename C = std::common_type_t<T, U>>
bool LT(T a, U b) {
  if constexpr (std::is_floating_point_v<C>) return C(a) < C(b) - EPS;
  return C(a) < C(b);
}

template<typename T, typename U>
bool LE(T a, U b) {
  return !LT(b, a);
}

// Specialized version of seg_intersection() from 7.2.3, simplified for include_boundary = false,
// i.e. we're detecting proper crossings of two non-parallel segments whose interiors intersect.
template<typename Pt>
bool proper_seg_intersection(const Pt &a, const Pt &b, const Pt &c, const Pt &d) {
  auto cross = [](const Pt &o, const Pt &p, const Pt &q) {
    return (p.x - o.x) * (q.y - o.y) - (p.y - o.y) * (q.x - o.x);  // Overflow warning.
  };
  auto c1 = cross(a, b, c), c2 = cross(a, b, d), c3 = cross(c, d, a), c4 = cross(c, d, b);
  return ((LT(c1, 0) && LT(0, c2)) || (LT(c2, 0) && LT(0, c1))) &&
         ((LT(c3, 0) && LT(0, c4)) || (LT(c4, 0) && LT(0, c3)));
}

template<typename Pt>
std::vector<std::tuple<Pt, Pt, Pt>> delaunay_triangulation(const std::vector<Pt> &p) {
  int n = static_cast<int>(p.size());
  std::vector<decltype(Pt::x)> z;
  z.reserve(n);
  for (const auto &[px, py] : p) {
    z.emplace_back(px * px + py * py);  // Overflow warning.
  }
  std::vector<std::tuple<Pt, Pt, Pt>> res;
  std::vector<std::tuple<int, int, int>> res_idx;
  for (int i = 0; i < n - 2; i++) {
    for (int j = i + 1; j < n; j++) {
      for (int k = i + 1; k < n; k++) {
        if (j == k) {
          continue;
        }
        // Overflow warning.
        auto nx = (p[j].y - p[i].y) * (z[k] - z[i]) - (p[k].y - p[i].y) * (z[j] - z[i]);
        auto ny = (p[k].x - p[i].x) * (z[j] - z[i]) - (p[j].x - p[i].x) * (z[k] - z[i]);
        auto nz = (p[j].x - p[i].x) * (p[k].y - p[i].y) - (p[k].x - p[i].x) * (p[j].y - p[i].y);
        if (LE(0, nz)) {
          continue;
        }
        std::vector<Pt> s1{p[i], p[j], p[k], p[i]};
        for (int m = 0; m < n; m++) {
          if (LT(0, nx * (p[m].x - p[i].x) + ny * (p[m].y - p[i].y) + nz * (z[m] - z[i]))) {
            goto skip;
          }
        }
        // Reject candidates whose edges properly cross already accepted triangles.
        for (auto [t0, t1, t2] : res_idx) {
          std::vector<Pt> s2{p[t0], p[t1], p[t2], p[t0]};
          for (int u = 0; u < 3; u++) {
            for (int v = 0; v < 3; v++) {
              if (proper_seg_intersection(s1[u], s1[u + 1], s2[v], s2[v + 1])) {
                goto skip;
              }
            }
          }
        }
        res.emplace_back(p[i], p[j], p[k]);
        res_idx.emplace_back(i, j, k);
      skip:;
      }
    }
  }
  return res;
}
\end{lstlisting}
\Needspace*{4\baselineskip}
\textbf{Example Usage}\par
\begin{lstlisting}[style=cpp]
#include <cassert>
using namespace std;

struct Point {
  double x, y;
  Point(double x = 0, double y = 0) : x(x), y(y) {}
  bool operator==(const Point &p) const { return x == p.x && y == p.y; }
  bool operator<(const Point &p) const { return x != p.x ? x < p.x : y < p.y; }
};

struct PointI {
  int x, y;
  PointI(int x = 0, int y = 0) : x(x), y(y) {}
  bool operator==(const PointI &p) const { return x == p.x && y == p.y; }
  bool operator<(const PointI &p) const { return x != p.x ? x < p.x : y < p.y; }
};

int main() {
  vector<Point> v{{1, 3}, {1, 2}, {2, 1}, {0, 0}, {-1, 3}};
  vector<tuple<Point, Point, Point>> t{
      {Point{1, 3}, Point{1, 2}, Point{-1, 3}},
      {Point{1, 3}, Point{2, 1}, Point{1, 2}},
      {Point{1, 2}, Point{2, 1}, Point{0, 0}},
      {Point{1, 2}, Point{0, 0}, Point{-1, 3}}
  };
  assert(delaunay_triangulation(v) == t);

  // Integer-coordinate inputs are triangulated using exact integer arithmetic.
  vector<PointI> iv{{1, 3}, {1, 2}, {2, 1}, {0, 0}, {-1, 3}};
  vector<tuple<PointI, PointI, PointI>> ti{
      {PointI{1, 3}, PointI{1, 2}, PointI{-1, 3}},
      {PointI{1, 3}, PointI{2, 1}, PointI{1, 2}},
      {PointI{1, 2}, PointI{2, 1}, PointI{0, 0}},
      {PointI{1, 2}, PointI{0, 0}, PointI{-1, 3}}
  };
  assert(delaunay_triangulation(iv) == ti);
  return 0;
}
\end{lstlisting}
\subsection{Delaunay Triangulation (Guibas-Stolfi)}
\label{sec:7.4.6}
Given a set $P$ of distinct two-dimensional points, a Delaunay triangulation of $P$ is a triangulation of the convex hull of $P$ such that no point of $P$ lies strictly inside the circumcircle of any triangle in the triangulation. The Delaunay triangulation is not necessarily unique when four or more points are cocircular.

This implementation produces one valid triangulation using the Guibas-Stolfi divide-and-conquer algorithm with a quad-edge data structure. The point set is split in half, each half is triangulated recursively, and the halves are stitched together from the bottom up by a sequence of cross edges, with circumcircle tests deciding each connecting edge and deleting invalidated ones.

The point type must provide exact lexicographic \inlinecode{operator\textless{}}, since sorting and duplicate removal require a strict ordering rather than epsilon equality. All arithmetic uses the point's coordinate type, so integer inputs yield an exact triangulation.

\begin{compactitem}
  \item \inlinecode{delaunay\_triangulation(p)} returns the triangles of one Delaunay triangulation as a vector of counterclockwise-oriented \inlinecode{Point\textless{}T\textgreater{}} triples (\inlinecode{std::tuple}). Duplicate points are ignored. If fewer than three non-collinear unique points remain, the result is empty.
\end{compactitem}

Overflow warning: the in-circle test grows like the fourth power of the coordinate magnitude and dominates the overflow budget. With 64-bit integer coordinates, it overflows once coordinates exceed roughly $14000$. For floating-point coordinates, the predicates are subject to the usual rounding error.

\textbf{Time Complexity}: $O(n \log n)$ per call, where $n$ is the number of input points.

\Needspace*{3\baselineskip}
\textbf{Space Complexity}:
\begin{compactitem}
  \item $O(n \log n)$ for allocated quad-edges in the worst case because deleted edges remain in the pool.
  \item $O(n)$ for the returned triangulation.
\end{compactitem}

\begin{lstlisting}[style=cpp]
#include <algorithm>
#include <cmath>
#include <tuple>
#include <type_traits>
#include <utility>
#include <vector>

const double EPS = 1e-9;

template<typename T, typename U, typename C = std::common_type_t<T, U>>
bool EQ(T a, U b) {
  if constexpr (std::is_floating_point_v<C>) return C(a) == C(b) || std::fabs(C(a) - C(b)) <= EPS;
  return C(a) == C(b);
}

template<typename T, typename U, typename C = std::common_type_t<T, U>>
bool LT(T a, U b) {
  if constexpr (std::is_floating_point_v<C>) return C(a) < C(b) - EPS;
  return C(a) < C(b);
}

template<typename T>
struct Point {
  T x, y;
  Point(T x = 0, T y = 0) : x(x), y(y) {}
  bool operator==(const Point &p) const { return x == p.x && y == p.y; }
  friend bool EQ(const Point &a, const Point &b) { return EQ(a.x, b.x) && EQ(a.y, b.y); }
  bool operator<(const Point &p) const { return std::tie(x, y) < std::tie(p.x, p.y); }
};

template<typename T>
T cross(const Point<T> &a, const Point<T> &b, const Point<T> &o) {
  return (a.x - o.x) * (b.y - o.y) - (a.y - o.y) * (b.x - o.x);  // Overflow warning.
}

template<typename T>
struct Edge {
  Point<T> origin;
  Edge *rot, *onext;
  bool primal, deleted, used;

  Edge *rev() const { return rot->rot; }
  Edge *lnext() const { return rot->rev()->onext->rot; }
  Edge *oprev() const { return rot->onext->rot; }
  Point<T> dest() const { return rev()->origin; }
};

template<typename T>
struct QuadEdgePool {
  std::vector<Edge<T> *> edges;

  QuadEdgePool() = default;
  QuadEdgePool(const QuadEdgePool &) = delete;
  QuadEdgePool &operator=(const QuadEdgePool &) = delete;

  ~QuadEdgePool() {
    for (Edge<T> *e : edges) {
      delete e;
    }
  }

  Edge<T> *make_edge(const Point<T> &u, const Point<T> &v) {
    Edge<T> *e1 = new Edge<T>, *e2 = new Edge<T>, *e3 = new Edge<T>, *e4 = new Edge<T>;
    edges.push_back(e1);
    edges.push_back(e2);
    edges.push_back(e3);
    edges.push_back(e4);
    e1->origin = u;
    e2->origin = v;
    e1->rot = e3;
    e2->rot = e4;
    e3->rot = e2;
    e4->rot = e1;
    e1->onext = e1;
    e2->onext = e2;
    e3->onext = e4;
    e4->onext = e3;
    e1->primal = e2->primal = true;
    e3->primal = e4->primal = false;
    e1->deleted = e2->deleted = e3->deleted = e4->deleted = false;
    e1->used = e2->used = e3->used = e4->used = false;
    return e1;
  }

  void delete_edge(Edge<T> *e) {
    splice(e, e->oprev());
    splice(e->rev(), e->rev()->oprev());
    e->deleted = e->rot->deleted = e->rev()->deleted = e->rev()->rot->deleted = true;
  }

  Edge<T> *connect(Edge<T> *a, Edge<T> *b) {
    Edge<T> *e = make_edge(a->dest(), b->origin);
    splice(e, a->lnext());
    splice(e->rev(), b);
    return e;
  }

  static void splice(Edge<T> *a, Edge<T> *b) {
    std::swap(a->onext->rot->onext, b->onext->rot->onext);
    std::swap(a->onext, b->onext);
  }
};

template<typename T>
bool left_of(const Point<T> &p, Edge<T> *e) {
  return LT(0, cross(e->origin, e->dest(), p));
}

template<typename T>
bool right_of(const Point<T> &p, Edge<T> *e) {
  return LT(cross(e->origin, e->dest(), p), 0);
}

template<typename T>
bool in_circle(const Point<T> &a, const Point<T> &b, const Point<T> &c, const Point<T> &d) {
  T adx = a.x - d.x, ady = a.y - d.y;
  T bdx = b.x - d.x, bdy = b.y - d.y;
  T cdx = c.x - d.x, cdy = c.y - d.y;
  T abdet = adx * bdy - bdx * ady;  // Overflow warning.
  T bcdet = bdx * cdy - cdx * bdy;
  T cadet = cdx * ady - adx * cdy;
  T alift = adx * adx + ady * ady;
  T blift = bdx * bdx + bdy * bdy;
  T clift = cdx * cdx + cdy * cdy;
  return LT(0, alift * bcdet + blift * cadet + clift * abdet);
}

template<typename T>
std::pair<Edge<T> *, Edge<T> *> build_triangulation(
    const std::vector<Point<T>> &p, int lo, int hi, QuadEdgePool<T> &pool
) {
  if (hi - lo == 1) {
    Edge<T> *a = pool.make_edge(p[lo], p[hi]);
    return {a, a->rev()};
  }
  if (hi - lo == 2) {
    Edge<T> *a = pool.make_edge(p[lo], p[lo + 1]);
    Edge<T> *b = pool.make_edge(p[lo + 1], p[hi]);
    pool.splice(a->rev(), b);
    int side = LT(0, cross(p[lo], p[lo + 1], p[hi]))
                   ? 1
                   : (LT(cross(p[lo], p[lo + 1], p[hi]), 0) ? -1 : 0);
    if (side == 0) {
      return {a, b->rev()};
    }
    Edge<T> *c = pool.connect(b, a);
    if (side == 1) {
      return {a, b->rev()};
    }
    return {c->rev(), c};
  }
  int mid = lo + (hi - lo) / 2;
  auto [ldo, ldi] = build_triangulation(p, lo, mid, pool);
  auto [rdi, rdo] = build_triangulation(p, mid + 1, hi, pool);
  while (true) {
    if (left_of(rdi->origin, ldi)) {
      ldi = ldi->lnext();
    } else if (right_of(ldi->origin, rdi)) {
      rdi = rdi->rev()->onext;
    } else {
      break;
    }
  }
  Edge<T> *base = pool.connect(rdi->rev(), ldi);
  if (ldi->origin == ldo->origin) {
    ldo = base->rev();
  }
  if (rdi->origin == rdo->origin) {
    rdo = base;
  }
  while (true) {
    Edge<T> *lcand = base->rev()->onext;
    if (right_of(lcand->dest(), base)) {
      while (in_circle(base->dest(), base->origin, lcand->dest(), lcand->onext->dest())) {
        Edge<T> *next = lcand->onext;
        pool.delete_edge(lcand);
        lcand = next;
      }
    }
    Edge<T> *rcand = base->oprev();
    if (right_of(rcand->dest(), base)) {
      while (in_circle(base->dest(), base->origin, rcand->dest(), rcand->oprev()->dest())) {
        Edge<T> *prev = rcand->oprev();
        pool.delete_edge(rcand);
        rcand = prev;
      }
    }
    bool lvalid = right_of(lcand->dest(), base), rvalid = right_of(rcand->dest(), base);
    if (!lvalid && !rvalid) {
      break;
    }
    if (!lvalid ||
        (rvalid && in_circle(lcand->dest(), lcand->origin, rcand->origin, rcand->dest()))) {
      base = pool.connect(rcand, base->rev());
    } else {
      base = pool.connect(base->rev(), lcand->rev());
    }
  }
  return {ldo, rdo};
}

template<typename T>
auto delaunay_triangulation(std::vector<Point<T>> p) {
  using Pt = Point<T>;
  using Triangle = std::tuple<Pt, Pt, Pt>;
  // Rotates triangle vertices so that the lexicographically smallest comes first, preserving the
  // cyclic (counterclockwise) order, to give each triangle a canonical, comparable representation.
  auto canonical = [](const Pt &a, const Pt &b, const Pt &c) -> Triangle {
    if (b < a && b < c) {
      return {b, c, a};
    }
    if (c < a && c < b) {
      return {c, a, b};
    }
    return {a, b, c};
  };
  std::sort(p.begin(), p.end());
  auto same_point = [](const Pt &a, const Pt &b) { return !(a < b) && !(b < a); };
  p.erase(std::unique(p.begin(), p.end(), same_point), p.end());
  std::vector<Triangle> res;
  if (p.size() < 3) {
    return res;
  }
  QuadEdgePool<T> pool;
  build_triangulation(p, 0, static_cast<int>(p.size()) - 1, pool);
  for (Edge<T> *start : pool.edges) {
    if (!start->primal || start->deleted || start->used) {
      continue;
    }
    std::vector<Pt> face;
    Edge<T> *curr = start;
    do {
      curr->used = true;
      face.push_back(curr->origin);
      curr = curr->lnext();
    } while (curr != start);
    if (face.size() == 3 && LT(0, cross(face[0], face[1], face[2]))) {
      res.push_back(canonical(face[0], face[1], face[2]));
    }
  }
  std::sort(res.begin(), res.end());
  return res;
}
\end{lstlisting}
\Needspace*{4\baselineskip}
\textbf{Example Usage}\par
\begin{lstlisting}[style=cpp]
#include <cassert>
#include <cstdint>
using namespace std;

using PointL = Point<int64_t>;
using PointD = Point<double>;

int main() {
  vector<PointD> v{{1, 3}, {1, 2}, {2, 1}, {0, 0}, {-1, 3}};
  vector<tuple<PointD, PointD, PointD>> t{
      {PointD{-1, 3}, PointD{0, 0}, PointD{1, 2}},
      {PointD{-1, 3}, PointD{1, 2}, PointD{1, 3}},
      {PointD{0, 0}, PointD{2, 1}, PointD{1, 2}},
      {PointD{1, 2}, PointD{2, 1}, PointD{1, 3}}
  };
  assert(delaunay_triangulation(v) == t);

  // Integer-coordinate inputs are triangulated using exact integer arithmetic.
  vector<PointL> iv{{1, 3}, {1, 2}, {2, 1}, {0, 0}, {-1, 3}};
  vector<tuple<PointL, PointL, PointL>> ti{
      {PointL{-1, 3}, PointL{0, 0}, PointL{1, 2}},
      {PointL{-1, 3}, PointL{1, 2}, PointL{1, 3}},
      {PointL{0, 0}, PointL{2, 1}, PointL{1, 2}},
      {PointL{1, 2}, PointL{2, 1}, PointL{1, 3}}
  };
  assert(delaunay_triangulation(iv) == ti);
  return 0;
}
\end{lstlisting}

\section{\texorpdfstring{3D Geometry}{7.5 3D Geometry}}
\setcounter{section}{5}
\setcounter{subsection}{0}
\subsection{3D Point}
\label{sec:7.5.1}
A 3D point class templated on the coordinate type \inlinecode{T}. Algebraic operations preserve \inlinecode{T}, while metric operations use \inlinecode{fp\_t}, which is \inlinecode{T} for floating-point coordinates and \inlinecode{double} otherwise.

\begin{compactitem}
  \item \inlinecode{TPoint3\textless{}T\textgreater{}(x = 0, y = 0, z = 0)} constructs point (\inlinecode{x}, \inlinecode{y}, \inlinecode{z}).
  \item Operators \inlinecode{+}, \inlinecode{-}, \inlinecode{*}, and \inlinecode{/} preserve the coordinate type except that division promotes to \inlinecode{fp\_t}. Comparisons are exact and lexicographic, while \inlinecode{EQ(p, q)} uses \inlinecode{EPS} for floating-point coordinates and remains exact for other types.
  \item \inlinecode{sqnorm()}, \inlinecode{dot(p)}, and \inlinecode{cross(p)} return exact algebraic results in coordinate type \inlinecode{T}.
  \item \inlinecode{norm()}, \inlinecode{phi()}, and \inlinecode{theta()} return the length, azimuth, and polar angle in \inlinecode{fp\_t}.
  \item \inlinecode{unit()} returns the unit vector and requires a nonzero point; \inlinecode{normal(p)} returns a unit normal and requires nonparallel vectors.
  \item \inlinecode{rotate(angle, axis)} rotates the point about a nonzero \inlinecode{axis} using Rodrigues' formula.
  \item \inlinecode{Point3I}, \inlinecode{Point3L}, \inlinecode{Point3D}, and \inlinecode{Point3LD} use \inlinecode{int}, \inlinecode{int64\_t}, \inlinecode{double}, and \inlinecode{long double} coordinates, respectively.
\end{compactitem}

Overflow warning: the exact products \inlinecode{dot()}, \inlinecode{cross()}, and \inlinecode{sqnorm()} grow like the squared coordinate magnitude. With \inlinecode{TPoint3\textless{}int\textgreater{}} these overflow a 32-bit \inlinecode{int} once coordinates exceed a few tens of thousands, so use \inlinecode{Point3L} (\inlinecode{TPoint3\textless{}int64\_t\textgreater{}}) for larger integer coordinates.

\textbf{Time Complexity}: $O(1)$ per operation.

\textbf{Space Complexity}: $O(1)$ for storage and auxiliary.

\begin{lstlisting}[style=cpp]
#include <cmath>
#include <cstdint>
#include <tuple>
#include <type_traits>

const double EPS = 1e-9;

template<typename T, typename U, typename C = std::common_type_t<T, U>>
bool EQ(T a, U b) {
  if constexpr (std::is_floating_point_v<C>) return C(a) == C(b) || std::fabs(C(a) - C(b)) <= EPS;
  return C(a) == C(b);
}

template<typename T>
struct TPoint3 {
  using fp_t = std::conditional_t<std::is_floating_point<T>::value, T, double>;

  T x, y, z;

  TPoint3(T x = 0, T y = 0, T z = 0) : x(x), y(y), z(z) {}
  bool operator==(const TPoint3 &p) const { return std::tie(x, y, z) == std::tie(p.x, p.y, p.z); }

  friend bool EQ(const TPoint3 &a, const TPoint3 &b) {
    return EQ(a.x, b.x) && EQ(a.y, b.y) && EQ(a.z, b.z);
  }

  bool operator!=(const TPoint3 &p) const { return !(*this == p); }
  bool operator<(const TPoint3 &p) const { return std::tie(x, y, z) < std::tie(p.x, p.y, p.z); }
  TPoint3 operator+(const TPoint3 &p) const { return {x + p.x, y + p.y, z + p.z}; }
  TPoint3 operator-(const TPoint3 &p) const { return {x - p.x, y - p.y, z - p.z}; }
  TPoint3 operator*(T k) const { return {x * k, y * k, z * k}; }
  TPoint3<fp_t> operator/(fp_t k) const { return {fp_t(x) / k, fp_t(y) / k, fp_t(z) / k}; }
  T dot(const TPoint3 &p) const { return x * p.x + y * p.y + z * p.z; }  // Overflow warning.
  T sqnorm() const { return x * x + y * y + z * z; }                     // Overflow warning.
  fp_t norm() const { return std::hypot(std::hypot(fp_t(x), fp_t(y)), fp_t(z)); }
  fp_t phi() const { return std::atan2(fp_t(y), fp_t(x)); }
  fp_t theta() const { return std::atan2(std::hypot(fp_t(x), fp_t(y)), fp_t(z)); }
  TPoint3<fp_t> unit() const { return *this / norm(); }

  TPoint3 cross(const TPoint3 &p) const {
    return {y * p.z - z * p.y, z * p.x - x * p.z, x * p.y - y * p.x};  // Overflow warning.
  }

  TPoint3<fp_t> normal(const TPoint3 &p) const { return cross(p).unit(); }

  TPoint3<fp_t> rotate(fp_t angle, const TPoint3 &axis) const {
    fp_t s = std::sin(angle), c = std::cos(angle);
    TPoint3<fp_t> u = axis.unit(), v{fp_t(x), fp_t(y), fp_t(z)};
    return u * v.dot(u) * (1 - c) + v * c - v.cross(u) * s;
  }
};

using Point3I = TPoint3<int>;
using Point3L = TPoint3<int64_t>;
using Point3D = TPoint3<double>;
using Point3LD = TPoint3<long double>;
\end{lstlisting}
\Needspace*{4\baselineskip}
\textbf{Example Usage}\par
\begin{lstlisting}[style=cpp]
#include <cassert>
using namespace std;

int main() {
  Point3I a(1, 0, 0), b(0, 1, 0);
  assert(a.dot(b) == 0);
  assert(a.cross(b) == Point3I(0, 0, 1));
  assert(a.sqnorm() == 1);

  Point3D p(1, 0, 0);
  assert(EQ(p.rotate(acos(-1.0) / 2, Point3D(0, 0, 1)), Point3D(0, 1, 0)));
  return 0;
}
\end{lstlisting}
\subsection{Polyhedron Volume}
\label{sec:7.5.2}
Given an oriented triangular mesh of a closed polyhedron, compute its signed volume. Each triangle contributes the signed volume of the tetrahedron formed by the triangle $abc$ and the origin $((a \times b) \cdot c) / 6$. If faces are oriented outward, the result is positive for the orientation used by \inlinecode{convex\_hull\_3d()} in \secref{7.5.3}; take \inlinecode{abs()} if only the magnitude is needed.

\begin{compactitem}
  \item \inlinecode{signed\_polyhedron\_volume(p, faces)} returns the signed volume. Each face triplet (\inlinecode{a}, \inlinecode{b}, \inlinecode{c}) represents the triangle with vertices \inlinecode{p[a]}, \inlinecode{p[b]}, and \inlinecode{p[c]}, and its indices must be valid.
\end{compactitem}

Overflow warning: for integral point types, each scalar triple product must fit in the point's coordinate arithmetic type before it is converted to \inlinecode{double}.

\textbf{Time Complexity}: $O(f)$ per call, where $f$ is the number of triangular faces.

\textbf{Space Complexity}: $O(1)$ auxiliary.

\begin{lstlisting}[style=cpp]
#include <algorithm>
#include <cassert>
#include <vector>

struct Point3D {
  double x, y, z;

  Point3D(double x = 0, double y = 0, double z = 0) : x(x), y(y), z(z) {}
  double dot(const Point3D &p) const { return x * p.x + y * p.y + z * p.z; }

  Point3D cross(const Point3D &p) const {
    return {y * p.z - z * p.y, z * p.x - x * p.z, x * p.y - y * p.x};
  }
};

template<typename Pt, typename F>
double signed_polyhedron_volume(const std::vector<Pt> &p, const std::vector<F> &faces) {
  assert(std::all_of(faces.begin(), faces.end(), [&](const F &face) {
    auto [a, b, c] = face;
    return 0 <= a && a < static_cast<int>(p.size()) && 0 <= b && b < static_cast<int>(p.size()) &&
           0 <= c && c < static_cast<int>(p.size());
  }));
  double volume = 0;
  for (const auto &[a, b, c] : faces) {
    volume += p[a].cross(p[b]).dot(p[c]);  // Overflow warning.
  }
  return volume / 6.0;
}
\end{lstlisting}
\Needspace*{4\baselineskip}
\textbf{Example Usage}\par
\begin{lstlisting}[style=cpp]
#include <cassert>
#include <cmath>
#include <tuple>
using namespace std;

bool EQ(double a, double b) {
  return fabs(a - b) < 1e-9;
}

int main() {
  using Face = tuple<int, int, int>;  // (a, b, c)
  vector<Point3D> p{{0, 0, 0}, {1, 0, 0}, {0, 1, 0}, {0, 0, 1}};
  vector<Face> faces{
      Face(0, 2, 1),
      Face(0, 1, 3),
      Face(0, 3, 2),
      Face(1, 2, 3),
  };
  assert(EQ(signed_polyhedron_volume(p, faces), 1.0 / 6.0));
  return 0;
}
\end{lstlisting}
\subsection{3D Convex Hull}
\label{sec:7.5.3}
The convex hull of points in 3D is the smallest convex polyhedron containing them. This incremental algorithm computes its triangular faces: start with a tetrahedron, delete every visible face for each new point, and stitch the horizon edges to that point. Returned faces are oriented outward.

\begin{compactitem}
  \item \inlinecode{convex\_hull\_3d(p)} returns triangular faces as \inlinecode{Face\{a, b, c\}} triples, where each face has vertices \inlinecode{p[a]}, \inlinecode{p[b]}, and \inlinecode{p[c]}. The input must contain at least four points, with no four coplanar. For coplanar or highly degenerate inputs, perturb the points slightly or use a more robust library.
\end{compactitem}

\textbf{Time Complexity}: $O(n^{2})$ per call, where $n$ is the number of points.

\textbf{Space Complexity}: $O(n^{2})$ auxiliary.

\begin{lstlisting}[style=cpp]
#include <array>
#include <cassert>
#include <utility>
#include <vector>

const double EPS = 1e-9;

struct Point3D {
  double x, y, z;

  Point3D(double x = 0, double y = 0, double z = 0) : x(x), y(y), z(z) {}
  Point3D operator-(const Point3D &p) const { return {x - p.x, y - p.y, z - p.z}; }
  Point3D operator*(double k) const { return {x * k, y * k, z * k}; }
  double dot(const Point3D &p) const { return x * p.x + y * p.y + z * p.z; }

  Point3D cross(const Point3D &p) const {
    return {y * p.z - z * p.y, z * p.x - x * p.z, x * p.y - y * p.x};
  }
};

struct Face {
  int a, b, c;
  Face(int a = 0, int b = 0, int c = 0) : a(a), b(b), c(c) {}
};

std::vector<Face> convex_hull_3d(const std::vector<Point3D> &p) {
  int n = static_cast<int>(p.size());
  assert(n >= 4);
  std::vector<std::vector<std::array<int, 2>>> edge(
      n, std::vector<std::array<int, 2>>(n, {-1, -1})
  );
  std::vector<Face> faces;
  auto normal = [&](const Face &f) { return (p[f.b] - p[f.a]).cross(p[f.c] - p[f.a]); };
  auto visible = [&](const Face &f, int i) { return normal(f).dot(p[i] - p[f.a]) > EPS; };
  auto add_edge = [&](int a, int b, int c) {
    if (a > b) {
      std::swap(a, b);
    }
    auto &e = edge[a][b];
    (e[0] == -1 ? e[0] : e[1]) = c;
  };
  auto remove_edge = [&](int a, int b, int c) {
    if (a > b) {
      std::swap(a, b);
    }
    auto &e = edge[a][b];
    (e[0] == c ? e[0] : e[1]) = -1;
  };
  auto edge_count = [&](int a, int b) {
    if (a > b) {
      std::swap(a, b);
    }
    return (edge[a][b][0] != -1) + (edge[a][b][1] != -1);
  };
  auto make_face = [&](int a, int b, int c, int inside) {
    Face f(a, b, c);
    if (normal(f).dot(p[inside] - p[a]) > 0) {
      std::swap(f.b, f.c);
    }
    add_edge(f.a, f.b, f.c);
    add_edge(f.b, f.c, f.a);
    add_edge(f.c, f.a, f.b);
    faces.push_back(f);
  };
  for (int i = 0; i < 4; i++) {
    for (int j = i + 1; j < 4; j++) {
      for (int k = j + 1; k < 4; k++) {
        make_face(i, j, k, 6 - i - j - k);
      }
    }
  }
  for (int i = 4; i < n; i++) {
    for (int j = 0; j < static_cast<int>(faces.size()); j++) {
      Face f = faces[j];
      if (visible(f, i)) {
        remove_edge(f.a, f.b, f.c);
        remove_edge(f.b, f.c, f.a);
        remove_edge(f.c, f.a, f.b);
        std::swap(faces[j--], faces.back());
        faces.pop_back();
      }
    }
    int old_faces = static_cast<int>(faces.size());
    for (int j = 0; j < old_faces; j++) {
      Face f = faces[j];
      if (edge_count(f.a, f.b) != 2) {
        make_face(f.a, f.b, i, f.c);
      }
      if (edge_count(f.b, f.c) != 2) {
        make_face(f.b, f.c, i, f.a);
      }
      if (edge_count(f.c, f.a) != 2) {
        make_face(f.c, f.a, i, f.b);
      }
    }
  }
  return faces;
}
\end{lstlisting}
\Needspace*{4\baselineskip}
\textbf{Example Usage}\par
\begin{lstlisting}[style=cpp]
using namespace std;

void assert_outward(const vector<Point3D> &p, const vector<Face> &faces) {
  for (const Face &f : faces) {
    assert(0 <= f.a && f.a < static_cast<int>(p.size()));
    assert(0 <= f.b && f.b < static_cast<int>(p.size()));
    assert(0 <= f.c && f.c < static_cast<int>(p.size()));
    assert(f.a != f.b && f.b != f.c && f.c != f.a);
    Point3D normal = (p[f.b] - p[f.a]).cross(p[f.c] - p[f.a]);
    for (const Point3D &q : p) {
      assert(normal.dot(q - p[f.a]) <= EPS);
    }
  }
}

int main() {
  vector<Point3D> tetra{{0, 0, 0}, {1, 0, 0}, {0, 1, 0}, {0, 0, 1}};
  auto faces = convex_hull_3d(tetra);
  // A tetrahedron has four triangular faces.
  assert(faces.size() == 4);
  assert_outward(tetra, faces);

  tetra.emplace_back(1, 1, 1);
  faces = convex_hull_3d(tetra);
  // Adding the opposite point makes a triangular bipyramid with six triangular faces.
  assert(faces.size() == 6);
  assert_outward(tetra, faces);
  return 0;
}
\end{lstlisting}
